% ===========================================================
% Diseñados para amar — Versión simplificada para pandoc/EPUB
% ===========================================================
\documentclass[11pt]{book}

% --- Idioma ---
\usepackage[spanish]{babel}

% --- Tipografía ---
\usepackage{fontspec}

% --- Imágenes ---
\usepackage{graphicx}
\graphicspath{{../img/}}

% --- Epígrafes ---
\usepackage{epigraph}

% --- Citas ---
\usepackage{csquotes}

% --- Listas ---
\usepackage{enumitem}

% --- Comando simplificado: Para caminar juntos ---
\newcommand{\paracaminar}{%
  \begin{center}\rule{0.3\textwidth}{0.4pt}\end{center}
  \subsection*{Para caminar juntos}
}

% --- Comando simplificado: viñeta de capítulo ---
\newcommand{\vinetacapitulo}[1]{%
  \begin{center}
    \includegraphics[width=0.5\textwidth]{#1.png}
  \end{center}
}

\begin{document}

% --- Prólogo ---
\chapter*{Prólogo}
\addcontentsline{toc}{chapter}{Prólogo}

\markboth{Prólogo}{Prólogo}

Estas páginas nacieron sin pretenderlo.

Durante algunos años he tenido la suerte de participar en cursos de preparación al matrimonio y de compartir estas ideas en tertulias con amigos, con parejas a punto de casarse y con matrimonios que llevan ya un buen trecho recorrido. Lo que aquí se recoge son, en esencia, esas conversaciones: reflexiones en voz alta sobre por qué estamos hechos para amar, qué tiene que ver nuestro cuerpo con el amor, y por qué el matrimonio cristiano, lejos de ser un corsé, es probablemente el regalo más grande que dos personas pueden hacerse.

No encontrarás aquí un tratado de teología ni un manual de autoayuda matrimonial. Son apuntes personales, con las limitaciones que eso implica. La curiosidad por este tema me ha llevado a explorar terrenos que no son los míos ---algo de biología, algo de neurociencia, algo de teología--- y lo que he encontrado por el camino es lo que intento compartir en estas páginas. Si en algún punto he simplificado demasiado o he sido impreciso, la responsabilidad es enteramente mía.

Lo que sí puedo asegurar es que todo lo que aquí se dice ha pasado por la experiencia. No son ideas encontradas en un manual y copiadas sin más. Son cosas que he visto funcionar en mi propio matrimonio, en el de amigos cercanos, y en los rostros de parejas que, al escuchar algunas de estas reflexiones, han dicho: ``Nunca lo había visto así.'' Si estos apuntes consiguen provocar ese mismo gesto en alguien, habrán cumplido su propósito.

Este libro cuenta una historia: la visión cristiana del matrimonio, apoyada en lo que la biología y la antropología nos van descubriendo. Hay otras historias, otras formas de amar, otros caminos. No las niego ni las descalifico. Simplemente cuento esta, con toda la honestidad de la que soy capaz.

La ciencia describe lo que describe: que estamos biológicamente configurados para el vínculo, que la vulnerabilidad reclama cuidado, que nuestra química nos empuja hacia el apego. La ciencia no demuestra que haya un Diseñador. Pero tampoco demuestra que no lo haya. Ante esos mismos datos, yo elijo creer que este diseño extraordinario apunta a Alguien que es Amor. No es una conclusión que pueda imponerse; es una apuesta personal, un acto de confianza. Acepto que otros, con igual honestidad, miren los mismos hechos y no den ese paso. El salto de la biología a la fe no es una demostración lógica: es una invitación. El lector es libre de aceptarla o no. Los datos biológicos siguen siendo verdad en cualquier caso. Pero si estás dispuesto a considerar que el \emph{cómo} puede revelar un \emph{para qué}, estas páginas te ofrecen un camino que merece la pena recorrer.

Estas páginas están vivas y quieren mejorar. Si tienes cualquier comentario, corrección o sugerencia, escríbeme a \texttt{jtorrecilla@gmail.com}.

\vspace{2em}

\begin{flushright}
\textit{Jesús Torrecilla Pinero}\\
\textit{Marzo de 2026}
\end{flushright}

% ===========================
% PARTE I: HECHOS PARA AMAR
% ===========================
\part{Hechos para amar}

\chapter{Nacer humano}

\epigraph{Todos los mamíferos nacen. Pero solo el ser humano necesita que alguien le sostenga la cabeza.}{}

Hay algo que casi nadie se para a pensar, y que sin embargo dice mucho de lo que somos: el parto humano es, con diferencia, el más difícil del reino animal.

No hablo solo del dolor, aunque también. Hablo de algo más profundo. De un conflicto que lleva inscrito en nuestro cuerpo millones de años de historia. Un conflicto entre dos cosas que nos hacen humanos ---caminar erguidos y pensar--- y que, en el momento de nacer, chocan de frente. Para entender qué tiene esto que ver con el amor, con el matrimonio y con la familia, hay que empezar por el principio. Y el principio es la sabana africana, hace unos cuatro millones de años.

\section{Ponerse de pie}

Imaginemos a una hembra de \textit{Australopithecus afarensis} ---quizá deberíamos llamarla por su nombre más famoso: Lucy.\footnote{El esqueleto de Lucy fue descubierto en Etiopía en 1974 por Donald Johanson. Tiene unos 3,2 millones de años de antigüedad y fue clave para entender el origen del bipedalismo.} Lucy medía poco más de un metro, pesaba unos treinta kilos y hacía algo que ninguno de sus antepasados había hecho de forma habitual: caminaba sobre dos piernas.

No fue un capricho. El bipedalismo traía ventajas enormes. En una sabana cada vez más abierta, donde los bosques retrocedían, caminar erguido permitía ver por encima de la hierba alta ---detectar al depredador antes de que fuera demasiado tarde---. Liberaba las manos para cargar alimentos, para transportar herramientas rudimentarias, para sostener a una cría. Y era más eficiente: caminar sobre dos piernas gasta menos energía que hacerlo sobre cuatro.\footnote{Estudios comparativos muestran que el bipedalismo humano consume aproximadamente un 75\% menos de energía que la locomoción cuadrúpeda de un chimpancé de tamaño equivalente. Cf. Sockol, Raichlen y Pontzer, ``Chimpanzee locomotor energetics and the origin of human bipedalism'', \textit{PNAS}, 2007.}

Pero todo tiene un precio. Para caminar erguidos de manera eficiente, la pelvis tuvo que cambiar de forma. Se acortó, se ensanchó por arriba y se estrechó por abajo, formando una especie de cuenco compacto que sostenía las vísceras y transfería el peso a las piernas. La pelvis de Lucy ya no se parecía a la de un chimpancé. Se parecía cada vez más a la nuestra.

Y eso iba a ser un problema.

\section{El cerebro que no para de crecer}

Mientras la pelvis se estrechaba, algo más ocurría. Lentamente, generación tras generación, el cerebro de nuestros antepasados comenzó a crecer.

No fue solo un aumento de tamaño. Fue un cambio cualitativo. La corteza cerebral se plegaba en circunvoluciones cada vez más complejas. Un cambio de dieta tuvo mucho que ver: cuando nuestros ancestros empezaron a consumir más proteína animal ---primero carroña, luego caza--- y a cocinar los alimentos con fuego,\footnote{La ``hipótesis del cocinado'' de Richard Wrangham propone que la cocción de alimentos fue clave para el aumento del cerebro humano, al hacer los nutrientes más accesibles y reducir la energía necesaria para la digestión. Cf. Wrangham, \textit{Catching Fire: How Cooking Made Us Human}, Basic Books, 2009.} el intestino pudo acortarse y la energía sobrante se redirigió al órgano más caro del cuerpo. El cerebro humano supone apenas el 2\% del peso corporal, pero consume el 20\% de la energía.\footnote{En comparación, el cerebro de un primate típico consume alrededor del 8-10\% de la energía corporal.}

Más neuronas, más conexiones, más capacidad. Aquello nos permitió fabricar herramientas cada vez mejores, planificar, comunicarnos con un lenguaje articulado, cooperar en grupo. Nos estábamos convirtiendo en humanos.

Pero un cerebro más grande necesita una cabeza más grande para alojarlo. El cráneo de un \textit{Homo sapiens} recién nacido tiene un volumen medio de unos 350 centímetros cúbicos ---y seguirá creciendo después del parto hasta los 1.400 aproximados del adulto---. Compare eso con los 150-170 de un chimpancé neonato, que además nace con un grado de desarrollo mucho mayor.

Tenemos, entonces, dos fuerzas evolutivas que nos hacen humanos y que tiran en direcciones opuestas: una pelvis que se estrecha para caminar, y una cabeza que crece para pensar. El punto de encuentro de esas dos fuerzas es el canal del parto. Y ahí es donde estalla lo que los antropólogos llaman el \textbf{dilema obstétrico}.\footnote{El término ``dilema obstétrico'' fue acuñado por el antropólogo Sherwood Washburn en 1960. Aunque investigaciones recientes matizan algunos aspectos de la hipótesis original, el conflicto básico entre bipedalismo y encefalización sigue siendo ampliamente aceptado. Cf. Washburn, ``Tools and Human Evolution'', \textit{Scientific American}, 1960.}

\section{El dilema obstétrico}

El conflicto es real y medible. En la mayor parte de los primates, la cabeza del bebé pasa por el canal del parto con cierta holgura. En los humanos, no. La cabeza del bebé se ajusta al canal pélvico como una llave en su cerradura ---con apenas milímetros de margen---. Por eso el parto humano es más largo, más doloroso y más peligroso que el de cualquier otro primate. Por eso la cabeza del bebé tiene que rotar durante el descenso, ejecutando una especie de giro helicoidal para adaptarse a la forma cambiante del canal.\footnote{Este mecanismo de rotación es exclusivo del parto humano. En otros primates, el bebé desciende sin necesidad de rotar. Es también la razón por la que el parto humano se beneficia de la asistencia: la cría nace mirando hacia atrás, lo que dificulta que la propia madre la reciba. Somos una de las pocas especies en las que el parto es normalmente asistido.} Por eso los huesos del cráneo del recién nacido no están soldados: las fontanelas permiten que la cabeza se deforme ligeramente para pasar.

La evolución tenía que resolver este conflicto, y lo hizo con una solución ingeniosa pero de consecuencias enormes: adelantar el parto.

Conviene ser honesto: no todos los investigadores explican este adelanto de la misma manera. Hay otra hipótesis, propuesta por Holly Dunsworth y sus colegas, que apunta no tanto a la pelvis como al metabolismo.\footnote{Dunsworth, H. et al., ``Metabolic hypothesis for human altriciality'', \textit{PNAS}, 2012. Dunsworth argumenta que el parto se desencadena cuando la madre alcanza el límite de su capacidad metabólica ---unas 2 a 2,5 veces su tasa basal---, independientemente de las dimensiones pélvicas.} Según esta visión, lo que limita la duración del embarazo no es que el bebé no quepa, sino que la madre no puede sostener energéticamente un feto cuyo cerebro crece a una velocidad sin precedentes en el reino animal. El embarazo se interrumpe cuando el coste metabólico se vuelve insostenible.

¿Quién tiene razón? Probablemente ambas hipótesis capturan parte de la verdad. Pero lo que importa para nuestro argumento es lo que las dos comparten: en ambos casos, la causa última es el cerebro humano ---demasiado grande para la pelvis o demasiado costoso para el metabolismo---, y en ambos casos el resultado es el mismo: nacemos antes de estar listos.

\section{El feto externo}

Si un bebé humano naciera con el mismo grado de madurez neurológica que un chimpancé recién nacido, la gestación debería durar entre 18 y 21 meses.\footnote{Cf. Portmann, A., ``Die Ontogenese des Menschen als Problem der Evolutionsforschung'', \textit{Verhandlungen der Schweizerischen Naturforschenden Gesellschaft}, 1941. Portmann fue el primero en calcular esta gestación teórica y en proponer el concepto de altricialidad secundaria.} Pero con 18 meses de desarrollo ---ya sea por las dimensiones de la cabeza o por el coste energético---, nuestro cuerpo no puede más. La solución: nacer a los nueve meses, a medio camino.

El resultado es que los bebés humanos llegan al mundo radicalmente inmaduros. Los biólogos lo llaman \textit{altricialidad secundaria}: una especie que, por su linaje, debería producir crías relativamente maduras al nacer (como los demás primates), pero que en la práctica produce crías tan indefensas como las de un ratón o un gorrión.

Pensemos en lo que esto significa. Un potro se pone de pie y camina a las pocas horas de nacer. Un cervatillo puede seguir a su madre por el bosque en su primer día. Incluso un chimpancé recién nacido tiene la fuerza suficiente para agarrarse al pelaje de su madre.

Un bebé humano no puede hacer nada de eso. No puede sostener la cabeza. No puede agarrarse. No puede desplazarse. No regula bien su temperatura. Su visión es borrosa. Su sistema inmunitario es inmaduro. Durante meses, es absolutamente incapaz de sobrevivir sin que alguien le cuide de forma constante.

El antropólogo Ashley Montagu lo expresó con una imagen que se ha hecho célebre: el primer año de vida del ser humano es una \textit{exterogestación}, una gestación fuera del útero.\footnote{Montagu, A., \textit{Touching: The Human Significance of the Skin}, Harper \& Row, 1971.} El bebé ha nacido, sí, pero su desarrollo está lejos de completarse. Su cerebro triplicará su tamaño en los primeros dos años, formando las conexiones que lo harán pensar, hablar, sentir.

Nacemos, en cierto modo, a medio hacer. Y esa es la clave de todo lo que viene después.

\section{Una especie que necesita amor para sobrevivir}

Pensemos ahora en lo que significaba esto en la sabana africana hace doscientos mil años. Una madre acaba de dar a luz. Tiene en brazos a una criatura que depende absolutamente de ella. Que necesita ser alimentada cada pocas horas. Que no puede quedarse sola ni un momento. Que hay que cargar a todas partes, porque no va a caminar en un año.

Esa madre necesita comer, necesita beber, necesita dormir. Necesita estar alerta a los depredadores. ¿Cómo puede hacer todo eso con un bebé en brazos?

La respuesta es que no puede. No sola.

Y aquí es donde la historia da un giro que lo cambia todo. Porque la evolución no solo adelantó el parto. También preparó algo para lo que venía después: mecanismos biológicos que empujan a cuidar.

Durante el parto y la lactancia, el cuerpo de la madre se inunda de oxitocina.\footnote{La oxitocina es un neuropéptido producido en el hipotálamo. Además de su papel en las contracciones uterinas y la eyección de leche, actúa en el cerebro generando sensaciones de confianza, calma y vinculación. Se la ha llamado ``hormona del vínculo'', aunque su función es más compleja: en determinados contextos también puede intensificar la desconfianza hacia los extraños, reforzando la protección del grupo cercano.} No es una hormona cualquiera. Es la hormona que genera vínculo, que empuja a cuidar, que hace que una madre sienta que daría la vida por esa criatura que acaba de conocer. El contacto piel con piel dispara más oxitocina. La lactancia la mantiene alta. El llanto del bebé activa en el cerebro de la madre circuitos de alarma y respuesta que son anteriores a cualquier decisión consciente.\footnote{Estudios de neuroimagen muestran que el llanto de un bebé activa la amígdala y el córtex cingulado anterior de la madre en milisegundos, antes de que haya procesamiento consciente. Cf. Swain et al., ``Brain basis of early parent-infant interactions'', \textit{Journal of Child Psychology and Psychiatry}, 2007.}

No es que la madre ``decida'' amar a su hijo. Es que su cuerpo la empuja a ello con una fuerza que precede a la voluntad. La biología no sustituye al amor libre ---una madre siempre puede rechazar a su cría, y el amor verdadero será siempre un acto de libertad---, pero le pone el viento a favor. Le da un empujón poderoso en la dirección correcta.

Y no solo a la madre. También al padre.

\section{El padre que se queda}

En la mayoría de los mamíferos, el macho se desentiende de la cría tras la cópula. Cumple su función reproductiva y desaparece. Pero en los humanos, algo distinto ocurrió.

Los estudios muestran que cuando un hombre se convierte en padre, su biología cambia.\footnote{Cf. Gettler et al., ``Longitudinal evidence that fatherhood decreases testosterone in human males'', \textit{PNAS}, 2011. Este estudio, realizado con más de 600 hombres filipinos, mostró que los niveles de testosterona disminuían significativamente al convertirse en padres, y más aún en quienes se implicaban activamente en el cuidado.} Los niveles de testosterona bajan ---la hormona que impulsa la competición y la búsqueda de nuevas parejas--- y suben los de oxitocina y vasopresina ---hormonas vinculadas al cuidado y la protección---. El cerebro del padre también se reorganiza: se activan circuitos de empatía y respuesta al llanto del bebé que antes estaban más silenciosos.

¿Por qué la evolución favoreció esto? Por una razón implacable: en un entorno donde los bebés son tan vulnerables y la madre está tan exigida, aquellos grupos donde el padre permanecía, donde ayudaba a alimentar y proteger, tenían más probabilidades de que sus crías sobrevivieran. Y esas crías crecían, se reproducían, y transmitían esa misma inclinación. No la inclinación a ser ``buenos'', en abstracto, sino la inclinación biológica a vincularse, a quedarse, a cuidar.

El amor de pareja, el deseo de estar juntos, la satisfacción de construir algo común, tienen una raíz que va más allá del romanticismo: son mecanismos que aseguran que una cría extremadamente vulnerable llegue a la edad adulta. El amor no es un adorno de la vida. Es lo que hizo posible que la vida continuara.

\section{La vulnerabilidad que nos abre}

Hay algo asombroso en todo esto, si nos paramos a pensarlo.

Nuestra grandeza como especie ---ese cerebro extraordinario que nos permite componer sinfonías, curar enfermedades, preguntarnos por el sentido de la vida--- tiene un precio: nacemos incompletos. Nacemos necesitados. Nacemos absolutamente dependientes del amor de otros.

Y esa dependencia no es un fallo de diseño. Es el diseño.

Porque es precisamente nuestra vulnerabilidad la que nos abre al otro. Si fuéramos autosuficientes al nacer, como una tortuga que rompe el cascarón y se arrastra sola hacia el mar, no necesitaríamos vínculos. No necesitaríamos familias. No necesitaríamos el abrazo de una madre ni la presencia de un padre. No necesitaríamos amor.

Pero no somos tortugas. Somos la especie que nace con la cabeza sin cerrar, que llora para llamar a alguien, que necesita años de cuidado antes de poder valerse por sí misma. Somos la especie que, para sobrevivir, tuvo que aprender a amar.

Y quizá eso diga algo profundo sobre lo que somos. Quizá nuestro cuerpo, con su pelvis estrecha y su cerebro enorme, con sus hormonas que empujan al vínculo y su llanto que activa el cuidado, nos esté diciendo algo que intuimos pero que a veces olvidamos: que no estamos hechos para vivir solos. Que el amor no es un añadido opcional a la existencia humana, sino su condición de posibilidad.

Esa es la primera lección que nos da el cuerpo. Y es solo el comienzo.

\paracaminar

\begin{enumerate}
  \item ¿Alguna vez os habéis planteado que el cuerpo humano está ``diseñado'' para la relación? ¿Qué os sugiere esa idea?

  \item Si habéis tenido la experiencia de tener un bebé en brazos ---propio o ajeno---, ¿qué sentisteis? ¿Reconocéis esa ``fuerza'' biológica que empuja a cuidar?

  \item Hablad juntos: ¿qué significa para vosotros que el amor no sea solo un sentimiento, sino algo inscrito en el cuerpo?
\end{enumerate}

\vinetacapitulo{ch01}

\chapter{Un cuerpo que habla}

\epigraph{El primer lenguaje que aprende un ser humano no tiene palabras. Se habla con la piel.}{}

En el capítulo anterior vimos que nacemos radicalmente inmaduros ---un feto externo, como decía Montagu--- y que esa vulnerabilidad no es un fallo sino el precio de nuestro cerebro extraordinario. Vimos también, casi de pasada, que la evolución no nos dejó solos ante esa fragilidad: preparó mecanismos que empujan a cuidar.

Ahora toca detenernos en esos mecanismos. Porque lo que ocurre en el cuerpo de una madre durante el parto, lo que le pasa al padre cuando sostiene a su hijo por primera vez, lo que sucede entre la piel de un recién nacido y la piel de quien lo abraza, es mucho más que biología. Es un lenguaje. Un lenguaje que nuestro cuerpo habla antes de que nosotros hayamos aprendido una sola palabra.

\section{La inundación}

En las últimas fases del parto, el cuerpo de la madre hace algo extraordinario. El hipotálamo, una estructura pequeña y antigua en la base del cerebro, comienza a liberar oxitocina en cantidades masivas.\footnote{Los niveles de oxitocina en plasma durante el parto pueden ser tres o cuatro veces superiores a los niveles basales, con picos aún mayores en el momento de la expulsión. Cf. Uvnäs-Moberg, K., \textit{The Oxytocin Factor}, Da Capo Press, 2003.} No es una liberación gradual. Es una inundación.

La oxitocina hace que el útero se contraiga ---esa es su función más conocida, la que aprovecha la medicina cuando induce un parto---. Pero hace mucho más. Actúa en el cerebro como un interruptor emocional. Reduce la ansiedad. Genera una sensación profunda de calma y confianza. Y, sobre todo, activa los circuitos cerebrales del vínculo: esos que hacen que una madre mire a la criatura que acaba de salir de su cuerpo y sienta algo que ninguna palabra describe del todo. Algo que no es exactamente decisión, ni exactamente instinto, sino una mezcla poderosa de ambos.

El torrente no se detiene con el parto. La lactancia lo mantiene. Cada vez que el bebé succiona, se libera más oxitocina.\footnote{La succión del pezón estimula las terminaciones nerviosas que envían señales al hipotálamo, desencadenando la liberación pulsátil de oxitocina. Este mecanismo tiene una doble función: provoca la eyección de leche (el ``reflejo de bajada'') y simultáneamente refuerza el vínculo emocional madre-hijo. Cf. Rilling y Young, ``The biology of mammalian parenting and its effect on offspring social development'', \textit{Science}, 2014.} No solo para que la leche fluya ---que también---, sino para que el vínculo se refuerce. La naturaleza, si se me permite la expresión, no se fía de una sola dosis. Insiste. Cada toma es un recordatorio biológico: \textit{esta criatura es tuya, cuídala}.

Y hay algo más. Durante el parto y la lactancia también se liberan endorfinas ---los opiáceos naturales del cuerpo--- que generan una sensación de bienestar profundo.\footnote{Cf. Matthiesen et al., ``Postpartum maternal oxytocin release by newborns: effects of infant hand massage and sucking'', \textit{Birth}, 2001.} Es como si el cuerpo recompensara a la madre por el esfuerzo del parto y por el acto de amamantar. Como si le dijera: \textit{esto que estás haciendo es lo correcto, sigue así}.

No es magia. Es química. Pero es una química que apunta con una precisión asombrosa hacia algo que todos reconocemos como profundamente humano: el amor de una madre por su hijo.

\section{La piel que siente}

Hay un momento que los obstetras conocen bien. Inmediatamente después del parto, si se coloca al recién nacido desnudo sobre el pecho desnudo de su madre ---lo que se llama contacto piel con piel---, ocurre algo notable.

El bebé, que hace un instante estaba llorando, se calma. Su ritmo cardíaco se estabiliza. Su temperatura corporal se regula ---la piel de la madre ajusta su temperatura para calentar o enfriar al bebé según lo necesite---.\footnote{Este fenómeno se conoce como ``sincronía térmica''. El pecho de la madre puede variar hasta dos grados centígrados para ajustarse a las necesidades térmicas del recién nacido. Cf. Ludington-Hoe, S., ``Skin-to-skin contact: A calming and reassuring experience for preterm infants'', \textit{Neonatal Network}, 2004.} Sus niveles de cortisol ---la hormona del estrés--- bajan drásticamente. Si se le deja ahí, sin prisa, el bebé es capaz de reptar lentamente hacia el pecho y encontrar el pezón por sí solo, guiado por el olfato.\footnote{Este comportamiento, documentado en vídeo por investigadores suecos en los años noventa, se conoce como \textit{breast crawl}. El recién nacido, colocado sobre el abdomen de la madre, es capaz de desplazarse hacia el pecho y comenzar a mamar sin ninguna ayuda externa en los primeros 60-90 minutos de vida. Cf. Widström et al., ``Newborn behaviour to locate the breast when skin-to-skin'', \textit{Acta Paediatrica}, 2011.}

No es solo el bebé el que responde. La madre también. El contacto piel con piel dispara en ella una nueva oleada de oxitocina. Activa circuitos cerebrales que los neurocientíficos asocian con la recompensa y la motivación.\footnote{Estudios con resonancia magnética funcional muestran que el contacto con el recién nacido activa en la madre el área tegmental ventral y el núcleo accumbens ---los mismos circuitos que se activan con otras experiencias de recompensa intensa---. Cf. Strathearn et al., ``Adult attachment predicts maternal brain and oxytocin response to infant cues'', \textit{Neuropsychopharmacology}, 2009.} En términos llanos: tocar a su bebé le produce placer. Un placer que no es frívolo sino profundamente funcional, porque la empuja a seguir en contacto, a seguir cuidando.

La ciencia ha medido estos efectos con rigor, y los resultados son consistentes: los bebés que reciben contacto piel con piel inmediato lloran menos, duermen mejor, mantienen niveles de glucosa más estables y tienen más éxito con la lactancia. Los prematuros que reciben el llamado ``método canguro'' ---contacto piel con piel prolongado--- ganan peso más rápido, regulan mejor la temperatura y tienen menos infecciones.\footnote{El método canguro fue desarrollado en Colombia en 1978 por los doctores Rey y Martínez como respuesta a la escasez de incubadoras. Los resultados fueron tan buenos que hoy la OMS lo recomienda como estándar de cuidado para prematuros en todo el mundo. Una revisión Cochrane de 2016 concluyó que reduce la mortalidad neonatal en un 40\%. Cf. Conde-Agudelo y Díaz-Rossello, ``Kangaroo mother care to reduce morbidity and mortality in low birthweight infants'', \textit{Cochrane Database of Systematic Reviews}, 2016.}

La piel, ese órgano enorme que nos envuelve, resulta ser mucho más que una barrera protectora. Es un órgano de comunicación. El primer canal por el que un ser humano recibe un mensaje, antes de entender una sola palabra, antes de reconocer una cara, es la piel. Y el mensaje que recibe es: \textit{estás a salvo, hay alguien aquí}.

\section{El cerebro del padre}

Pero este capítulo quedaría incompleto ---y sería injusto--- si se limitara a la madre. Porque uno de los hallazgos más reveladores de la neurociencia reciente es que el cerebro del padre también cambia.

Ya lo mencionamos brevemente en el capítulo anterior: cuando un hombre se convierte en padre, sus niveles de testosterona bajan y los de oxitocina suben. Pero el cambio va más allá de las hormonas. Estudios con neuroimagen muestran que la paternidad reorganiza literalmente el cerebro masculino.\footnote{Cf. Abraham et al., ``Father's brain is sensitive to childcare experiences'', \textit{PNAS}, 2014. Este estudio comparó la actividad cerebral de madres, padres heterosexuales y padres homosexuales primarios, y encontró que la implicación activa en el cuidado ---más que el sexo biológico--- era lo que determinaba la activación de los circuitos parentales.}

En los padres que se implican en el cuidado, se activan regiones asociadas con la empatía, la lectura emocional y la vigilancia. La amígdala ---que procesa las amenazas--- se sensibiliza al llanto del bebé. El córtex prefrontal ---que planifica y toma decisiones--- se conecta más intensamente con las áreas emocionales. Es como si el cerebro del hombre se recableara para incorporar una nueva prioridad: \textit{proteger a esta criatura}.

Y hay un dato que me parece especialmente elocuente. La bajada de testosterona no es un accidente ni un efecto secundario indeseable. Es proporcional a la implicación: cuanto más tiempo pasa un padre con su bebé, más baja la testosterona.\footnote{Cf. Gettler et al., ``Longitudinal evidence that fatherhood decreases testosterone in human males'', \textit{PNAS}, 2011. Los padres que dormían en la misma habitación que sus hijos mostraban los niveles más bajos de testosterona.} No es que ser padre ``debilite'' al hombre, como cierta cultura querría hacernos creer. Es que lo reconfigura. La testosterona alta es útil para competir, para buscar pareja, para enfrentarse a rivales. Pero cuando ya tienes pareja y tienes una cría vulnerable, lo que necesitas no es competición: es presencia, atención, cuidado. Y el cuerpo lo sabe.

Hay aquí algo que debería darnos que pensar. El cuerpo del hombre no está diseñado solo para engendrar. Está diseñado también para quedarse. Para cuidar. Para ser padre, no solo progenitor. La biología no nos convierte en padres automáticamente ---eso sería determinismo, y los seres humanos no somos máquinas---. Pero nos pone el viento a favor. Nos da un empujón en la dirección correcta, igual que se lo da a la madre.

\section{Un cuerpo configurado para la relación}

Si damos un paso atrás y miramos el panorama completo, lo que encontramos es impresionante.

Tenemos un parto que libera hormonas del vínculo. Una lactancia que las mantiene. Una piel que, al tocar otra piel, dispara cascadas de bienestar y confianza. Un cerebro materno que se reorganiza para la vigilancia y el cuidado. Un cerebro paterno que hace lo mismo cuando el padre se implica. Un bebé cuyo llanto activa circuitos de alarma en los adultos de su alrededor. Un rostro infantil ---ojos grandes, frente ancha, mejillas redondas--- que dispara ternura en prácticamente todos los seres humanos.\footnote{El etólogo Konrad Lorenz describió en 1943 el \textit{Kindchenschema} (``esquema infantil''): un conjunto de rasgos faciales que los humanos percibimos como ``tiernos'' y que activan conductas de cuidado. Estudios posteriores han confirmado que estos rasgos provocan respuestas cerebrales de recompensa incluso en personas que no son padres. Cf. Lorenz, K., ``Die angeborenen Formen möglicher Erfahrung'', \textit{Zeitschrift für Tierpsychologie}, 1943.}

Todo esto no es casualidad. Es un sistema. Un sistema en el que cada pieza encaja con las demás para producir un resultado: que los adultos cuiden de una cría que no puede cuidarse sola.

Podemos explicarlo en términos evolutivos, y la explicación es sólida: los grupos donde estos mecanismos funcionaban mejor criaban más hijos que llegaban a la edad adulta, y esos hijos transmitían los mismos mecanismos. Selección natural. Pero la explicación evolutiva, siendo verdadera, no agota el significado de lo que vemos. Porque lo que vemos es un cuerpo que no está diseñado para la soledad. Un cuerpo que habla, con cada una de sus fibras, un lenguaje de relación.

Y ese lenguaje no se limita a la relación madre-hijo o padre-hijo. Como vamos a ver ahora, marca para siempre la forma en que nos relacionamos con todos los demás.

\section{Bowlby y el descubrimiento del apego}

En 1950, la Organización Mundial de la Salud encargó a un psiquiatra británico llamado John Bowlby un informe sobre la salud mental de los niños sin hogar en la Europa de posguerra.\footnote{Bowlby, J., \textit{Maternal Care and Mental Health}, OMS, 1951. Este informe fue el germen de lo que después se convertiría en la teoría del apego, desarrollada en su trilogía \textit{Attachment and Loss} (1969, 1973, 1980).} Lo que Bowlby encontró cambió para siempre nuestra forma de entender el desarrollo humano.

Miles de niños habían sido separados de sus padres durante la guerra ---evacuados, huérfanos, internados en instituciones---. Muchos de ellos habían recibido alimento suficiente, abrigo, atención médica. Pero algo les faltaba. No crecían bien. Algunos dejaban de hablar, se balanceaban compulsivamente, se negaban a comer. Otros se mostraban incapaces de establecer relaciones con los cuidadores. Las tasas de mortalidad en los orfanatos eran alarmantemente altas, incluso cuando las condiciones materiales eran aceptables.\footnote{El pediatra René Spitz documentó este fenómeno en su investigación sobre lo que llamó ``hospitalismo'': niños institucionalizados que, pese a recibir alimentación y cuidado médico adecuados, desarrollaban un deterioro físico y psicológico severo por falta de contacto humano. Algunos morían. Cf. Spitz, R., ``Hospitalism: An inquiry into the genesis of psychiatric conditions in early childhood'', \textit{The Psychoanalytic Study of the Child}, 1945.}

Bowlby comprendió algo que hoy nos parece evidente pero que entonces era revolucionario: \textbf{un niño no necesita solo alimento y protección. Necesita una relación}. Necesita una figura estable, predecible, que responda a sus llamadas. Necesita lo que Bowlby llamó una \textit{base segura}: alguien a quien agarrarse cuando tiene miedo, desde cuya cercanía puede explorar el mundo.

Y esto no era un lujo emocional. Era una necesidad biológica. Bowlby lo planteó en términos evolutivos: en la sabana donde evolucionamos, un bebé separado de su cuidador era un bebé muerto. El llanto, la búsqueda de proximidad, la angustia ante la separación, no eran caprichos ni ``mañas''. Eran un sistema de supervivencia, tallado por millones de años de selección natural.

Ese sistema es lo que hoy llamamos \textbf{apego}.

\section{Lo que aprendemos sin saber que lo aprendemos}

Unos años después de Bowlby, una psicóloga canadiense llamada Mary Ainsworth ideó un experimento elegante para observar el apego en acción.\footnote{El procedimiento, conocido como ``Situación Extraña'' (\textit{Strange Situation}), fue diseñado por Ainsworth a finales de los años sesenta. Se basa en observar la reacción de un niño de entre 12 y 18 meses ante una secuencia de separaciones y reencuentros con su madre en presencia de un desconocido. Cf. Ainsworth et al., \textit{Patterns of Attachment}, Lawrence Erlbaum, 1978.}

La situación era sencilla: una madre entraba con su hijo pequeño en una habitación con juguetes. Un desconocido entraba. La madre salía. El desconocido intentaba consolar al niño. La madre volvía. Los investigadores observaban las reacciones del niño en cada fase.

Lo que Ainsworth descubrió fue que no todos los niños reaccionaban igual. Y que las diferencias no eran aleatorias: tenían que ver con la calidad de la relación con la madre durante los meses anteriores. Ainsworth identificó tres patrones principales:

\subsection{Apego seguro}

La mayoría de los niños ---alrededor del 60-65\%--- mostraban lo que Ainsworth llamó \textit{apego seguro}. Estos niños usaban a la madre como base para explorar: jugaban con confianza mientras ella estaba presente, se angustiaban al verla partir, y al volver a verla, iban hacia ella, se consolaban rápidamente y volvían a jugar.

Estos niños tenían madres que, en general, respondían a sus necesidades de forma consistente: cuando el bebé lloraba, acudían; cuando buscaba contacto, lo daban; cuando necesitaba espacio, lo permitían. No madres perfectas ---eso no existe---. Madres suficientemente atentas.\footnote{El concepto de ``madre suficientemente buena'' (\textit{good enough mother}) fue acuñado por el pediatra y psicoanalista Donald Winnicott, y refleja una idea liberadora: no se necesita la perfección para criar bien. Se necesita presencia, atención y una respuesta razonablemente consistente. Cf. Winnicott, D. W., \textit{Playing and Reality}, Tavistock, 1971.}

\subsection{Apego inseguro}

Otros niños mostraban patrones distintos. Algunos, llamados \textit{evitativos}, parecían indiferentes a la partida de la madre y la ignoraban al volver. No es que no les importara ---sus niveles de cortisol estaban tan altos como los de los demás---, sino que habían aprendido que expresar la necesidad no servía de nada. Sus madres tendían a ser menos receptivas al contacto físico, más distantes emocionalmente.

Otros, llamados \textit{ambivalentes} o \textit{ansiosos}, se mostraban angustiadísimos cuando la madre se iba, pero al volver no se consolaban: se aferraban a ella con furia, llorando, resistiéndose a ser calmados. Habían aprendido que la respuesta de la madre era impredecible ---a veces atenta, a veces ausente--- y habían desarrollado una estrategia de hiperactivación: gritar más fuerte, llorar más tiempo, para asegurarse de que alguien respondiera.\footnote{Investigaciones posteriores añadieron un cuarto patrón, el apego \textit{desorganizado}, observado en niños cuyo cuidador era al mismo tiempo fuente de consuelo y fuente de miedo (como en casos de maltrato). Cf. Main y Solomon, ``Discovery of a new, insecure-disorganized/disoriented attachment pattern'', en Brazelton y Yogman (eds.), \textit{Affective Development in Infancy}, Ablex, 1986.}

\section{Del bebé al adulto}

Aquí viene lo que de verdad importa para quienes estáis leyendo este libro pensando en vuestro matrimonio.

Los estilos de apego no se quedan en la infancia. Nos acompañan. Investigaciones de las últimas décadas han mostrado que los patrones que aprendemos en los primeros años ---esa forma de estar con el otro, esa expectativa sobre si los demás van a responder cuando los necesitemos--- tienden a reproducirse en las relaciones adultas, incluida la relación de pareja.\footnote{Cf. Hazan y Shaver, ``Romantic love conceptualized as an attachment process'', \textit{Journal of Personality and Social Psychology}, 1987. Este artículo pionero trasladó la teoría del apego al ámbito de las relaciones adultas y abrió una de las líneas de investigación más fructíferas de la psicología de la pareja.}

Una persona con apego seguro tiende a sentirse cómoda con la intimidad. Puede pedir ayuda sin sentirse débil. Puede dar espacio al otro sin sentirse abandonada. Confía, en general, en que el otro va a estar ahí.

Una persona con apego evitativo tiende a la independencia excesiva. Le cuesta pedir, le cuesta mostrar vulnerabilidad, se retira emocionalmente cuando las cosas se ponen intensas. No es que no necesite al otro; es que aprendió, muy pronto, que mostrar la necesidad era arriesgado.

Una persona con apego ansioso tiende a la hipervigilancia emocional. Necesita confirmación constante de que el otro la quiere. Un mensaje que tarda en llegar se convierte en fuente de angustia. La distancia se interpreta como rechazo. No es que sea ``dependiente'' por elección; es que su sistema de apego aprendió que la disponibilidad del otro era incierta, y se puso en alerta permanente.

Pero ---y esto es fundamental--- \textbf{el apego no es una condena}. Los estilos de apego son tendencias, no destinos. Se pueden conocer, comprender y, con esfuerzo y con la relación adecuada, modificar. De hecho, una relación de pareja segura es uno de los contextos más poderosos para sanar heridas de apego.\footnote{El concepto de ``apego ganado'' (\textit{earned secure attachment}) describe precisamente a personas que, pese a haber tenido experiencias de apego inseguro en la infancia, desarrollan un estilo seguro en la edad adulta ---a menudo a través de una relación de pareja estable, terapia, u otras relaciones significativas---. Cf. Roisman et al., ``Earned-security in retrospect'', \textit{Development and Psychopathology}, 2002.} Cuando alguien te responde de forma constante, cuando está ahí día tras día, tu sistema de apego ---ese sistema antiguo, tallado para la supervivencia--- empieza a relajarse. Empieza a confiar. Empieza a creer que puede soltar la guardia.

\section{Lo que el cuerpo nos dice}

Si recorremos todo lo que hemos visto en este capítulo ---la oxitocina del parto, la piel que comunica, el cerebro del padre que se transforma, el sistema de apego que nos configura--- lo que emerge es una imagen coherente y asombrosa.

Nuestro cuerpo no es un accidente. No es un vehículo neutro que nos transporta por la vida. Es un cuerpo que habla. Que dice, con cada mecanismo, con cada hormona, con cada circuito neuronal: \textit{estás hecho para la relación. Estás hecho para el vínculo. Estás hecho para cuidar y para ser cuidado}.

La biología no es todo, por supuesto. Si lo fuera, bastaría con la química para explicar el amor, y sabemos que no basta. Hay algo en el amor que desborda la biología, algo que le da sentido, algo que la eleva. Algo que hace que una madre que adopta a un niño ---sin el torrente hormonal del parto--- pueda amarlo con la misma intensidad. Algo que hace que un padre al que la biología no le ha puesto las cosas fáciles elija quedarse cada mañana.

Pero la biología es el suelo sobre el que se construye todo lo demás. Y es un suelo que tiene una orientación clara. Como si alguien ---o algo--- hubiera querido que fuera así.

A eso ---a quién o qué hay detrás de ese diseño--- llegaremos más adelante. Por ahora, basta con escuchar lo que el cuerpo dice. Y lo que dice es bastante claro.

\paracaminar

\begin{enumerate}
  \item Conocer los estilos de apego puede ayudar a entenderse mejor como pareja. ¿Os reconocéis en alguno de ellos? ¿Cómo reaccionáis cada uno cuando sentís que el otro se distancia o cuando necesitáis cercanía?

  \item Hablad juntos sobre vuestra experiencia con el contacto físico ---no solo el sexual, sino el cotidiano: la mano en el hombro, el abrazo al llegar, dormir cerca---. ¿Qué papel tiene en vuestra relación? ¿Os resulta fácil o difícil?

  \item Si el cuerpo está ``configurado para la relación'', ¿qué consecuencias tiene eso para la forma en que queréis construir vuestra familia?
\end{enumerate}

\vinetacapitulo{ch02}

\chapter{La química del vínculo}

\epigraph{El enamoramiento es la forma que tiene la naturaleza de que dos personas se miren el tiempo suficiente como para decidir quedarse.}{Helen Fisher}

Casi todo el mundo reconoce la experiencia, aunque no sepa explicarla. Te ha pasado, o te está pasando, o conoces a alguien a quien le pasa: de repente, una persona que antes era una más entre miles ocupa el centro de todo. No puedes dejar de pensar en ella. La recuerdas al despertar y al acostarte. Cada mensaje suyo te produce un pequeño estallido de alegría. Comes menos, duermes peor, y sin embargo te sientes lleno de energía. El mundo entero parece más luminoso, más intenso, más lleno de sentido.

Estás enamorado.

Y lo que sientes no es una metáfora. Es química. Literalmente.

\section{El cerebro enamorado}

En el año 2005, la antropóloga Helen Fisher y su equipo publicaron un estudio que marcó un antes y un después en la comprensión científica del amor.\footnote{Fisher, H., Aron, A., y Brown, L.L., ``Romantic love: a mammalian brain system for mate choice'', \textit{Philosophical Transactions of the Royal Society B}, 2006. Fisher llevaba décadas estudiando el amor desde la antropología y la neurociencia, y es autora de varios libros sobre el tema, entre ellos \textit{Why We Love} (2004) y \textit{Anatomy of Love} (1992).} Tomaron a un grupo de personas que declaraban estar intensamente enamoradas, las metieron en un escáner de resonancia magnética funcional, y les mostraron fotografías de la persona amada. Lo que vieron en las pantallas fue revelador.

El cerebro de una persona enamorada se ilumina de una forma muy específica. No se activan principalmente las zonas de la emoción, como cabría esperar. Lo que se enciende con fuerza es el \textit{área tegmental ventral} y el \textit{núcleo caudado}: regiones profundas del cerebro, muy antiguas evolutivamente, que forman parte del llamado sistema de recompensa.\footnote{El sistema de recompensa es un circuito cerebral que refuerza las conductas necesarias para la supervivencia ---comer, beber, reproducirse--- generando una sensación de placer que nos impulsa a repetirlas. Que el enamoramiento active este sistema sugiere que, para el cerebro, el vínculo de pareja es tan esencial como alimentarse.} Son las mismas zonas que se activan ante una recompensa intensa, ante algo que el cerebro interpreta como vital, como necesario. Como algo que no puede dejar pasar.

Y el combustible de esa activación tiene nombre: \textbf{dopamina}.

\section{El cóctel de la pasión}

La dopamina es un neurotransmisor ---una molécula que las neuronas usan para comunicarse entre sí--- asociado a la motivación, la búsqueda de recompensa y el deseo. Cuando estamos enamorados, el cerebro libera dopamina en cantidades muy superiores a lo habitual. Es la dopamina la que produce esa sensación de euforia, esa energía inagotable, ese enfoque casi obsesivo en la persona amada.\footnote{Los niveles de dopamina en el enamoramiento se han comparado con los que producen ciertas drogas estimulantes. No es una exageración: el sistema de recompensa activado es el mismo. La diferencia es que el enamoramiento activa el circuito de forma natural y con un propósito biológico. Cf. Fisher, H., \textit{Why We Love: The Nature and Chemistry of Romantic Love}, Henry Holt, 2004.}

Pero la dopamina no actúa sola. Junto a ella sube la \textbf{noradrenalina}, que pone al cuerpo en estado de alerta ---de ahí las manos sudorosas, el corazón acelerado, la dificultad para dormir, la pérdida de apetito---. Y ocurre algo más, quizá lo más revelador: baja la \textbf{serotonina}.

La serotonina es el neurotransmisor de la calma, de la regulación emocional, del equilibrio. Cuando sus niveles descienden, aparecen los pensamientos intrusivos, la rumiación, la incapacidad de dejar de pensar en algo. ¿Les suena? Un estudio célebre de la psiquiatra Donatella Marazziti, de la Universidad de Pisa, comparó los niveles de serotonina de personas recién enamoradas con los de pacientes diagnosticados de trastorno obsesivo-compulsivo. El resultado fue asombroso: eran prácticamente iguales.\footnote{Marazziti, D., Akiskal, H.S., Rossi, A., y Cassano, G.B., ``Alteration of the platelet serotonin transporter in romantic love'', \textit{Psychological Medicine}, 1999. Los enamorados presentaban niveles de serotonina un 40\% más bajos que los del grupo de control, similares a los de pacientes con TOC.}

Sí: el enamoramiento, bioquímicamente hablando, se parece a una obsesión.

No dejas de pensar en esa persona. La idealizas. Repasas mentalmente cada conversación, cada gesto, cada palabra. Cualquier estímulo ---una canción, un olor, una calle--- te la trae a la memoria. Tu capacidad de concentrarte en otras cosas disminuye notablemente. Si alguien te dijera que te comportas como un obsesivo, te parecería exagerado. Pero tu cerebro, si pudiera hablar, confirmaría el diagnóstico.

\section{Una locura con sentido}

Ahora bien: antes de alarmarnos, conviene hacer una pregunta. ¿Por qué? ¿Por qué la evolución diseñó un estado tan intenso, tan absorbente, tan parecido a la locura?

La respuesta tiene que ver con algo que ya vimos en el primer capítulo: nuestros bebés nacen extraordinariamente vulnerables. Necesitan años de cuidado. Y para que una cría humana sobreviva, no basta con una madre: hace falta una pareja que se comprometa. Que se quede. Que colabore.

El enamoramiento es el mecanismo que la naturaleza utiliza para que dos personas se elijan mutuamente y se mantengan juntas el tiempo suficiente como para establecer un vínculo profundo. Es una especie de turbo biológico: durante unos meses ---a veces algo más de un año, raramente más de tres--- el cerebro concentra toda su maquinaria motivacional en una sola persona. Te hace sentir que esa persona es única, insustituible, la más importante del mundo. Y en cierto modo tiene razón: porque lo que está en juego, evolutivamente, es la supervivencia de la especie.

Helen Fisher propuso que el amor humano funciona en realidad como tres sistemas cerebrales interconectados pero distintos:\footnote{Fisher, H., ``Lust, Attraction, and Attachment in Mammalian Reproduction'', \textit{Human Nature}, 1998. Fisher desarrolló ampliamente este modelo en \textit{Why We Love} (2004) y \textit{Anatomy of Love} (edición revisada, 2016).}

\begin{itemize}
    \item El \textbf{deseo sexual} (\textit{lust}), impulsado fundamentalmente por las hormonas sexuales ---testosterona y estrógenos---, que nos orienta hacia la búsqueda de pareja en general.
    \item La \textbf{atracción romántica} (\textit{attraction}), mediada por la dopamina y la noradrenalina, que concentra nuestra atención en una persona concreta y genera el estado que llamamos enamoramiento.
    \item El \textbf{apego de pareja} (\textit{attachment}), sostenido por la oxitocina y la vasopresina, que produce la sensación de calma, seguridad y conexión profunda con esa persona a largo plazo.
\end{itemize}

Los tres sistemas pueden activarse juntos, pero no siempre lo hacen. Puedes sentir deseo por alguien sin estar enamorado. Puedes estar enamorado de alguien y sentir deseo por otro. Puedes sentir apego profundo por tu pareja de muchos años sin experimentar ya la tormenta del enamoramiento inicial. La vida sentimental humana es compleja precisamente porque estos tres sistemas no siempre van de la mano.

Pero cuando los tres coinciden en una misma persona ---cuando deseas, te enamoras y te vinculas con la misma persona---, algo extraordinario ocurre. Y es ahí donde entra la segunda parte de la historia.

\section{De la tormenta a la calma}

El enamoramiento no dura para siempre. No puede durar para siempre. Un estado de activación tan intenso consume una cantidad enorme de recursos. Si el cerebro mantuviera indefinidamente los niveles de dopamina del enamoramiento, acabaríamos agotados, incapaces de concentrarnos en otra cosa, funcionando al límite.\footnote{Diversos estudios estiman que la fase de enamoramiento intenso dura entre 12 y 18 meses, aunque con variaciones individuales importantes. Cf. Tennov, D., \textit{Love and Limerence: The Experience of Being in Love}, Stein and Day, 1979.}

Y aquí es donde muchas parejas se asustan. Porque cuando la tormenta de dopamina empieza a amainar, cuando ya puedes dormir con normalidad y recuperas el apetito, cuando ya no sientes ese vuelco en el estómago cada vez que ves a tu pareja, es tentador pensar que el amor se ha terminado. Que ``se ha apagado la llama''. Que ya no sientes lo mismo.

Pero lo que ocurre es algo muy distinto. Lo que ocurre es que el cerebro está haciendo una transición. Está pasando del sistema de atracción al sistema de apego. Y las protagonistas ahora son otras moléculas: la \textbf{oxitocina} y la \textbf{vasopresina}.

Ya encontramos la oxitocina en el capítulo anterior, vinculada al parto y la lactancia. Pero su papel va mucho más allá. La oxitocina se libera con el contacto físico ---un abrazo, una caricia, la relación sexual---, y genera sensaciones de confianza, calma y bienestar.\footnote{Cf. Uvnäs-Moberg, K., \textit{The Oxytocin Factor: Tapping the Hormone of Calm, Love, and Healing}, Da Capo Press, 2003.} No produce el estallido eufórico de la dopamina. Produce algo más silencioso pero más estable: la sensación de estar en casa. De estar con la persona adecuada. De seguridad.

La vasopresina, por su parte, está especialmente vinculada a la conducta protectora y al vínculo de pareja en los varones. Uno de los estudios más fascinantes sobre esta molécula se realizó con un animal inesperado: el topillo de las praderas.\footnote{Young, L.J. y Wang, Z., ``The neurobiology of pair bonding'', \textit{Nature Neuroscience}, 2004. Los topillos de las praderas (\textit{Microtus ochrogaster}) son uno de los pocos mamíferos monógamos. Sus primos, los topillos de montaña (\textit{Microtus montanus}), son promiscuos. La diferencia clave está en la distribución de receptores de vasopresina en el cerebro.} Los topillos de las praderas son uno de los pocos mamíferos que forman parejas estables de por vida. ¿La clave? Tienen una alta densidad de receptores de vasopresina en las zonas de recompensa del cerebro. Cuando se les bloquean esos receptores, dejan de formar vínculos de pareja. Cuando se aumentan artificialmente en una especie de topillo promiscuo, este empieza a vincularse con una sola hembra.

No somos topillos, por supuesto. Pero la molécula es la misma, y el mecanismo básico se conserva: la vasopresina facilita la fidelidad y la protección del vínculo.

\section{El amor que permanece}

Lo que la ciencia nos muestra, entonces, es que el amor humano no es un único fenómeno, sino un proceso. Un proceso que empieza con una tormenta y, si todo va bien, desemboca en un puerto.

La tormenta es espectacular. Es intensa, es absorbente, es memorable. Pero es, por definición, temporal. El puerto no tiene la misma adrenalina, pero tiene algo que la tormenta no puede dar: profundidad, estabilidad, la certeza de que hay alguien a tu lado cuando la vida se pone difícil.

Y aquí viene lo verdaderamente interesante: el apego no es una versión rebajada del enamoramiento. No es ``conformarse'' con menos. Es un sistema biológico distinto, diseñado para un propósito distinto. El enamoramiento sirve para elegir; el apego sirve para permanecer. Y ambos son necesarios.

Los estudios de parejas que llevan décadas juntas y que declaran seguir profundamente enamoradas muestran algo notable: en sus cerebros, los circuitos de apego están muy activos, como cabría esperar, pero además siguen mostrando activación en las áreas de recompensa asociadas a la dopamina.\footnote{Acevedo, B.P., Aron, A., Fisher, H.E., y Brown, L.L., ``Neural correlates of long-term intense romantic love'', \textit{Social Cognitive and Affective Neuroscience}, 2012. Este estudio escaneó los cerebros de personas que llevaban una media de 21 años de relación y seguían reportando amor intenso. Las áreas de recompensa seguían activas, pero sin la activación de las áreas de ansiedad y obsesión típicas del enamoramiento temprano.} Es decir: el amor maduro conserva la intensidad, pero pierde la ansiedad. Se parece más a una llama constante que a un incendio.

\section{Más allá de la química}

Sería tentador detenernos aquí y concluir que el amor es, al final, solo una cuestión de moléculas. Que somos marionetas de la dopamina y la oxitocina, y que lo que llamamos amor no es más que un sofisticado programa biológico de vinculación.

Pero hay algo que no encaja en esa explicación. Algo que la biología sola no puede explicar.

Porque la biología nos empuja, sí. Nos pone el viento a favor. Pero no nos obliga. Hay personas que sienten el torbellino del enamoramiento y deciden no seguirlo. Hay parejas que pierden el apego biológico y eligen quedarse. Hay hombres y mujeres que, contra toda la química de su cuerpo, deciden ser fieles a una promesa que hicieron un día. Y hay otros que, con toda la oxitocina del mundo circulando por sus venas, deciden marcharse.

El ser humano es el único animal que puede decirle que no a sus hormonas. El único que puede elegir contra su instinto. Y el único que puede elegir \textit{a favor} de su instinto no porque le arrastre, sino porque libremente decide que esa es la dirección correcta.

Eso es lo que hace al amor humano algo radicalmente distinto al vínculo de pareja de los topillos. No es que nuestra oxitocina sea diferente. Es que nosotros somos más que nuestra oxitocina. Somos capaces de tomar esa corriente biológica que nos empuja hacia el otro y convertirla en una decisión. En un proyecto. En una promesa.

El enamoramiento te dice: ``No puedo dejar de pensar en ti''. El apego te dice: ``Me siento seguro a tu lado''. Pero hay algo más ---algo que no se reduce a la bioquímica--- que dice: ``Te elijo. Hoy. Y mañana. Y pasado mañana. No porque mis hormonas me lo pidan, sino porque quiero construir mi vida contigo''.

Ese ``algo más'' es, quizá, lo más humano que hay en nosotros. La biología nos regala un impulso poderoso. Pero convertir ese impulso en un amor que dura toda la vida requiere algo que ninguna molécula puede proporcionar: la libertad de elegir.

Y aquí se abre una pregunta que merece la pena hacerse. Si nuestro cuerpo está tan finamente diseñado para el vínculo ---con sistemas que nos empujan a elegir, a apegarnos, a permanecer---, y si al mismo tiempo somos libres para acoger o rechazar ese impulso... ¿de dónde viene esa doble dotación? ¿Quién nos hizo capaces de amar \textit{y} de elegir amar?

Pero eso es materia para más adelante. De momento, basta con asombrarse.

\paracaminar

\begin{enumerate}
  \item Recordad juntos los primeros tiempos de vuestra relación: las mariposas, la dificultad para pensar en otra cosa, la euforia. Ahora que sabéis que había dopamina, noradrenalina y serotonina baja detrás de todo aquello, ¿cambia en algo lo que sentisteis? ¿O lo hace más fascinante?

  \item ¿Habéis experimentado ya la transición del enamoramiento al apego? Si es así, ¿cómo la vivisteis? ¿Hubo un momento en que pensasteis que el amor se estaba ``apagando'', o supisteis reconocer que estaba cambiando de forma?

  \item Hablad juntos: ¿qué significa para vosotros pasar de ``no puedo dejar de pensar en ti'' a ``te elijo cada día''? ¿Qué tiene de distinto un amor elegido frente a un amor sentido?
\end{enumerate}

\vinetacapitulo{ch03}

\chapter{El amor no es un cuento}

\epigraph{Amar no es mirarse el uno al otro, sino mirar juntos en la misma dirección.}{Antoine de Saint-Exupéry}

Todos hemos visto la escena. Él y ella se miran a los ojos. Suena la música. El mundo se detiene. Se besan. Aparecen las letras: ``Y vivieron felices para siempre''. Fin.

La película acaba justo donde empieza lo interesante.

Porque si algo hemos aprendido en los capítulos anteriores, es que nuestro cuerpo está extraordinariamente preparado para el encuentro. Tenemos hormonas que nos empujan a vincularnos, circuitos neuronales que nos hacen sentir euforia junto a la persona que nos atrae, mecanismos biológicos que favorecen la permanencia. Todo eso es real, está inscrito en nuestra naturaleza, y es asombroso.

Pero hay un salto enorme entre sentir todo eso y construir un amor que dure. Un salto que las películas se saltan alegremente ---y que la vida, con su habitual falta de banda sonora, nos obliga a dar con los pies en el suelo.

\section{El mito del amor de película}

No es que las películas de amor sean malas. Algunas son estupendas. El problema es cuando convertimos la ficción en modelo de lo que el amor debería ser. Y eso, sin darnos cuenta, lo hacemos más de lo que pensamos.

El guion suele ser parecido. Dos personas se encuentran. Hay una chispa inmediata, un flechazo que ninguno de los dos puede controlar. Quizá hay un obstáculo ---un malentendido, una distancia, un rival--- que se resuelve en el tercer acto con una carrera bajo la lluvia o una declaración épica en un aeropuerto. Y luego: la felicidad. El amor verdadero ha triunfado.

¿Qué nos enseña ese guion, si lo miramos con un poco de distancia?

Primero, que el amor es algo que te \textit{pasa}. Que llega sin que lo busques, te arrolla como un tren y ya está. Si no sientes mariposas en el estómago, si no pierdes el sueño, si no te tiemblan las piernas, entonces no es ``el de verdad''. El amor auténtico se reconoce porque es irresistible.

Segundo, que el amor no debería costar esfuerzo. Si hay que trabajar en la relación, si hay que pedir perdón, si hay que negociar quién friega los platos, algo falla. El amor de verdad fluye. Es natural. Es fácil.

Y tercero ---quizá lo más dañino---, que existe una única persona perfecta para ti en algún lugar del mundo. Tu ``media naranja'', tu alma gemela, la pieza que encaja sin necesidad de limar bordes. Y si un día descubres que la persona que tienes al lado no es exactamente esa pieza... bueno, quizá te equivocaste. Quizá hay que seguir buscando.\footnote{El sociólogo Zygmunt Bauman analiza con lucidez esta mentalidad consumista aplicada a las relaciones en \textit{Amor líquido}, Fondo de Cultura Económica, 2005. Cuando el otro deja de satisfacernos, lo ``devolvemos'' y buscamos un modelo mejor.}

Dicho así, suena absurdo. Pero conviene ser honestos: ¿cuántas relaciones se han roto no porque hubiera un problema grave, sino porque alguien esperaba sentir permanentemente lo que sintió los primeros meses? ¿Cuántas personas han abandonado algo bueno ---algo real--- persiguiendo una intensidad que, por definición, no puede durar para siempre?

\section{Cuando la química se confunde con el amor}

Aquí es donde lo que vimos en el capítulo anterior cobra una importancia decisiva. Recordemos: el enamoramiento tiene un sustrato biológico poderosísimo. Dopamina a raudales, noradrenalina que nos mantiene en alerta, serotonina baja que nos hace pensar obsesivamente en la otra persona. Es un cóctel químico diseñado por la evolución para que dos personas se acerquen, se deseen, quieran estar juntas. Y funciona. Vaya si funciona.

Pero tiene fecha de caducidad.

No es que el sentimiento desaparezca del todo. Es que cambia. La intensidad del inicio ---ese no poder pensar en otra cosa, esa euforia, ese ``me muero si no te veo''--- se atenúa con el tiempo. Y tiene que ser así. Ningún organismo puede mantenerse indefinidamente en ese estado de excitación sin colapsarse. Los estudios sugieren que la fase de enamoramiento intenso dura, en la mayoría de los casos, entre uno y tres años.\footnote{Cf. Fisher, H., \textit{Why We Love: The Nature and Chemistry of Romantic Love}, Henry Holt, 2004. Fisher distingue tres sistemas cerebrales del amor ---impulso sexual, atracción romántica y apego--- con sustratos neuroquímicos diferenciados.}

Y cuando baja la marea ---cuando la dopamina deja de inundar cada encuentro, cuando puedes volver a concentrarte en el trabajo, cuando duermes otra vez de un tirón--- llega el momento de la verdad. Porque ahí es donde muchas personas sienten que algo se ha roto. ``Ya no siento lo mismo'', dicen. ``Se ha acabado la magia.''

Pero lo que se ha acabado no es el amor. Lo que se ha acabado es la fase química del enamoramiento. Son cosas distintas. Tan distintas como el fogonazo que enciende la chimenea y el fuego lento que calienta la casa durante todo el invierno. El fogonazo es espectacular, sí. Pero es el fuego sostenido el que te salva del frío.

Confundir el enamoramiento con el amor es como confundir el hambre con la nutrición. El hambre te lleva a la mesa. Pero lo que te alimenta es comer ---y comer bien, día tras día, no un atracón puntual seguido de ayuno.

\section{Amar es un verbo}

Erich Fromm escribió hace más de medio siglo algo que sigue siendo revolucionario por lo sencillo: el amor no es un sustantivo, es un verbo. No es algo que se \textit{tiene}, sino algo que se \textit{hace}.\footnote{Fromm, E., \textit{El arte de amar}, Paidós, 1956. Fromm argumenta que el amor es un arte que requiere conocimiento, esfuerzo y práctica, igual que cualquier otra actividad humana seria.}

Esto cambia todo.

Si el amor fuera solo un sentimiento ---algo que te pasa, que te invade, que no controlas---, entonces no tendría mucho sentido prometer amor. Sería como prometer que mañana va a hacer sol. Puedes desearlo, pero no depende de ti.

Sin embargo, si el amor es también ---y sobre todo--- una decisión, una acción, algo que haces con tu voluntad libre, entonces sí puedes prometerlo. Porque puedes decidir cuidar al otro aunque hoy no sientas mariposas. Puedes elegir escuchar cuando estás cansado. Puedes optar por el perdón cuando sería más fácil llevar la cuenta de los agravios. Puedes levantarte cada mañana y renovar la decisión de querer a esa persona concreta, con sus virtudes y con sus manías, con su luz y con sus sombras.

Eso no significa que el sentimiento no importe. Claro que importa. El deseo, la ternura, la emoción de estar con alguien que te hace bien: todo eso es parte del amor y sería triste prescindir de ello. Pero el sentimiento solo no basta. Sin la decisión, el sentimiento es como un barco sin timón: va donde lo lleve el viento.

Fijaos en algo curioso. Cuando un padre o una madre dicen ``quiero a mi hijo'', nadie entiende que eso signifique que sienten una emoción agradable cada vez que lo miran. Hay noches en que el niño llora y tú estás agotado. Hay adolescencias que ponen a prueba la paciencia de un santo. Hay momentos en que el amor parental se parece más a la determinación que a la ternura. Y sin embargo, nadie duda de que eso sea amor. Lo es, y del bueno.

¿Por qué, entonces, con el amor de pareja exigimos una montaña rusa emocional permanente? ¿Por qué si un día el corazón no se acelera pensamos que algo va mal?

\section{Las tres caras del amor}

Los griegos, que tenían una palabra para casi todo, usaban al menos tres términos distintos donde nosotros usamos uno solo: ``amor''.\footnote{La distinción es más amplia aún. El griego contaba también con \textit{storgé} (amor familiar, afecto natural) y otros matices. Aquí nos centramos en las tres dimensiones más relevantes para el amor de pareja.} Y esa distinción no es un capricho filológico: ilumina algo que la experiencia confirma.

\subsection*{Eros: el deseo que busca}

El eros es la primera chispa. Es la atracción, el deseo, la pasión. Es esa fuerza que te lleva a querer estar cerca del otro, a tocarlo, a mirarlo. Es el componente del amor que siente el cuerpo antes de que la cabeza pueda explicar lo que pasa.

El eros es bueno. No es algo de lo que avergonzarse ni algo que haya que domesticar como a una bestia. Es parte de nuestra naturaleza, y ya hemos visto que tiene raíces biológicas profundas. Sin eros, probablemente dos personas nunca se acercarían lo suficiente como para descubrir lo que tienen en común.

Pero el eros tiene una limitación. Es, por naturaleza, un amor que busca para sí. Deseo al otro porque me hace sentir bien, porque me completa, porque me hace feliz. No hay nada malo en eso ---¿quién no quiere ser feliz?---, pero si el amor se queda ahí, se vuelve frágil. Porque el día que el otro deje de hacerme sentir bien ---porque está enfermo, o cansado, o simplemente porque los años han pasado y la novedad se ha agotado---, ¿qué queda?

\subsection*{Philia: la amistad que comparte}

El segundo rostro del amor es lo que los griegos llamaban \textit{philia}: la amistad profunda. Es el amor que nace de la vida compartida, de los intereses comunes, del placer de estar juntos no solo por la atracción, sino por la compañía.\footnote{C.\,S. Lewis dedica un capítulo espléndido a la amistad en \textit{Los cuatro amores} (1960). Para Lewis, la amistad es el menos biológico de los amores y, quizá por eso, el más cercano a lo angélico: ``Los amigos se miran el uno al otro; los enamorados se miran a los ojos.''}

Hay parejas que se desean pero no se gustan. Hay parejas con una química explosiva que, si les quitas la cama, no tienen de qué hablar. El eros sin philia es un fuego que consume pero no calienta. La philia es lo que convierte a la persona que amas en tu amigo, tu confidente, tu compañero de viaje. La persona con la que puedes estar en silencio sin que el silencio pese.

Pensadlo un momento: en un matrimonio que dure toda la vida, ¿cuántas horas se pasan conversando, comiendo juntos, decidiendo cosas prácticas, riendo, discutiendo, resolviendo problemas, sentados en el mismo sofá? Muchas más que en cualquier otra actividad. Si esa persona no es también tu amigo, tienes un problema.

La philia crece donde el eros la planta. Pero necesita su propio cuidado: tiempo compartido, intereses cultivados en común, conversación honesta. Parejas que dejan de hablar ---de hablar de verdad, no de logística doméstica--- están dejando morir la philia sin darse cuenta.

\subsection*{Ágape: el amor que se entrega}

Hay un tercer nivel, y es el más difícil. El más exigente. Y, paradójicamente, el que hace que el amor sea verdaderamente grande.

El ágape es el amor como don. No el amor que busca (eros) ni el amor que comparte (philia), sino el amor que se da sin condiciones. Es el amor que dice: ``Quiero tu bien, aunque me cueste. Aunque no me lo agradezcas. Aunque hoy no me apetezca.''

Es la madre que se levanta a las tres de la mañana a atender al bebé. Es el marido que acompaña a su mujer a la quimioterapia sin faltar un día. Es la esposa que sostiene a su marido cuando ha perdido el trabajo y la confianza en sí mismo. No porque sea agradable hacerlo, sino porque el otro lo necesita y yo puedo darlo.

Fromm lo expresó con una fórmula que merece la pena recordar: el amor inmaduro dice ``te amo porque te necesito''; el amor maduro dice ``te necesito porque te amo''.\footnote{Fromm, \textit{El arte de amar}, cap. II. La fórmula es sencilla pero profundamente reveladora: en el amor maduro, la necesidad nace del amor y no al revés.} La diferencia es abismal. En el primer caso, el otro es un medio para satisfacer mi necesidad. En el segundo, mi necesidad nace de haber descubierto en el otro un valor que merece todo lo que puedo dar.

El ágape no anula al eros ni a la philia. Los necesita. Un matrimonio solo de ágape, sin deseo ni amistad, sería sacrificio heroico pero no plenitud. Las tres dimensiones se necesitan como las patas de un trípode: quita una y todo se cae.

Pero es el ágape el que sostiene al amor cuando las otras dos dimensiones pasan por sus inevitables valles. Y todos los matrimonios tienen valles. Es el ágape lo que te lleva a elegir al otro cuando elegirlo cuesta. Y es ahí, precisamente ahí, donde el amor revela su verdadero rostro.

\section{Conocerte para poder darte}

Hay una verdad incómoda que conviene mirar de frente: no se puede dar lo que no se tiene.

Si no me conozco ---si no sé cuáles son mis heridas, mis miedos, mis necesidades profundas, mis mecanismos de defensa---, lo que llamaré ``amor'' será muchas veces otra cosa disfrazada. Será dependencia emocional vestida de entrega. Será miedo a la soledad vestido de compromiso. Será la búsqueda inconsciente de que el otro repare lo que yo no he querido mirar en mí mismo.

Esto no es un reproche. Es una realidad que conviene nombrar antes de prometer algo tan grande como ``te amaré toda la vida''. Porque esa promesa exige un sujeto que sepa quién es y qué está prometiendo.

La madurez afectiva no es un estado al que se llega un buen día y ya está. Es un camino. Pero sí hay señales que indican si estamos en ese camino o si estamos evitándolo. Quien se conoce, sabe distinguir entre lo que necesita y lo que desea. Sabe poner nombre a sus emociones sin que lo arrastren. Es capaz de estar solo sin sentir pánico ---y precisamente por eso es capaz de estar con otro sin asfixiarlo---.\footnote{La psicóloga Harriet Lerner lo resume con claridad: ``Solo puedes ser realmente íntimo con otro si no te has perdido a ti mismo.'' Cf. Lerner, H., \textit{La danza de la intimidad}, Urano, 1990.}

Quien no se conoce, en cambio, tiende a buscar en la pareja lo que debería buscar en sí mismo. Y eso es cargar sobre los hombros del otro un peso que no le corresponde. Nadie puede ser tu terapeuta, tu padre ausente, tu autoestima y tu pareja al mismo tiempo. Es demasiado. Y antes o después, la estructura se resquebraja.

Decimos muchas veces que el matrimonio es ``darse al otro''. Y es verdad. Pero solo puedes darte si te tienes. Solo puedes hacer un regalo de ti mismo si sabes qué hay en el paquete.

\section{La paradoja de la libertad}

Hay una objeción que aparece una y otra vez, especialmente entre los más jóvenes: ``¿Para qué comprometerse? ¿No es mejor mantener las opciones abiertas?''

La lógica parece impecable. Si me comprometo con una persona, me cierro a todas las demás. Si no me comprometo, puedo elegir cada día lo que más me convenga. Máxima libertad, máxima flexibilidad.

Pero la experiencia dice otra cosa. Y conviene escucharla.

Pensemos en alguien que, por miedo a equivocarse, nunca elige una carrera profesional. Va probando un poco de esto, un poco de aquello. Mantiene ``las opciones abiertas''. A los cuarenta años, tiene mucha experiencia superficial y ninguna maestría. No ha profundizado en nada. No es el profesional brillante que podría haber sido, porque serlo habría requerido años de dedicación a una sola cosa.

Con el amor pasa algo parecido. El compromiso no es la tumba de la libertad; es su cauce. Como un río: el agua que se desborda por todas partes no llega a ningún sitio; el agua que fluye entre riberas puede mover molinos, regar campos, llegar al mar.

Cuando eliges a una persona ---cuando dices ``tú, y no otra''---, no estás limitando tu capacidad de amar. La estás concentrando. La estás dirigiendo. Y al hacerlo, puedes llegar a una profundidad que el que va saltando de relación en relación nunca conocerá.

Hay algo más, y es quizá lo más sorprendente. El compromiso genera una forma de libertad que la indecisión no puede dar: la libertad de mostrarte como eres.

Pensad en la diferencia entre una primera cita y una cena con alguien que llevas veinte años queriendo. En la primera cita te muestras en tu mejor versión. Elegiste la camisa perfecta, preparaste tres anécdotas graciosas, te cuidas de no decir nada inconveniente. Estás representando un papel.

Después de veinte años, no. Después de veinte años puedes estar en pijama, con la cara lavada, diciendo tonterías, mostrando tus inseguridades, y sabes que el otro no se va a ir. Esa seguridad ---la certeza de que el otro ha elegido quedarse y va a seguir eligiéndolo--- es lo que te permite dejar de actuar y empezar a vivir.

El filósofo Gabriel Marcel decía que la fidelidad creadora no es la repetición monótona de una decisión pasada, sino la renovación constante de un compromiso vivo.\footnote{Cf. Marcel, G., \textit{Être et avoir}, Aubier, 1935. Marcel distingue entre la fidelidad inerte ---la del que se queda por inercia--- y la fidelidad creadora, que reinventa cada día las razones para quedarse.} No te quedas porque lo prometiste hace años y ``no queda otra''. Te quedas porque cada día descubres nuevas razones para elegir lo que ya elegiste. Y eso ---elegir libremente lo que ya es tuyo--- es una de las formas más puras de la libertad.

\section{Un amor que crece}

Terminemos volviendo al principio. El cuento dice que el amor verdadero es el del flechazo, el del corazón desbocado, el del ``lo supe desde el primer momento''. Y que si no es así, no es de verdad.

La realidad dice otra cosa. La realidad dice que los amores más hondos, los más sólidos, los que dan más calor, son los que se han construido ladrillo a ladrillo. Los que han pasado por crisis y han salido más fuertes. Los que han aprendido a pedir perdón y a concederlo. Los que se han reinventado cuando la vida les cambió las reglas del juego.

Hay una belleza particular en las parejas que llevan décadas juntas. No la belleza de la portada de revista, sino la belleza de algo que ha sido probado por el tiempo y ha resistido. La belleza de dos personas que se miran y, debajo de las arrugas y las canas, siguen viendo al otro que eligieron ---y lo siguen eligiendo hoy.

Eso no lo da un fogonazo. Lo da un fuego que alguien ha tenido la paciencia y el coraje de alimentar, día tras día, año tras año.

El amor no es un cuento. Es algo mucho mejor que un cuento. Porque los cuentos terminan, y el amor verdadero no tiene por qué hacerlo.

\paracaminar

\begin{enumerate}
  \item ¿Reconocéis en vuestra propia historia algún ``mito'' sobre el amor que hayáis tenido que desmontar? ¿Qué esperabais del amor antes de vivirlo y qué habéis descubierto después?

  \item De las tres dimensiones del amor ---eros, philia, ágape---, ¿cuál sentís que está más presente en vuestra relación ahora mismo? ¿Cuál necesita más cuidado?

  \item ¿En qué momento concreto habéis experimentado que elegir al otro ---cuando costaba, cuando no apetecía--- ha hecho crecer vuestro amor en vez de reducirlo?
\end{enumerate}

\vinetacapitulo{ch04}

\chapter{Dos, y no uno solo}

\epigraph{Lo contrario del amor no es el odio, sino la indiferencia. Y la indiferencia solo es posible ante lo que no me interpela, ante lo que es igual a mí.}{Adaptado de Elie Wiesel}

Pocos temas generan tanta tensión hoy como la diferencia entre el varón y la mujer. Y se entiende. Durante siglos, esa diferencia se utilizó para justificar la subordinación de la mujer: si somos ``complementarios'', decía el argumento, cada uno tiene su lugar ---y el lugar de ella resultaba ser siempre el segundo---. Cualquier persona honesta tiene que reconocer que eso ocurrió, que hizo mucho daño, y que la reacción contra esa injusticia ha sido en gran parte legítima.

Pero hay una pregunta que merece la pena hacerse con calma, sin trincheras y sin eslóganes: ¿y si el problema no fuera la diferencia, sino lo que hicimos con ella? ¿Y si la complementariedad, bien entendida, no tuviera nada que ver con la desigualdad, sino con algo mucho más hondo ---con la posibilidad misma del encuentro---?

Este capítulo no pretende zanjar un debate que sigue abierto. Pretende algo más modesto: mirar la diferencia sexual con ojos limpios. No como un campo de batalla, sino como un lugar de asombro.

\section{Lo que la complementariedad no es}

Conviene empezar por despejar el terreno. Cuando aquí hablamos de complementariedad no estamos diciendo que el varón sea la pieza A y la mujer la pieza B de un puzzle donde cada uno tiene una función predeterminada. No estamos diciendo que ella sea la emoción y él la razón, que ella cuide y él provea, que ella sienta y él decida. Esos moldes rígidos han causado mucho sufrimiento y, además, son falsos: cualquiera que conozca a hombres y mujeres reales sabe que la realidad es infinitamente más rica que el estereotipo.

Tampoco estamos diciendo que la diferencia sexual sea un sistema jerárquico donde uno vale más que el otro. La igual dignidad del varón y la mujer no es una concesión moderna ni un gesto de cortesía: es un dato antropológico fundamental. Somos igualmente humanos, igualmente dignos, igualmente capaces de razón, de libertad y de amor.\footnote{La igualdad radical en dignidad entre varón y mujer es afirmada hoy tanto desde la antropología filosófica como desde las ciencias sociales. Desde una perspectiva creyente, el Génesis lo expresa con nitidez: ``varón y mujer los creó'' (Gn 1,27), ambos imagen de Dios, sin jerarquía alguna en el acto creador.}

Lo que sí estamos diciendo es algo diferente: que el varón y la mujer, siendo iguales en dignidad, no son idénticos. Y que esa no-identidad, lejos de ser un problema, es una riqueza. Quizá la mayor riqueza de la condición humana.

\section{Un dato que precede a la cultura}

Vivimos en un tiempo que tiende a pensar que toda diferencia entre hombres y mujeres es cultural, construida, y por tanto modificable a voluntad. Es una postura comprensible: si las diferencias son solo culturales, eliminar la desigualdad es cuestión de cambiar la cultura. Y algo de razón tiene: muchas de las diferencias que durante siglos se consideraron ``naturales'' ---que las mujeres no podían estudiar, votar, o dirigir--- eran efectivamente culturales y, con justicia, han sido superadas.

Pero la ciencia nos invita a no simplificar demasiado. Porque junto a las diferencias que son culturales hay otras que no lo son ---o al menos, no del todo---.

Empecemos por lo más básico. Cada célula del cuerpo de un varón tiene un par de cromosomas XY; cada célula de una mujer, XX. Esa diferencia genética no se limita a los órganos reproductivos: influye en la expresión de miles de genes a lo largo de todo el organismo.\footnote{Cf. Arnold, A. P., ``A general theory of sexual differentiation'', \textit{Journal of Neuroscience Research}, 2017. Arnold argumenta que los efectos de los cromosomas sexuales van mucho más allá de las gónadas y afectan a prácticamente todos los tejidos del cuerpo.} Durante el desarrollo embrionario, la exposición a diferentes niveles de hormonas ---especialmente testosterona y estrógenos--- organiza el cerebro de formas que, sin ser determinantes, son significativas.

Los estudios de neurociencia muestran diferencias estadísticas entre cerebros masculinos y femeninos: en el tamaño relativo de ciertas estructuras, en la conectividad entre hemisferios, en la densidad de receptores de determinados neurotransmisores.\footnote{El estudio más citado es el de Ingalhalikar et al., ``Sex differences in the structural connectome of the human brain'', \textit{PNAS}, 2014, que encontró patrones de conectividad distintos entre cerebros masculinos y femeninos. Es un hallazgo sugerente pero debatido: otros investigadores han señalado que los tamaños de efecto son pequeños y que el solapamiento entre individuos de ambos sexos es mayor que la diferencia entre las medias. Para una revisión más amplia y matizada, cf. Halpern, D.\,F., \textit{Sex Differences in Cognitive Abilities}, Psychology Press, 2012, donde se concluye que las diferencias cognitivas entre sexos son reales pero modestas, parcialmente biológicas y parcialmente moldeadas por la experiencia.} El perfil hormonal es diferente a lo largo de toda la vida: los niveles de testosterona, estrógenos, progesterona y oxitocina varían notablemente entre varones y mujeres, y esas hormonas influyen en la percepción del dolor, en la respuesta al estrés, en los patrones de sueño, en la forma de procesar las emociones.\footnote{Cf. Taylor et al., ``Biobehavioral responses to stress in females: Tend-and-befriend, not fight-or-flight'', \textit{Psychological Review}, 2000. Este influyente artículo mostró que la respuesta al estrés en mujeres tiende a seguir un patrón de ``cuidar y buscar alianzas'' (\textit{tend and befriend}), mediado por la oxitocina, a diferencia del clásico modelo masculino de ``luchar o huir''.}

Nada de esto significa que los hombres sean ``de Marte'' y las mujeres ``de Venus''. La variabilidad dentro de cada sexo es enorme, y la cultura moldea poderosamente la expresión de estas tendencias. Un hombre puede ser profundamente empático y una mujer enormemente analítica, sin que ninguno de los dos sea menos hombre o menos mujer por ello. Las diferencias biológicas son tendencias estadísticas, no destinos individuales.

Pero sería igualmente ingenuo negar que esas tendencias existen. Y lo más interesante no es que existan, sino para qué existen.

Porque la diferencia entre varón y mujer no se agota en los cromosomas, las hormonas o las conexiones neuronales. Todo eso es real, pero es el signo, no el significado. Lo que esos datos biológicos están señalando es algo más profundo: que ser varón y ser mujer son dos maneras distintas de ser persona. No solo dos cuerpos diferentes, sino dos formas de estar en el mundo, de percibir, de abrirse a la realidad. Y cuando esas dos formas se encuentran de verdad, ocurre algo que la soledad ---por rica que sea--- no puede dar: cada uno descubre en el otro dimensiones de lo humano a las que, desde sí mismo, no tenía acceso. La diferencia no es un accidente biológico que hay que tolerar; es la condición que hace posible el regalo más grande que una persona puede hacer: entregarse a alguien que es genuinamente otro, y recibir de ese otro lo que uno no podía darse a sí mismo. Sin alteridad real, el amor se convierte en un espejo ---bonito, quizá, pero estéril---. Solo ante alguien verdaderamente distinto puedo salir de mí, arriesgarme, y descubrir que dar la vida tiene sentido.

\section{La alteridad como condición del encuentro}

Aquí es donde la reflexión se pone interesante. Porque si nos preguntamos por qué la naturaleza ---o, si se prefiere, la evolución--- ha mantenido la diferencia sexual en una especie tan sofisticada como la nuestra, la respuesta no es solo ``para la reproducción''. La reproducción sexual existe en organismos que no tienen nada parecido a lo que nosotros llamamos encuentro. Pero en el ser humano, la diferencia sexual hace posible algo más: hace posible que haya un \textit{otro} de verdad.

Pensemos en esto despacio. ¿Qué pasa cuando me encuentro con alguien que es exactamente igual a mí? Que, en el fondo, no me encuentro con nadie. Me encuentro con un espejo. Y un espejo no me interpela, no me desafía, no me saca de mí mismo. Me confirma. Me deja donde estoy.

El verdadero encuentro ---ese que me transforma, que me obliga a salir de mi mundo y entrar en el de otro--- solo es posible cuando hay alteridad, cuando el otro es realmente \textit{otro}. Alguien que ve las cosas de un modo que yo no había previsto. Que siente de una manera que no es la mía. Que me revela dimensiones de la realidad a las que yo solo, desde mi perspectiva, no tengo acceso.

La diferencia sexual es, quizá, la forma más radical y más cotidiana de alteridad que experimentamos los seres humanos. No la única, pero sí la más constitutiva. Cuando un varón y una mujer se encuentran de verdad ---no como estereotipos, sino como personas concretas, con toda su complejidad---, cada uno tiene la oportunidad de descubrir un modo de ser humano que no es el suyo. Y eso es un regalo, no una amenaza.\footnote{Emmanuel Levinas desarrolló filosóficamente esta idea de la alteridad como condición del encuentro ético. Para Levinas, el rostro del otro ---irreductiblemente diferente de mí--- es lo que funda la relación ética y me saca de mi soledad. Cf. Levinas, E., \textit{Totalidad e infinito}, Sígueme, 1977.}

Es importante decir algo aquí con honestidad. Al hablar de la riqueza de la diferencia sexual en el matrimonio, no estamos negando que exista amor verdadero, generoso y sacrificado en otras formas de relación. La experiencia nos muestra personas que aman con autenticidad y entrega a personas de su mismo sexo, y ese amor ---que es libre, que busca el bien del otro, que es capaz de sacrificio--- merece respeto y comprensión.\footnote{El \textit{Catecismo de la Iglesia Católica} pide expresamente que las personas homosexuales «deben ser acogidas con respeto, compasión y delicadeza» y que «se evitará, respecto a ellos, todo signo de discriminación injusta'' (CIC, n. 2358). El Papa Francisco ha insistido en múltiples ocasiones en que la Iglesia no puede ser una aduana sino una casa abierta.} Lo que este capítulo explora es algo específico: la riqueza particular que la diferencia sexual aporta al encuentro conyugal. Del mismo modo, tampoco excluimos otras vocaciones orientadas al amor: la vocación sacerdotal, la vida consagrada, quienes eligen el celibato por el Reino o por otras razones ---todas ellas son caminos auténticos de entrega y de amor---. Este libro habla de una vocación particular, la del matrimonio entre varón y mujer, que tiene una característica singular: es la más común dentro del género humano. La inmensa mayoría de las personas que han vivido, viven y vivirán están llamadas a amar de esta manera concreta. Hablar de ella con profundidad no es menospreciar las demás. Es, sencillamente, mirar con asombro lo que la diferencia entre varón y mujer hace posible, sin que eso signifique que el amor solo exista ahí.

\section{La riqueza que nace de la diferencia}

Bajemos de la filosofía a la vida concreta. Cualquier pareja con algunos años de convivencia sabe de qué estamos hablando, aunque quizá nunca lo haya formulado así.

Ella llega a casa preocupada por un problema en el trabajo. No quiere que él se lo resuelva ---quiere contárselo, necesita que la escuche, que valide lo que siente---. Él, mientras tanto, escucha el problema y automáticamente empieza a buscar soluciones. No porque no le importe lo que ella siente, sino porque así funciona su cabeza: si hay un problema, hay que resolverlo. Los dos tienen razón. Los dos están expresando formas legítimas de afrontar la dificultad. Y los dos necesitan aprender del otro: ella, a veces, necesita que le ayuden a actuar; él, a veces, necesita que le ayuden a parar y simplemente estar.

¿Es este ejemplo un estereotipo? Parcialmente, sí. No todas las mujeres procesan así ni todos los hombres son así. Pero la investigación en psicología de la comunicación muestra que estas tendencias existen como patrones estadísticos, y que buena parte de los conflictos de pareja nacen precisamente de no entender que el otro no está fallando ---está siendo diferente---.\footnote{Cf. Tannen, D., \textit{You Just Don't Understand: Women and Men in Conversation}, William Morrow, 1990. Tannen, lingüista de la Universidad de Georgetown, documentó extensamente las diferencias en los estilos comunicativos de hombres y mujeres, subrayando que la mayoría de los malentendidos nacen de asumir que el otro debería comunicarse como uno mismo.}

La diferencia obliga a salir de uno mismo. Obliga a escuchar, a intentar comprender una lógica que no es la propia, a aceptar que mi forma de ver el mundo no es la única forma de verlo. Y eso, que a veces resulta agotador, es exactamente lo que nos hace crecer.

Una pareja donde los dos son idénticos ---donde piensan igual, sienten igual, reaccionan igual--- sería una pareja cómoda, quizá, pero estéril. No en el sentido biológico, sino en el sentido humano: no habría fricción creativa, no habría ese choque fecundo entre dos maneras de ver que obliga a construir algo nuevo, algo que ninguno de los dos habría podido construir solo.

\section{El cuerpo que habla}

Hay algo más. La diferencia sexual no es solo una diferencia psicológica o de carácter. Es, ante todo, una diferencia inscrita en el cuerpo. Y el cuerpo humano, como hemos visto en los capítulos anteriores, no es un accidente: es un lenguaje.

Nuestro cuerpo habla. Habla cuando una madre abraza a su hijo y la oxitocina inunda a los dos. Habla cuando el llanto de un bebé activa circuitos de alarma en el cerebro de sus padres. Habla cuando el contacto piel con piel entre dos personas que se quieren baja el cortisol y sube la sensación de seguridad.

Y habla también a través de la diferencia sexual. El cuerpo del varón y el de la mujer ``dicen'' cosas diferentes. No mejores ni peores: diferentes. La configuración hormonal de ella la inclina ---en términos estadísticos, insisto--- hacia una mayor sensibilidad a las señales emocionales del otro, hacia una mayor facilidad para el lenguaje verbal de las emociones, hacia una respuesta al estrés que busca la conexión.\footnote{Cf. la ya mencionada hipótesis \textit{tend and befriend} de Taylor et al. Estudios posteriores han confirmado que el sistema oxitocina-estrógenos favorece respuestas prosociales ante el estrés, especialmente en mujeres.} La configuración hormonal de él lo inclina hacia una mayor orientación espacial, hacia una respuesta al estrés más orientada a la acción, hacia una protectividad que se expresa muchas veces más con hechos que con palabras.\footnote{Cf. Archer, J., ``The reality and evolutionary significance of human psychological sex differences'', \textit{Biological Reviews}, 2019. Archer ofrece una revisión equilibrada de las diferencias psicológicas entre sexos, distinguiendo las que tienen base evolutiva sólida de las que son más culturales.}

Estas inclinaciones no son cárceles. Son puntos de partida. Cada persona concreta las modula, las amplía, las matiza con su historia, su educación, su libertad. Pero negar que existen es negar lo que el cuerpo dice. Y en un libro que ha empezado escuchando lo que el cuerpo tiene que decirnos sobre el amor, sería incoherente dejar de escucharlo aquí.

\section{Ni uniformidad ni guerra}

Si tuviéramos que resumir lo que hemos visto en este capítulo, podríamos decirlo así: la diferencia sexual humana se mueve entre dos peligros. Uno es antiguo: usar la diferencia para justificar la desigualdad, asignando roles rígidos y jerarquías donde no las hay. El otro es más reciente: negar la diferencia en nombre de la igualdad, como si reconocer que varones y mujeres no son idénticos fuera automáticamente una forma de discriminación.\footnote{Cf. Badinter, E., \textit{XY: De l'identité masculine}, Odile Jacob, 1992; y en una línea diferente, Halpern, D. F., \textit{Sex Differences in Cognitive Abilities}, Psychology Press, 2012. Ambas autoras, desde perspectivas distintas, ilustran la tensión entre reconocer las diferencias y evitar que se conviertan en herramientas de opresión.}

La realidad es más sutil que ambos extremos. Somos iguales en dignidad y diferentes en el modo de ser. Y esa diferencia no es un error a corregir ni una jerarquía a imponer: es una invitación al encuentro.

Pensemos en una orquesta. Que el violín y el violonchelo sean instrumentos diferentes no significa que uno sea superior al otro. Pero tampoco significa que den la misma nota. La belleza de la música nace precisamente de que cada uno aporta un timbre, un registro, una textura que el otro no tiene. Si todos los instrumentos fueran iguales, no habría armonía: habría unísono. Y el unísono, por potente que sea, es infinitamente más pobre que la sinfonía.

El varón y la mujer son como dos instrumentos distintos que, cuando afinan juntos, producen algo que ninguno de los dos podría producir solo. Eso es la complementariedad bien entendida. No la media naranja que me completa porque estoy incompleto, sino la otra voz que enriquece la canción precisamente porque suena diferente.

\section{Ante el misterio del otro}

Hay una última cosa que conviene decir, y quizá sea la más importante.

La diferencia sexual nos sitúa, permanentemente, ante el misterio del otro. Porque por mucho que conviva con mi pareja, por mucho que la conozca, por mucho que hayamos compartido, siempre habrá algo en ella ---o en él--- que se me escapa. Una forma de sentir que no es la mía. Una manera de estar en el mundo que no puedo experimentar desde dentro. Un misterio que no se agota.

Y eso, lejos de ser una frustración, es la garantía de que el amor no se aburre. Porque lo que se posee del todo deja de interesar. Lo que se domina completamente deja de fascinar. Solo lo que conserva algo de misterio sigue invitándome a acercarme, a descubrir, a asombrarme.

La diferencia sexual nos recuerda que el otro no es una extensión de mí mismo. No es mi reflejo, ni mi complemento mecánico, ni la pieza que me faltaba. Es alguien. Alguien irreductiblemente distinto, irreductiblemente libre, irreductiblemente otro. Y es justamente por eso que puedo amarlo de verdad. Porque el amor, el verdadero, no es fusión: es encuentro. Y para que haya encuentro, tiene que haber dos.

Dos, y no uno solo.

\paracaminar

\begin{enumerate}
  \item Pensad juntos: ¿en qué momentos concretos de vuestra relación habéis experimentado la diferencia del otro como una riqueza? ¿Y en qué momentos os ha costado más aceptarla?

  \item ¿Hay algo de vuestra pareja que, por más que lo intentéis, no termináis de entender del todo? ¿Podéis ver ese ``misterio'' no como un problema, sino como una invitación a seguir conociéndoos?

  \item Hablad con sinceridad: ¿alguna vez habéis caído en la tentación de querer que el otro piense, sienta o reaccione como vosotros? ¿Qué pasaría si, en lugar de eso, agradeciérais la diferencia?
\end{enumerate}

\vinetacapitulo{ch05}

\chapter{Diseñados para amar}

\epigraph{Dios es amor, y quien permanece en el amor permanece en Dios.}{1 Juan 4,16}

Hagamos una pausa.

Llevamos cinco capítulos caminando juntos, y quizá convenga detenerse un momento para mirar el paisaje que hemos recorrido. No para repetirlo, sino para verlo de golpe, como quien sube a una colina y contempla el camino que serpentea a sus pies.

\section{Lo que el cuerpo nos ha dicho}

Empezamos por el parto. Vimos que nacemos de la forma más difícil del reino animal, radicalmente inmaduros, con la cabeza sin cerrar y sin poder sostenernos solos. Que esa fragilidad no es un fallo sino el precio de nuestro cerebro extraordinario. Y que, precisamente por nacer así, necesitamos que alguien nos sostenga, nos alimente, nos quiera ---durante años--- para llegar a ser lo que somos.

Después descubrimos que el cuerpo no nos deja solos ante esa tarea. La oxitocina inunda a la madre en el parto y la lactancia, empujándola al vínculo. El padre cambia también: baja la testosterona, sube la oxitocina, se activan circuitos cerebrales que antes dormían. El cuerpo de ambos se transforma para cuidar.

Vimos que nos enamoramos con una química que se parece a la obsesión ---dopamina, noradrenalina, serotonina baja--- y que después, si el vínculo madura, la oxitocina y la vasopresina construyen algo más estable, más hondo: un apego que nos empuja a quedarnos. Que el amor no es solo un sentimiento bonito, sino que tiene un sustrato biológico que nos inclina a la permanencia.

Desmontamos el mito del amor romántico de las películas y descubrimos que el amor real no se encuentra: se construye. Que es decisión además de sentimiento. Que tiene capas ---deseo, amistad, entrega--- y que crece cuando se elige cada día, no cuando se espera que todo fluya solo.

Y vimos, por último, que la diferencia entre varón y mujer no es un accidente ni un problema, sino la condición misma del encuentro. Que solo puedo encontrarme de verdad con quien es distinto de mí. Que la alteridad no amenaza al amor: lo hace posible.

Todo esto lo hemos visto sin necesidad de abrir la Biblia, sin citar un solo dogma, sin invocar ninguna autoridad religiosa. Lo hemos leído en los huesos, en las hormonas, en las neuronas. En el cuerpo.

Y el cuerpo dice, con una claridad que asombra: estás hecho para la relación. Estás configurado para el vínculo. Estás inclinado al amor.

\section{La pregunta que faltaba}

Pero hay una pregunta que el cuerpo, por sí solo, no puede responder.

¿Por qué?

¿Por qué estamos hechos así? ¿Por qué la evolución ---o lo que sea que haya detrás--- nos ha configurado como seres que necesitan amar y ser amados para sobrevivir, para crecer, para ser plenamente humanos?

La biología ofrece una respuesta, y es honesta: porque así sobrevivimos mejor. Porque las crías cuidadas por padres vinculados tenían más probabilidades de llegar a la edad adulta. Porque la selección natural favoreció a los grupos que cooperaban, que cuidaban, que permanecían juntos. Es una respuesta verdadera. Pero ¿es suficiente?

Pensemos un momento. Si todo se reduce a supervivencia, entonces el amor es un truco ---un mecanismo elegante, sí, pero un truco al fin y al cabo---. La oxitocina que inunda a una madre cuando abraza a su hijo sería solo un recurso para que la especie no se extinga. El deseo de permanecer junto a alguien, de construir una vida compartida, de ser fiel, sería una ilusión útil. Una estrategia de la evolución disfrazada de sentimiento.

¿De verdad es eso todo?

Hay algo en nosotros que se rebela contra esa explicación. Algo que intuye que el amor que sentimos ---el de verdad, el que nos desborda, el que nos hace mejores, el que a veces duele--- es algo más que bioquímica al servicio de la reproducción. Que cuando una madre mira a su hijo recién nacido y siente que daría la vida por él, ahí hay algo que trasciende la selección natural. Que cuando dos personas se prometen fidelidad para siempre, están haciendo algo que ningún gen les ha pedido.

La ciencia puede describir el cómo. Puede explicar los mecanismos, las hormonas, los circuitos neuronales. Pero el porqué último ---la razón de que existamos así, configurados para el amor--- queda fuera de su alcance. No porque la ciencia falle, sino porque esa pregunta pertenece a otro orden.

Y es aquí donde este libro da un paso.

\section{El que nos diseña es Amor}

Hay una frase en la primera carta de san Juan que, cuando se lee despacio, lo cambia todo. Una frase breve, de apenas tres palabras en griego ---\textit{ho theós agápe estín}---, que contiene quizá la afirmación más radical de toda la historia del pensamiento:

\begin{quotation}
\textit{Dios es amor.}\footnote{1~Jn 4,8. El texto completo dice: ``El que no ama no ha conocido a Dios, porque Dios es amor.'' No dice que Dios \textit{tiene} amor, o que \textit{da} amor, o que \textit{manda} amar. Dice que Dios \textit{es} amor. Es una afirmación ontológica: define la naturaleza misma de Dios.}
\end{quotation}

No dice que Dios es poderoso (que lo es). No dice que Dios es sabio (que lo es). Dice que Dios es amor. No como una cualidad entre otras, sino como definición de lo que Él es. Amor no es algo que Dios hace; es lo que Dios es.

Y ahora viene la segunda pieza. El libro del Génesis dice que Dios creó al ser humano ``a su imagen y semejanza''.\footnote{Gn 1,27. La expresión hebrea \textit{tselem} (imagen) y \textit{demut} (semejanza) indica que el ser humano refleja algo de Dios mismo, no en un sentido físico, sino en su capacidad de conocer, de querer y de entregarse. Cf. Concilio Vaticano~II, Constitución pastoral \textit{Gaudium et Spes}, n.~12: ``Creyentes y no creyentes están generalmente de acuerdo en este punto: todos los bienes de la tierra deben ordenarse en función del hombre, centro y cima de todos ellos.''} Si Dios es amor, entonces su imagen ---el ser humano--- tiene que ser un ser capaz de amar. No puede ser de otro modo. Un espejo refleja lo que tiene delante; y si lo que tiene delante es amor, el reflejo será un ser hecho para amar.

Ahora volvamos la vista a todo lo que hemos recorrido en estos capítulos. La pelvis que se estrecha. El cerebro que crece. El bebé que nace indefenso. La oxitocina que inunda a la madre. La testosterona del padre que baja. La química que nos enamora. El apego que nos hace quedarnos. La diferencia que nos abre al otro.

¿Y si todo eso no fuera solo evolución?

¿Y si fuera la firma del Diseñador?

Quiero ser honesto: lo que acabo de hacer no es una demostración científica, sino una invitación. Los hechos biológicos que hemos recorrido juntos son sólidos y se sostienen por sí mismos, creas o no en un Diseñador. Pero hay un momento en que uno mira todo eso y siente que puede significar algo más ---y dar ese paso es un acto de confianza personal, no una conclusión de laboratorio---. Aquí es donde entra mi fe, y tú eres libre de acompañarme o de quedarte con la ciencia a solas, que ya dice mucho.

No digo que la evolución no sea real ---lo es, y lo hemos visto con rigor---. Digo que la evolución, siendo verdadera, puede no ser la explicación última. Que detrás del cómo biológico puede haber un porqué más profundo. Que las mismas fuerzas que la ciencia describe con precisión pueden ser, al mismo tiempo, el modo en que un Dios que es amor ha querido configurar a la criatura hecha a su imagen.

Estamos diseñados para amar. No solo porque la selección natural lo favoreció, sino porque Alguien que es Amor nos pensó así desde el principio.\footnote{El Catecismo de la Iglesia Católica lo expresa con claridad: ``Siendo a imagen de Dios, el ser humano tiene la dignidad de persona; no es solamente algo, sino alguien. Es capaz de conocerse, de poseerse y de darse libremente y entrar en comunión con otras personas'' (CIC, n.~357). Y el Concilio Vaticano~II: ``El hombre, única criatura terrestre a la que Dios ha amado por sí misma, no puede encontrar su propia plenitud si no es en la entrega sincera de sí mismo a los demás'' (\textit{Gaudium et Spes}, n.~24).}

\section{La libertad, condición del amor}

Pero Dios no se conformó con hacernos capaces de amar. Hizo algo más arriesgado, más hermoso y más terrible a la vez: nos hizo libres.

¿Por qué? Porque un amor obligado no es amor. Es obediencia, es rutina, es a lo sumo costumbre. Pero no es amor. El amor, para ser amor de verdad, necesita poder decir que no.

Pensemos en algo sencillo. Si alguien te dice ``te quiero'' porque no tiene más remedio, porque está programado para decirlo, porque no puede evitarlo... ¿te sentirías querido? Probablemente no. Sentirías que eres el destinatario de un automatismo, no de un afecto. Lo que hace que un ``te quiero'' valga es que la persona que lo dice podría no decirlo. Podría irse, podría callarse, podría elegir a otro. Pero te elige a ti. Libremente.

Dios quiere ser amado así. No por obligación, no por inercia, no por programación genética. Libremente. Y para eso nos da algo que ningún otro ser en la creación tiene de la misma manera: conciencia y libertad.\footnote{Cf. \textit{Gaudium et Spes}, n.~17: ``La dignidad humana requiere, por tanto, que el hombre actúe según su conciencia y libre elección, es decir, movido e inducido por convicción interna personal y no bajo la presión de un ciego impulso interior o de la mera coacción externa.''} La capacidad de conocer lo que hacemos y de elegir hacerlo ---o no hacerlo---.

¿Es esto un riesgo? Claro que lo es. La libertad implica que podemos elegir mal. Que podemos dañar, mentir, abandonar, traicionar. Que podemos mirar al amor de frente y darle la espalda. La historia humana está llena de libertad mal usada, y las heridas que deja son reales.

Pero la alternativa sería peor. La alternativa sería un mundo de autómatas que ``aman'' porque no pueden hacer otra cosa. Un mundo sin traiciones, sí, pero también sin fidelidad ---porque la fidelidad solo tiene sentido cuando podrías ser infiel---. Un mundo sin egoísmo, pero también sin generosidad ---porque la generosidad solo existe cuando podrías quedarte con todo---. Un mundo sin dolor, quizá, pero también sin alegría verdadera.

Dios no crea la libertad a regañadientes, como un defecto inevitable. La crea como condición necesaria. Sin libertad, no habría imagen de Dios. Sin libertad, no habría amor. Sin libertad, seríamos otra cosa ---quizá más eficiente, quizá más segura--- pero no seríamos humanos.\footnote{Santo Tomás de Aquino desarrolla esta idea con precisión: el ser humano es imagen de Dios precisamente en cuanto es principio de sus propias acciones, al tener libre albedrío y dominio sobre sus actos. Cf. \textit{Summa Theologiae}, I-II, prólogo.}

\section{La primera pregunta}

Y aquí es donde todo lo que hemos dicho aterriza en algo concreto. Tan concreto como una iglesia, dos personas de pie ante un altar, y un sacerdote que abre la boca para hacer la primera pregunta del rito del matrimonio.

No pregunta: ``¿Os queréis?'' No pregunta: ``¿Estáis enamorados?'' No pregunta: ``¿Prometéis ser fieles?'' Todo eso vendrá después. Lo primero, lo que va antes de cualquier otra cosa, es esto:

\begin{quotation}
\textit{¿Venís a contraer matrimonio sin ser coaccionados, libre y voluntariamente?}\footnote{Ritual del Matrimonio, n.~60 (edición típica en español). Las tres preguntas del escrutinio son, en este orden: libertad, fidelidad y apertura a la vida. El orden no es casual.}
\end{quotation}

¿Venís libremente?

La Iglesia lleva dos mil años celebrando matrimonios, y en todo ese tiempo no ha cambiado el orden de las preguntas. Primero la libertad. Antes que el amor, antes que la fidelidad, antes que los hijos. Porque sin libertad, nada de lo que sigue es real.

Si no vienes libremente, tu ``sí'' no es un sí. Es un sonido vacío. Si no vienes libremente, tu promesa de fidelidad es papel mojado. Si no vienes libremente, tu entrega es una representación teatral.

La libertad es el suelo sobre el que se construye todo lo demás. Y no es casualidad que sea así. Es coherencia perfecta con lo que somos: seres hechos a imagen de un Dios que es amor, y que por tanto solo pueden amar de verdad cuando eligen libremente hacerlo.

Pero conviene entender bien qué es la libertad, porque tendemos a confundirla con hacer lo que apetece en cada momento. Y no es eso. Pensad en alguien que decide correr una maratón. Desde el momento en que toma esa decisión, su vida cambia: se levanta temprano a entrenar, se somete a dietas, renuncia a salir, aguanta sesiones que no le apetecen nada. ¿Ha dejado de ser libre? Al contrario: es \emph{más} libre que antes, porque está usando su libertad para algo que ha elegido y que le trasciende. O pensad en el opositor que se encierra a estudiar durante dos, tres, cuatro años. Renuncia a vacaciones, a vida social, a dormir hasta tarde. ¿Es un esclavo? No. Es alguien que ha decidido libremente labrarse un futuro, y que acepta el precio porque la decisión es suya y el horizonte merece la pena.

Pues el matrimonio funciona exactamente igual. Elegir a una persona para toda la vida no te quita libertad: la pone al servicio de algo grande. Las renuncias que implica ---y las habrá, todos los días--- no son cadenas que te impone nadie. Son el entrenamiento del maratoniano, la disciplina del opositor: el precio libre de una decisión libre.

Decía san Josemaría Escrivá que ``porque quiero'' es la razón más sobrenatural que existe.\footnote{San Josemaría Escrivá de Balaguer (1902-1975), fundador del Opus Dei. Esta idea recorre su predicación sobre la libertad como fundamento de la vida cristiana. Cf. \textit{Amigos de Dios}, nn.~24-38.} Y tenía razón. Porque en esas dos palabras se manifiesta lo que nos hace imagen de Dios: la libertad. No ``porque debo'' ---que es obediencia, que es norma, que es el mínimo---. Sino ``porque quiero'' ---que es elección, que es don, que es lo más parecido a cómo actúa Dios---.

Y cuando dos personas se miran a los ojos en el altar y dicen ``sí, quiero'', están pronunciando la frase más sobrenatural que puede decir un ser humano. No ``sí, debo''. No ``sí, me toca''. ``Sí, quiero.'' Libremente. Como Dios nos ama: porque quiere, no porque necesite hacerlo.

\section{El puente}

Hemos llegado, entonces, a un lugar nuevo.

Durante la primera parte de este libro, hemos mirado al ser humano desde la ciencia: la biología, la neurociencia, la psicología, la antropología. Y hemos descubierto algo asombroso: que todo en nosotros ---huesos, hormonas, neuronas, emociones--- apunta en la misma dirección. Hacia el vínculo. Hacia el cuidado. Hacia el amor.

Ahora sabemos por qué. Estamos diseñados para amar porque el que nos diseña es Amor. Y nos hace libres para que ese amor sea verdadero.

Con esto en las manos, la segunda parte de este libro se abre de forma natural. Si estamos diseñados para amar ---si esa es nuestra vocación más profunda---, entonces la pregunta siguiente es inevitable: ¿cómo? ¿Cómo se concreta ese amor? ¿Qué forma toma? ¿Qué nos ayuda a vivirlo bien, y qué lo desfigura?

La fe cristiana tiene una propuesta. No la impone; la ofrece. Y lo que ofrece no es un conjunto de normas sobre lo que se puede y no se puede hacer, sino una mirada ---una forma de entender el amor, el matrimonio y la familia que, curiosamente, encaja con lo que el cuerpo ya nos había dicho desde el principio---.

Eso es lo que vamos a explorar ahora.

Un aviso antes de seguir. Si la primera parte de este libro ha podido sorprender ---porque quizá no esperabas encontrar oxitocina y topillos en un libro sobre matrimonio---, la segunda parte transitará un territorio más conocido: el de la fe cristiana, sus sacramentos, su visión del amor. Sobre esto hay mucho y muy bueno ya escrito ---de Juan Pablo II a Francisco, de C.\,S. Lewis a Benedicto XVI---. No pretendo inventar nada: pretendo reunir, con la misma honestidad de los capítulos anteriores, lo que otros han dicho mejor que yo, y ponerlo en diálogo con lo que el cuerpo ya nos había enseñado. Si la primera parte descubría, la segunda reconoce. Y reconocer algo que ya estaba ahí puede ser, a su manera, tan revelador como descubrirlo por primera vez.

\paracaminar

\begin{enumerate}
  \item Cuando leéis que ``Dios es amor'' y que estamos hechos a su imagen, ¿qué resuena en vosotros? ¿Os parece una idea lejana o conecta con algo que habéis experimentado?

  \item Hablad juntos sobre la libertad en vuestra relación. ¿Sentís que os elegís libremente cada día? ¿Hay algo ---presión social, inercia, miedo--- que a veces nubla esa libertad?

  \item La primera pregunta del rito matrimonial es ``¿Venís libremente?'' ¿Por qué creéis que va antes que las demás? ¿Qué os dice eso sobre lo que la Iglesia entiende por amor?
\end{enumerate}

\vinetacapitulo{ch06}


% ===========================
% PARTE II: LLAMADOS A AMAR
% ===========================
\part{Llamados a amar}

\chapter{No es bueno estar solo}

\epigraph{«Varón y mujer los creó.»}{Génesis 1,27}

Hay un texto que casi todo el mundo cree conocer, y que casi nadie ha leído con cuidado.

Me refiero, claro, al relato de la creación. Todos tenemos en la cabeza una especie de resumen que va más o menos así: Dios crea al hombre, lo pone en el jardín, ve que está solo, le quita una costilla y fabrica a la mujer. Fin. La mujer aparece después, como un añadido, casi como un parche. Como si Dios hubiera hecho al varón, se hubiera dado cuenta de que le faltaba algo, y hubiera improvisado una solución.

Pero eso no es lo que dice el Génesis. O, mejor dicho, no es \emph{todo} lo que dice. Porque el Génesis tiene dos relatos de la creación, no uno. Y cuando los leemos juntos, con calma, lo que encontramos es bastante más hondo ---y bastante más hermoso--- de lo que el resumen sugiere.

\section{Lo que el Génesis dice de verdad}

El primer relato está en el capítulo 1. Es el más solemne, el más cósmico. Dios va creando por etapas ---la luz, el firmamento, las aguas, los animales--- y cada vez dice: «Es bueno.» Hasta que llega al ser humano. Y entonces el texto hace algo distinto. No dice «hágase el hombre» como dijo «hágase la luz». Dice algo mucho más íntimo:

\begin{quotation}
\textit{«Hagamos al hombre a nuestra imagen y semejanza. [\ldots] Y creó Dios al hombre a su imagen, a imagen de Dios lo creó; varón y mujer los creó.»}\footnote{Gn 1,26-27. El término hebreo \textit{adam} no designa aquí al varón, sino al ser humano, a la humanidad. Solo en el segundo relato (Gn 2) se introduce la diferenciación \textit{ish} (varón) e \textit{ishshah} (mujer).}
\end{quotation}

Quiero que nos detengamos aquí un momento, porque esta frase contiene algo que cambia toda la lectura. Fijémonos: Dios no crea primero al varón y después a la mujer. Crea al ser humano ---\textit{adam}, la humanidad--- y lo crea «varón y mujer». Ambos a la vez. Ambos imagen de Dios. No hay un primero y un segundo. No hay un original y una copia. Hay un ser humano que, desde el principio, es relación: varón \emph{y} mujer.

Solo después de crear así al ser humano, Dios dice algo que no había dicho de nada en toda la creación: «No es bueno que el hombre esté solo.»\footnote{Gn 2,18. Es la primera vez en todo el relato de la creación que algo «no es bueno». Juan Pablo II desarrolla extensamente el significado de esta «soledad originaria» en sus catequesis sobre la Teología del Cuerpo. Cf. Juan Pablo II, \textit{Hombre y mujer los creó. Catequesis sobre el amor humano}, audiencias generales de septiembre a noviembre de 1979.}

Y aquí está la clave: si la mujer fuera simplemente un complemento que viene a llenar un hueco, a solucionar un problema de diseño, Dios habría creado primero al varón solo y luego habría tenido que corregir su propia obra. Pero no es eso lo que pasa. Lo que pasa es que el ser humano es, desde su misma creación, un ser relacional. La capacidad de amar, la necesidad del otro, no es algo que se le añade después: está inscrita en su origen. Es constitutiva. La mujer no viene a arreglar un defecto del varón. Ni el varón a completar una carencia de la mujer. El «no es bueno que el hombre esté solo» no es la constatación de un fallo, sino la revelación de una verdad: estamos hechos para el encuentro.

\section{La biología ya lo sabía}

Si habéis leído la primera parte de este libro, todo esto os sonará. Y es que resulta que la fe y la ciencia, cada una desde su lenguaje, dicen lo mismo.

Hemos visto que nacemos radicalmente inmaduros: un «feto externo» que necesita meses de cuidado constante para sobrevivir.\footnote{Cf. capítulo 1 de este libro.} Hemos visto que nuestro cuerpo está equipado con mecanismos hormonales ---oxitocina, vasopresina, dopamina--- que nos empujan a vincularnos, a cuidar, a quedarnos.\footnote{Cf. capítulos 2 y 3.} Hemos visto que la diferencia entre varón y mujer no es un accidente cultural, sino un dato inscrito en nuestra biología que hace posible la complementariedad y el encuentro.\footnote{Cf. capítulo 5.}

La biología nos dice: no estáis hechos para vivir solos. Necesitáis al otro. El amor no es un lujo, es una condición de supervivencia.

Y ahora el Génesis llega y le pone nombre a lo que el cuerpo ya sabía: «No es bueno que el hombre esté solo.» No es bueno, no porque le falte compañía, sino porque sin el otro no puede amar. Y sin amar, no puede ser plenamente lo que es: imagen de Dios.

\section{Imagen de un Dios que es comunión}

Aquí entramos en algo que la tradición cristiana ha ido comprendiendo poco a poco a lo largo de siglos, y que quizá sea una de las intuiciones más profundas de toda la teología.

Cuando el Génesis dice que el ser humano es «imagen de Dios», ¿qué quiere decir? ¿Que nos parecemos a Dios físicamente? Obviamente no. ¿Que somos inteligentes, libres, creativos? Sí, también. Pero hay algo más radical.

Dios, para la fe cristiana, no es un ser solitario. Dios es Trinidad: Padre, Hijo y Espíritu Santo. No tres dioses, sino un solo Dios que es, en sí mismo, relación, comunión, entrega. El Padre se entrega totalmente al Hijo. El Hijo se entrega totalmente al Padre. Y esa entrega mutua es tan real, tan fecunda, que «produce» una tercera persona: el Espíritu Santo, que es el Amor en persona.\footnote{«Dios, que es amor y vive en sí mismo un misterio de comunión personal de amor, crea al género humano a su imagen [\ldots] inscribiendo en la humanidad del hombre y de la mujer la vocación y, consiguientemente, la capacidad y la responsabilidad del amor y de la comunión.» Juan Pablo II, Exhortación apostólica \textit{Familiaris Consortio}, n. 11 (1981).}

Dios no \emph{tiene} amor. Dios \emph{es} amor.\footnote{1~Jn 4,8.} Y ser imagen de ese Dios significa estar hecho para amar como Él ama: entregándose, saliendo de uno mismo, haciéndose don para el otro.

Por eso el ser humano no puede realizarse en solitario. No es una cuestión de debilidad o de dependencia psicológica. Es una cuestión de identidad. Somos imagen de un Dios-comunión: solo amando nos parecemos a Él. Solo en la entrega al otro encontramos quiénes somos.

\section{La costilla, la carne, el grito}

Ahora sí podemos entender el segundo relato del Génesis ---el de la costilla--- en toda su profundidad. Porque no es un cuento ingenuo sobre cirugía divina. Es una de las páginas más densas jamás escritas sobre lo que significa encontrar al otro.

Dios hace desfilar ante el hombre a todos los animales para que les ponga nombre.\footnote{Gn 2,19-20.} En la mentalidad bíblica, poner nombre es conocer, es comprender la esencia de algo. El hombre nombra a cada animal, los conoce, los entiende. Pero en ninguno se reconoce. Ninguno es «como él». Hay una soledad que ningún animal puede llenar ---por mucho que queramos a nuestro perro, que conste---.\footnote{Juan Pablo II llama a esta experiencia la «soledad originaria»: el ser humano se descubre distinto de todo lo demás, se reconoce como persona, como ser que busca un «tú». Cf. Juan Pablo II, audiencia general del 10 de octubre de 1979.}

Y entonces Dios actúa. No toma barro del suelo, como hizo para formar al primer ser humano. Toma algo del hombre mismo: una costilla, que en hebreo (\textit{tsela}) puede significar también «costado», «lado».\footnote{El término hebreo \textit{tsela'} aparece en otros libros del Antiguo Testamento con el significado de «lado» (por ejemplo, en la descripción de los laterales del Arca de la Alianza, Ex 25,12). La traducción como «costilla» es posible, pero empobrece el sentido: lo que Dios toma no es un hueso, sino un lado, una mitad.} La imagen es poderosa: la mujer no viene de fuera, no es un cuerpo extraño. Sale del mismo ser del hombre. Comparte su sustancia, su dignidad, su naturaleza.

Y cuando el hombre despierta y la ve, se produce el primer poema de la historia humana. Las primeras palabras que un ser humano pronuncia en toda la Biblia no son una oración, ni un mandamiento, ni una queja. Son un grito de asombro, de reconocimiento, de alegría:

\begin{quotation}
\textit{«¡Esta sí es hueso de mis huesos y carne de mi carne!»}\footnote{Gn 2,23.}
\end{quotation}

\textit{Esta sí.} En hebreo, \textit{zo't happa'am} ---«¡por fin!», «¡esta vez sí!»---. Después de haber visto pasar a todos los animales sin reconocerse en ninguno, el hombre ve a la mujer y por fin se encuentra. No encuentra a alguien igual que él ---si fuera igual, no habría encuentro---. Encuentra a alguien que es distinto y a la vez de su misma carne. Alguien a quien puede mirar cara a cara. Alguien a quien puede amar.

Hay un detalle que pasa desapercibido y que es precioso: el hombre solo se reconoce a sí mismo como varón (\textit{ish}) cuando reconoce a la mujer (\textit{ishshah}). Antes de ella, no tenía nombre propio. Era simplemente \textit{adam}, el ser humano. Solo ante el otro descubre quién es. Solo en la relación se revela la identidad.\footnote{«El hombre sólo se define como ``varón'' ante la ``mujer''; sólo la presencia del ``tú'' revela plenamente su ``yo''.» Cf. Juan Pablo II, audiencia general del 7 de noviembre de 1979.}

\section{El matrimonio como sacramento}

«Por eso dejará el hombre a su padre y a su madre, se unirá a su mujer, y serán una sola carne.»\footnote{Gn 2,24. Jesús mismo citará este versículo en Mt 19,5 para fundamentar la indisolubilidad del matrimonio, remitiendo a sus oyentes no a la Ley de Moisés, sino «al principio».}

Así termina el relato del Génesis. Y así comienza la teología del matrimonio.

Cuando dos personas se casan por la Iglesia, hay algo que mucha gente no sabe, o que sabe pero no ha terminado de entender: los esposos no «reciben» el sacramento del matrimonio como quien recibe un paquete. Los esposos \emph{son} los ministros del sacramento.\footnote{«Los esposos, como ministros de la gracia de Cristo, se confieren mutuamente el sacramento del Matrimonio, expresando ante la Iglesia su consentimiento.» \textit{Catecismo de la Iglesia Católica}, n. 1623. A diferencia de los demás sacramentos, el ministro del matrimonio no es el sacerdote o el diácono, sino los propios contrayentes. El sacerdote es testigo cualificado de la Iglesia, pero no administra el sacramento.} Son ellos quienes se lo dan mutuamente. El sacerdote no casa a nadie: es testigo. Testigo cualificado, sí, en nombre de la Iglesia. Pero quien realiza el sacramento son el hombre y la mujer que se dicen el uno al otro: «Yo te recibo como esposo. Yo te recibo como esposa.»

Esto no es un detalle litúrgico. Es una revolución.

Significa que el matrimonio no es un trámite administrativo con flores y arras. No es una bendición decorativa que la Iglesia añade a lo que la pareja ya tenía. Es un acto en el que dos personas, con su libertad y ante Dios, se convierten en signo vivo de algo mucho más grande que ellas mismas: el amor de Cristo por su Iglesia.\footnote{Ef 5,31-32: «"Por eso dejará el hombre a su padre y a su madre y se unirá a su mujer, y los dos serán una sola carne." Gran misterio es este, y yo lo refiero a Cristo y a la Iglesia.»}

Y ese signo no es abstracto. Ese signo es la vida entera de los esposos. Cada gesto de cuidado, cada perdón concedido, cada noche en vela junto a un hijo enfermo, cada renuncia silenciosa, cada abrazo después de una discusión: todo eso es sacramento. Todo eso hace visible, aquí y ahora, un amor que viene de más arriba.

\subsection*{Alianza, no contrato}

Hay una diferencia fundamental entre un contrato y una alianza, y entenderla cambia la forma de mirar el matrimonio.\footnote{Cf. Scott Hahn, \textit{La cena del Cordero}, Rialp, 2001. Hahn desarrolla con claridad la diferencia entre contrato y alianza en el contexto bíblico.}

Un contrato intercambia bienes y servicios: yo te doy esto, tú me das aquello. Si una parte incumple, la otra queda libre. Un contrato se puede rescindir. Un contrato protege intereses.

Una alianza intercambia personas: yo me doy a ti, tú te das a mí. No hay letra pequeña. No hay cláusula de rescisión. Una alianza crea un vínculo que transforma a quienes la contraen. «Ya no son dos, sino una sola carne»,\footnote{Mt 19,6.} dice Jesús. No dijo «un solo patrimonio» ni «un solo proyecto». Dijo «una sola carne»: una sola vida entregada.

El matrimonio cristiano es una alianza. Y cuando los novios se dicen «sí», están haciendo lo mismo ---salvando las infinitas distancias--- que hizo Dios con su pueblo: vincularse para siempre, sin condiciones, por amor. «Yo seré tu Dios y tú serás mi pueblo.»\footnote{Jr 31,33; cf. Os 2,21-22: «Te desposaré conmigo para siempre; te desposaré en justicia y en derecho, en amor y en compasión.»} Los profetas, de hecho, usaron constantemente la imagen del matrimonio para hablar de la relación entre Dios e Israel.\footnote{Toda la tradición profética emplea la metáfora nupcial: Oseas, Isaías, Jeremías, Ezequiel. El Cantar de los Cantares se ha leído, desde antiguo, como canto del amor entre Dios y su pueblo. No es casualidad que Jesús inaugure su vida pública en una boda (Jn 2,1-11).} No al revés. No es que el matrimonio sea «como» la alianza de Dios. Es que la alianza de Dios se expresa como un matrimonio, porque el matrimonio es la imagen más cercana que tenemos los humanos del amor total.

\section{Una cruz a la que abrazar}

Quiero terminar este capítulo con algo que no está en ningún manual de teología, pero que a mí me ha dado mucho que pensar. Es algo que suelo compartir en las charlas con novios, y que siempre genera un silencio especial en la sala.

Jesús dijo: «Nadie va al Padre sino por mí.»\footnote{Jn 14,6.} Y también dijo: «El que quiera seguirme, que tome su cruz y me siga.»\footnote{Cf. Mt 16,24; Mc 8,34; Lc 9,23.}

Normalmente, cuando pensamos en la Cruz, pensamos en sufrimiento. En sacrificio. En algo pesado, ominoso, doloroso. «Llevar la cruz» se ha convertido en sinónimo de aguantar, de soportar, de cargar con algo que preferirías no tener.

Pero paremos un momento. ¿Qué es la Cruz, en realidad?

La Cruz es el lugar donde Cristo entrega su vida. Sí, con dolor. Sí, con sufrimiento. Pero la Cruz no es principalmente sufrimiento. La Cruz es principalmente \emph{entrega}. Es el lugar donde Jesús dice, con todo su cuerpo y toda su sangre: «Te amo hasta el final. Doy mi vida por ti.»\footnote{Cf. Jn 13,1: «Habiendo amado a los suyos que estaban en el mundo, los amó hasta el extremo.»}

Y ahora pensemos en el matrimonio.

¿No es el matrimonio, precisamente, el lugar donde dos personas entregan su vida por amor? No de golpe, no con un gesto dramático, sino día a día, gota a gota. En el café que preparas por la mañana sabiendo cómo le gusta. En el cansancio compartido después de acostar a los niños. En la renuncia a tener razón para poder tener paz. En el «me quedo contigo» de cada día, que es menos espectacular que el del día de la boda pero infinitamente más real.

Si la Cruz es entrega de vida por amor, y el matrimonio es entrega de vida por amor\ldots{} entonces mi mujer es mi cruz. Mi marido es mi cruz. No en el sentido de «qué cruz de marido me ha tocado» ---que también tiene su gracia---. En el sentido más profundo: la persona a la que amo es la razón por la que entrego mi vida. Es mi camino de santidad. Es mi puerta al Padre.

«El que quiera seguirme, que tome su cruz.» Y Dios, que nos conoce bien, nos lo pone fácil. Nos da una cruz a la que podemos abrazar cada noche. Una cruz que nos devuelve el abrazo. Una cruz que nos dice «te quiero» y que nos necesita tanto como nosotros a ella. Y aquí está la clave: que devuelve el abrazo. Una cruz que aplasta, que humilla, que hace daño, no es la cruz de Cristo: es la cruz del verdugo. La cruz del matrimonio es siempre entrega mutua, nunca sufrimiento impuesto por uno sobre el otro.\footnote{Sobre las situaciones de violencia o abuso en el matrimonio, véase el capítulo 8 de este libro.}

Hay gente que busca la santidad en grandes penitencias, en ayunos heroicos, en gestas extraordinarias. Y está bien. Pero para los casados, el camino es más sencillo y más difícil a la vez: amar a la persona que tienes al lado. Amarla cuando es fácil y cuando no lo es. Amarla cuando te enamora y cuando te saca de quicio. Amarla como Cristo ama a su Iglesia: «entregándose por ella».\footnote{Ef 5,25.}

Esa es la cruz del matrimonio. Y es, quizá, la forma más hermosa de llevar la cruz que existe.

\paracaminar

\begin{enumerate}
  \item Leed juntos Génesis 1,26-27 y Génesis 2,18-24. ¿Qué os dice sobre vuestro matrimonio que el ser humano sea, desde el principio, «varón y mujer», creado para la relación?

  \item ¿Habéis pensado alguna vez en vuestro matrimonio como una \emph{alianza} y no como un contrato? ¿Qué diferencia práctica tiene eso en vuestra vida diaria?

  \item La imagen de la Cruz como entrega de amor: ¿cómo la recibís? ¿Podéis identificar momentos concretos en los que «entregar la vida» por el otro os haya hecho crecer como personas?
\end{enumerate}

\vinetacapitulo{ch07}

\chapter{La promesa que nadie te va a hacer}

\epigraph{«Lo que Dios ha unido, que no lo separe el hombre.»}{Mateo 19,6}

Antes de entrar en materia, quiero hacer una pausa y contarte algo que a mí me sorprendió cuando lo descubrí. Porque los tres capítulos que vienen ahora ---este y los dos siguientes--- giran en torno a tres palabras que suelen producir urticaria: indisolubilidad, fidelidad y fecundidad. Tres palabras que, cuando alguien las escucha juntas, tiende a pensar en reglas, prohibiciones y corsés. En cosas que la Iglesia exige y que la vida moderna ha dejado atrás.

Y sin embargo.

Si abres el Código Civil ---no el Derecho Canónico, no el catecismo, sino el Código Civil, el que regula el matrimonio para todos, creyentes o no---, te encuentras con algo llamativo. El artículo 68 dice, textualmente, que los cónyuges están obligados a vivir juntos, guardarse fidelidad y socorrerse mutuamente.\footnote{Art. 68 del Código Civil español: «Los cónyuges están obligados a vivir juntos, guardarse fidelidad y socorrerse mutuamente. Deberán, además, compartir las responsabilidades domésticas y el cuidado y atención de ascendientes y descendientes y otras personas dependientes a su cargo.»} Y el artículo 67 añade que deben respetarse y actuar en interés de la familia.\footnote{Art. 67 del Código Civil español: «Los cónyuges deben respetarse y ayudarse mutuamente y actuar en interés de la familia.»}

Fidelidad. Vida en común. Socorro mutuo. Interés de la familia. ¿Te suenan?

Son, con otras palabras, las mismas columnas sobre las que se levanta el matrimonio cristiano. No porque la Iglesia haya impuesto sus ideas al Estado ---que es lo que muchos piensan---, sino por algo mucho más interesante: porque estos principios son tan humanos, están tan inscritos en la naturaleza del amor, que el derecho civil los ha reconocido por su propio peso. La ley no necesitó consultar ningún catecismo para llegar a estas conclusiones. Bastó con observar qué es lo que hace que el amor funcione, qué es lo que necesita una familia para sostenerse, qué es lo que dos personas se deben la una a la otra cuando deciden compartir la vida.

Piénsalo un momento. Si algo tan «religioso» como la fidelidad aparece como obligación legal en un código civil, quizá no sea tan religioso como pensábamos. Quizá sea, simplemente, humano. Quizá la biología, la ley y la fe estén diciendo lo mismo con distintos acentos: que el amor de verdad tiene una forma, y que esa forma incluye permanencia, exclusividad y apertura a la vida.

Esto no debería sorprendernos. En la primera parte de este libro vimos que nuestro cuerpo ya viene equipado para amar así: con mecanismos hormonales que nos empujan a vincularnos, a quedarnos, a cuidar.\footnote{Cf. capítulos 2 y 3 de este libro.} Si el cuerpo habla de permanencia, si la ley habla de permanencia y si la fe habla de permanencia\ldots{} quizá sea porque la permanencia no es un invento de nadie, sino algo que está en la naturaleza misma del amor.

Dicho esto, vamos al grano. Porque este capítulo va sobre la primera de esas tres columnas. La más temida. La más malentendida. Y, cuando se mira bien, la más hermosa.

Hablemos de la indisolubilidad.

\section{El vínculo que no se cuestiona}

Pero antes, pensemos un momento en algo que damos por sentado y que, curiosamente, nadie discute.

De todos los vínculos familiares que tienes, absolutamente todos son indisolubles. Haga lo que haga tu hijo, seguirá siendo tu hijo. Siempre. Aunque os peleéis, aunque no os habléis durante años, aunque te decepcione de la peor manera posible: es tu hijo. Lo mismo con tu padre, con tu madre, con tus hermanos, con tus primos. No puede pasar nada ---nada--- que haga que tu primo deje de ser tu primo. No hay papeles que firmar, no hay juez que lo decrete, no hay ley que lo permita. Los vínculos de sangre son para siempre, y a nadie se le ocurre cuestionarlo.

Y ninguno de esos vínculos lo elegiste. Te vinieron dados. No decidiste nacer en esa familia, no escogiste a tus padres, no firmaste ningún contrato con tus hermanos. Simplemente, así son las cosas.

Y sin embargo ---aquí viene la paradoja---, pensamos que podemos deshacer el único vínculo familiar que decidimos crear libremente. El único que elegimos con plena conciencia, con libertad, delante de testigos. El único para el que, como hemos visto en la primera parte de este libro, estamos biológicamente diseñados. El que garantiza la supervivencia de nuestras crías y, como veremos, nuestra propia santidad.

Los vínculos que no elegimos son indisolubles sin discusión. El vínculo que elegimos libremente, ese sí creemos que podemos deshacerlo. ¿No es extraño?

Quizá la pregunta no sea si la indisolubilidad es razonable. Quizá la pregunta sea por qué nos parece razonable la disolubilidad.

\section{La cadena}

Hay pocas palabras que generen más rechazo inmediato que «indisolubilidad». Suena a cerrojo. A candado sin llave. A puerta que se cierra de golpe y ya no se abre jamás.

La reacción es comprensible. Vivimos en una cultura que celebra la libertad de elegir y, sobre todo, la libertad de dejar de elegir. Podemos cambiar de trabajo, de ciudad, de teléfono, de look, de pareja. La posibilidad de salir es lo que nos hace sentirnos libres. ¿Qué clase de libertad es la que te dice «esto es para siempre, sin marcha atrás»?

Y la objeción más frecuente, la que todo el mundo lleva en la punta de la lengua, es esta: «¿Y si no funciona?»

Es una pregunta razonable. Todos conocemos matrimonios rotos. Todos hemos visto el dolor de una separación, la amargura de quien se siente atrapado, el sufrimiento de los hijos. Ante eso, la indisolubilidad parece no solo anticuada, sino cruel. ¿Cómo puedes exigir a alguien que permanezca en algo que le hace daño?

Entiendo la objeción. De verdad. No la despacho con un gesto de la mano. El sufrimiento en un matrimonio roto es real, y merece todo el respeto y toda la compasión del mundo.\footnote{El Papa Francisco dedica todo el capítulo VIII de \textit{Amoris Laetitia} a acompañar la fragilidad, reconociendo las «situaciones imperfectas» y pidiendo a la Iglesia que sea «hospital de campaña» más que «aduana»: «Se trata de integrar a todos, se debe ayudar a cada uno a encontrar su propia manera de participar en la comunidad eclesial.» Francisco, Exhortación apostólica \textit{Amoris Laetitia}, n. 297 (2016).}

Pero quiero proponerte algo. Quiero pedirte que dejes en suspenso, solo por unos minutos, esa objeción. Que no la olvides ---la recogeremos después---, pero que antes mires la indisolubilidad desde otro ángulo. Uno que quizá no hayas considerado.

\section{¿Y si fuera al revés?}

¿Y si la indisolubilidad no fuera una cadena que te ata?

¿Y si fuera un regalo que te dan?

Piénsalo así. Imagina que alguien te dice: «Yo me quedo contigo. Pase lo que pase. En lo bueno y en lo malo. Cuando estés guapo y cuando no lo estés. Cuando seas encantador y cuando seas insoportable. Cuando todo vaya bien y cuando todo se tuerza. No me voy. No te voy a dejar. Nunca.»

¿Quién no querría escuchar eso?

No estoy hablando de una fantasía romántica. Estoy hablando de la necesidad más profunda del ser humano: la seguridad de ser amado sin condiciones. De saber que hay alguien en el mundo que ha decidido no irse. Que ha puesto su libertad ---lo más valioso que tiene--- al servicio de una promesa: «Contigo, hasta el final.»

Eso es la indisolubilidad. No es lo que la Iglesia te exige. Es lo que tu cónyuge te regala.

\section{La promesa que el mundo no sabe hacer}

Vivimos rodeados de promesas con fecha de caducidad. «Te quiero» dura mientras las cosas vayan bien. «Estoy contigo» lleva una letra pequeña que dice «salvo que encuentre algo mejor». Las relaciones se viven con una especie de periodo de prueba permanente: si funciona, genial; si no, cada uno por su lado.

Y parece muy razonable. Muy maduro. Muy libre.

Pero ¿lo has pensado bien? ¿Cómo vives cuando sabes que la otra persona puede irse en cualquier momento?

Vives auditando. Vigilando. Midiendo. ¿Le sigo gustando? ¿Sigo siendo suficiente? ¿Estará pensando en otro, en otra? ¿He engordado demasiado? ¿He dicho algo que no debía? ¿Estoy a la altura?

Vives, en el fondo, con miedo.

Y el miedo es lo contrario de la libertad.

Ahora imagina lo otro. Imagina que esa persona te ha dicho «no me voy», y lo ha dicho de verdad, y tú le crees. Que sabes, con la certeza que da una promesa pública y solemne, que mañana va a seguir ahí. Y pasado mañana. Y dentro de diez años. Y cuando tengas arrugas y el pelo blanco y ya no te acuerdes de dónde has puesto las llaves.

¿Qué pasa entonces?

Entonces puedes quitarte la máscara.

\section{La libertad de ser tú}

Solo en la seguridad de que el otro no se va puedes mostrarte como eres de verdad.

Piénsalo. En las primeras citas ---todos lo hemos hecho--- te presentas en tu mejor versión. Te peinas, te arreglas, mides tus palabras, controlas tus manías. Escondes lo que crees que no gusta. Muestras lo que crees que impresiona. Estás actuando. No mintiendo, exactamente, pero sí editando. Poniendo filtro.

Y tiene sentido. Cuando la otra persona aún no se ha comprometido, cuando todavía está decidiendo si se queda o se va, tú no puedes permitirte el lujo de ser vulnerable. Mostrar tus debilidades es un riesgo: ¿y si al verlas decide que no mereces la pena?

Ahora imagina que la decisión ya está tomada. Que hay una promesa. Que esa persona ha dicho, delante de su familia, de sus amigos y de Dios: «No me voy.»

Entonces puedes dejar de actuar. Puedes decir «estoy agotado y hoy no puedo con nada» sin miedo a que te dejen. Puedes confesar «me equivoqué, perdóname» sin miedo a que sea el fin. Puedes mostrar tus cicatrices, tus inseguridades, tus rincones más oscuros, y saber que la otra persona no va a salir corriendo. Porque ya decidió quedarse.

Eso no es una cadena. Eso es el hogar.

Hay una diferencia enorme entre una relación en la que estás siempre haciendo casting ---buscando que te renueven el contrato--- y una en la que sabes que tienes tu sitio. En la primera vives tenso, vigilante, midiendo cada paso. En la segunda puedes respirar. Puedes crecer. Puedes ser tú.

Y lo paradójico es que solo cuando eres tú de verdad ---sin filtros, sin máscaras, sin poses--- el otro puede amarte de verdad. Porque ama lo que eres, no lo que finges ser.

La indisolubilidad no limita la libertad. La hace posible.

\section{La biología ya lo sabía}

Si has leído la primera parte de este libro, quizá estés pensando: «Esto me suena.» Y con razón.

Cuando hablamos de la química del vínculo,\footnote{Cf. capítulo 3 de este libro.} vimos que el enamoramiento ---esa tormenta de dopamina, noradrenalina y serotonina baja--- tiene fecha de caducidad. Dura entre uno y tres años, aproximadamente. Pasado ese tiempo, el cóctel cambia. Y lo que aparece no es vacío, sino algo distinto y más profundo: el apego.

El apego se construye con oxitocina y vasopresina, las hormonas de la confianza, la calma y la permanencia. Son las mismas hormonas que vinculan a una madre con su bebé. Y hacen algo que la dopamina del enamoramiento no puede hacer: crear un vínculo duradero, sereno, basado no en la excitación sino en la seguridad.

La biología, literalmente, nos empuja de la pasión al apego. Del «no puedo vivir sin ti» al «quiero construir contigo». Del terremoto a la tierra firme.

¿No es curioso? Nuestro propio cuerpo está diseñado para la permanencia. La evolución ha invertido millones de años en crear un sistema que nos lleva, paso a paso, del deseo inicial al vínculo estable. Porque nuestras crías ---esas criaturas indefensas que necesitan años de cuidado--- necesitan padres que se queden. No padres que se enamoren un rato y luego se vayan.

La indisolubilidad no es una ocurrencia de la Iglesia. Es la palabra que la fe le pone a algo que el cuerpo ya sabe.

\section{El regalo que transforma}

Hay algo más en la indisolubilidad, algo que es fácil de pasar por alto y que a mí me parece decisivo.

Cuando alguien te promete que no te va a dejar nunca, esa promesa te cambia. No en abstracto, no en teoría. Te cambia de verdad. Porque al recibirla, nace en ti algo que antes no estaba: el deseo de merecerla.

Nadie te lo pide. Nadie te lo exige. Pero cuando alguien ha puesto su vida entera en tus manos ---cuando ha apostado todo lo que tiene a una carta, y esa carta eres tú---, algo dentro de ti se moviliza. Quieres estar a la altura. No por obligación, sino por amor. No porque te castiguen si fallas, sino porque no quieres fallar a quien te ha dado tanto.

«Me ha prometido que se queda. Pues yo quiero ser alguien que merezca que se quede.»

¿Ves cómo funciona? La indisolubilidad no es un mandato que te aplasta. Es un regalo que te eleva. Te convierte en mejor persona, no porque te obligue, sino porque te inspira. Como el que recibe una confianza enorme y, precisamente por eso, se esfuerza en no defraudarla.

Es lo que ocurre con todas las grandes promesas. El padre que le dice a su hijo «siempre estaré aquí» no se está poniendo una cadena: se está dando una razón para ser fuerte. El amigo que dice «cuenta conmigo pase lo que pase» no está firmando una condena: está eligiendo ser mejor persona. La promesa de permanencia no te encadena a la miseria. Te empuja hacia arriba.

\section{Pero ¿y si no funciona?}

Lo prometí: íbamos a volver a esta pregunta. Porque no sería honesto esquivarla.

No todo matrimonio es un camino de rosas. Hay matrimonios donde hay dolor real, incluso violencia. Hay situaciones de las que hay que salir, y la Iglesia lo sabe. No estoy diciendo que haya que aguantar cualquier cosa en nombre de la indisolubilidad. Eso sería una caricatura ---y una crueldad---.

Pero fijaos en algo que pasa constantemente y que dice más que cualquier argumento teórico. Cuando una pareja se separa ---da igual si fue de mutuo acuerdo o entre abogados, da igual si se llevan bien o mal---, escuchadlos hablar. Prestad atención a cómo se refieren al otro. Es llamativo: muy a menudo, incluso años después de la separación, siguen diciendo «mi marido», «mi mujer». No «mi ex». No «la persona con la que estuve». \emph{Mi marido. Mi mujer.} Lo dicen con cariño o con rabia, con nostalgia o con amargura, pero lo dicen. Algo en ellos ---más fuerte que la distancia, más fuerte que el conflicto, más fuerte que los papeles del juzgado--- sigue reconociendo un vínculo que no han podido romper del todo. Como si el lenguaje supiera lo que la razón a veces no quiere admitir: que lo que hubo entre ellos no fue un contrato que se rescindió, sino un lazo que, de alguna manera misteriosa, sigue ahí. El vínculo es más fuerte de lo que pensamos. Más fuerte, incluso, de lo que quienes lo viven están dispuestos a reconocer.

Pero sería deshonesto dejar las cosas aquí, en la anécdota entrañable del lenguaje que delata un vínculo. Porque la realidad del fracaso matrimonial es más compleja y, a veces, más oscura.

Hay matrimonios que se rompen y que dejan heridas profundas. Hay personas que han vivido la indisolubilidad no como un regalo sino como una losa, como una condena a permanecer en algo que les hacía daño. Hay separaciones que no son fracasos sino actos de dignidad: la dignidad de quien dice «hasta aquí» ante la violencia, el desprecio o la destrucción sistemática de la persona. Y hay que tener la honestidad de nombrarlo.

La Iglesia, que a veces parece hablar solo del ideal, sabe esto perfectamente. La separación física está prevista en el derecho canónico precisamente para proteger a quien sufre.\footnote{CIC, can. 1153 §1: «Si uno de los cónyuges pone en grave peligro espiritual o corporal al otro o a la prole, o de otro modo hace demasiado dura la vida en común, proporciona al otro un motivo legítimo para separarse.»} Y existe un proceso de nulidad que no es un divorcio católico encubierto, sino algo muy distinto: la constatación, tras un examen riguroso, de que lo que parecía un matrimonio no llegó a serlo nunca ---porque faltó libertad, madurez suficiente, o un consentimiento verdadero---.\footnote{Cf. CIC, cann. 1095--1107. La declaración de nulidad no dice «este matrimonio ha terminado», sino «este matrimonio nunca existió válidamente». Es una distinción sutil pero fundamental.} No se trata de buscar una puerta trasera, sino de reconocer una verdad.

Y hay algo más que conviene decir, porque se dice poco: el fracaso de un matrimonio no es el final de la historia de amor de esa persona. Como en cualquier otra dimensión de la vida ---la profesional, la vocacional, la personal---, el fracaso puede ser una herida, pero también un punto de partida. Hay personas separadas que han encontrado, a través de ese dolor, una profundidad humana y espiritual que no tenían antes. El Papa Francisco lo dice con una claridad que la Iglesia no siempre ha tenido: esas personas «no están excomulgadas», y la comunidad cristiana debe ser para ellas «hospital de campaña» y no «aduana».\footnote{Francisco, \textit{Amoris Laetitia}, nn. 243, 291, 296--297 (2016). Todo el capítulo VIII de la exhortación está dedicado a «acompañar, discernir e integrar la fragilidad».}

Nada de esto debilita la indisolubilidad. La hace más real, más humana, más honesta. Porque la promesa «para siempre» no se pronuncia en un mundo perfecto sino en este, donde las personas son frágiles, cometen errores y a veces se equivocan incluso en lo más importante. Reconocer esa fragilidad no es renunciar al ideal; es amarlo con los ojos abiertos.

Lo que estoy diciendo es otra cosa. Estoy diciendo que la pregunta «¿y si no funciona?», cuando se hace \emph{antes} de casarse, como miedo preventivo, como excusa para no comprometerse, revela algo importante: revela que estamos mirando la indisolubilidad desde el lado equivocado.

La estamos mirando como una cláusula contractual que te atrapa. Pero no es eso. Es una declaración de amor que te sostiene.

¿Y si, en lugar de preguntar «qué pasa si no funciona», preguntáramos «qué pasa si sí funciona»? ¿Qué pasa si alguien de verdad cumple esa promesa? ¿Qué pasa si tienes la experiencia de ser amado así, sin reservas, sin salida de emergencia, sin condiciones?

Te diré lo que pasa: que tu vida se transforma. Que por primera vez puedes bajar la guardia del todo. Que descubres una libertad que no habías conocido: la libertad de ser amado no por lo que haces, sino por lo que eres. La libertad de no tener que ganarte cada día el derecho a ser querido.

Y esa experiencia ---la de ser amado incondicionalmente por alguien que ha elegido quedarse--- es quizá la imagen más cercana que tenemos del amor de Dios.\footnote{«El amor conyugal comporta una totalidad en la que entran todos los elementos de la persona ---reclamo del cuerpo y del instinto, fuerza del sentimiento y de la afectividad, aspiración del espíritu y de la voluntad---; mira a una unidad profundamente personal que, más allá de la unión en una sola carne, conduce a no tener más que un corazón y un alma.» Juan Pablo II, \textit{Familiaris Consortio}, n. 13.}

\section{Una promesa que viene de más arriba}

Quiero terminar con algo que quizá ya intuyes.

Si la indisolubilidad es un regalo tan grande, ¿de dónde viene? ¿Es solo una decisión humana, un acto de voluntad heroica? ¿O hay algo más?

Cuando Jesús habla del matrimonio, no inventa nada nuevo. Hace algo mucho más audaz: vuelve al principio. Le preguntan por el divorcio, y en lugar de discutir la ley de Moisés, remite a sus oyentes al Génesis: «Al principio no fue así.»\footnote{Mt 19,8. Los fariseos le preguntan por qué Moisés permitió el divorcio, y Jesús responde: «Por la dureza de vuestro corazón os permitió Moisés repudiar a vuestras mujeres; pero al principio no fue así.»} Al principio, cuando Dios pensó el amor entre un hombre y una mujer, lo pensó así: para siempre. No como castigo. No como prueba. Como regalo.

«Lo que Dios ha unido, que no lo separe el hombre.»\footnote{Mt 19,6.}

Fíjate: no dice «lo que la Iglesia ha unido». Dice «lo que Dios ha unido». El vínculo no lo crea la ceremonia, ni el sacerdote, ni el registro civil. El vínculo viene de arriba. Y precisamente por eso tiene una solidez que ninguna promesa puramente humana podría tener.

Esto es lo que la tradición cristiana llama el vínculo sacramental. Cuando dos bautizados se casan, su unión no solo es un compromiso entre ellos: es un signo del amor de Cristo por su Iglesia.\footnote{Cf. Ef 5,31-32. El Concilio de Trento definió la sacramentalidad del matrimonio y su indisolubilidad intrínseca. Cf. Concilio de Trento, Sesión XXIV (1563), Decreto \textit{Tametsi}. El \textit{Catecismo de la Iglesia Católica}, n. 1640, dice: «El vínculo matrimonial es establecido por Dios mismo, de modo que el matrimonio celebrado y consumado entre bautizados no puede ser disuelto jamás.»} Un amor que no tiene fecha de caducidad. Un amor que no se retira cuando el otro falla. Un amor que permanece.

Y la buena noticia es esta: no estáis solos en esa promesa. No se trata de que vosotros, con vuestra sola fuerza de voluntad, tengáis que sostener un compromiso de por vida. Se trata de que Dios mismo está en ese vínculo, sosteniéndolo desde dentro. La gracia del sacramento no es un adorno: es una fuerza real que trabaja en vuestro matrimonio, que os empuja a amar cuando ya no os quedan fuerzas, que os sostiene cuando todo parece venirse abajo.\footnote{«La gracia propia del sacramento del Matrimonio está destinada a perfeccionar el amor de los cónyuges, a fortalecer su unidad indisoluble. Por medio de esta gracia "se ayudan mutuamente a santificarse".» \textit{Catecismo de la Iglesia Católica}, n. 1641.}

No os piden que seáis héroes. Os ofrecen una mano que os sostiene.

\section{La promesa más hermosa}

Déjame resumirlo con una imagen.

Imagina dos formas de cruzar un río. En la primera, cruzas sobre piedras sueltas. Cada paso es incierto. No sabes si la próxima piedra aguantará tu peso. Avanzas con miedo, midiendo, calculando. Si la piedra se mueve, te caes.

En la segunda, cruzas por un puente de piedra. Sólido. Firme. Construido para durar. No tienes que pensar en cada paso. Puedes mirar alrededor, disfrutar del paisaje, caminar con confianza. El puente aguanta.

La indisolubilidad es el puente.

No te quita la libertad de caminar. Te da una superficie firme sobre la que hacerlo. No te obliga a llegar al otro lado: te permite llegar.

Y lo más hermoso de todo es que ese puente no te lo construyes tú solo. Lo construís juntos. Piedra a piedra. Día a día. Con cada «me quedo» silencioso, con cada perdón concedido, con cada mañana en la que decides amar aunque no tengas ganas. Y sobre ese puente, sostenido por una gracia que viene de más arriba, podéis caminar seguros. Juntos. Hasta el final.

Esa es la promesa que el mundo no sabe hacer. La promesa que suena a cadena y resulta ser alas. La promesa que nadie te va a hacer\ldots{} excepto la persona que de verdad te ama.

\paracaminar

\begin{enumerate}
  \item ¿Cómo percibíais la indisolubilidad antes de leer este capítulo: como una carga o como un regalo? ¿Ha cambiado algo en vuestra forma de verla? Hablad con sinceridad.

  \item Pensad en vuestra relación: ¿hay momentos en los que os habéis «puesto la máscara» por miedo a que el otro no aceptara lo que veía? ¿Qué pasaría si pudierais quitárosla del todo, con la seguridad de que el otro no se va?

  \item Intentad deciros, mirándoos a los ojos, la promesa de este capítulo con vuestras propias palabras. No tiene que ser solemne ni perfecto. Solo verdadero.
\end{enumerate}

\vinetacapitulo{ch08}

\chapter{Creer en ti más que tú mismo}

\epigraph{La fidelidad no nace de una prohibición. Nace de una fe.}{}

Vamos con la segunda palabra difícil. Si en el capítulo anterior hablamos de la indisolubilidad y descubrimos que, mirada desde el lado correcto, es el regalo más grande que alguien puede hacerte, ahora toca hablar de algo que suele provocar todavía más resistencia: la fidelidad.

Y es que la fidelidad tiene mala prensa. No entre quienes ya llevan años casados y han descubierto lo que significa. Sino entre quienes la miran desde fuera, desde antes, desde la cultura que nos rodea. Para muchos, fidelidad es sinónimo de restricción. De renuncia. De jaula.

El argumento suele formularse más o menos así: «¿Me estás diciendo que no puedo estar con nadie más? ¿En toda mi vida? ¿Nunca?» Dicho así, suena a condena. Como si el matrimonio fuera una puerta que se cierra ---para siempre--- sobre un mundo de posibilidades. Como si casarse fuera aceptar una limitación que nadie en su sano juicio querría.

Y si la fidelidad fuera eso ---una prohibición sexual disfrazada de virtud---, tendrían razón en resistirse. Nadie quiere vivir encerrado. Nadie quiere que le digan lo que no puede hacer durante el resto de su vida.

Pero ¿y si la fidelidad no fuera eso? ¿Y si fuera exactamente lo contrario?

\section{Una palabra que viene de la fe}

Hagamos algo que siempre ayuda: ir a la raíz. La palabra \textit{fidelidad} viene del latín \textit{fidelitas}, que a su vez viene de \textit{fides}. Y \textit{fides} no significa prohibición, ni restricción, ni renuncia. \textit{Fides} significa \textbf{fe}.\footnote{El término latino \textit{fides} abarca un campo semántico amplio: fe, confianza, lealtad, crédito, palabra dada. En el derecho romano, la \textit{fides} era el fundamento de todo pacto: la confianza en que el otro cumplirá lo acordado. Cf. Benveniste, É., \textit{Le vocabulaire des institutions indo-européennes}, vol. 1, Minuit, 1969.}

Fidelidad es, literalmente, un acto de fe. No fe en Dios ---aunque también---, sino fe en una persona. En ti.

Cuando alguien te es fiel, lo que te está diciendo no es «me resigno a no estar con otros». Lo que te está diciendo es: «Creo en ti. Creo en lo que eres. Creo en lo que puedes llegar a ser. Y apuesto mi vida entera a esa creencia.»

Es una declaración de fe, no un acto de renuncia.

¿Y si fuera al revés?

¿Y si, en lugar de mirar la fidelidad como algo que tú \emph{das} ---renunciando, sacrificándote, privándote---, la miráramos como algo que tú \emph{recibes}? Porque cuando tu mujer te es fiel, o cuando tu marido te es fiel, lo que está ocurriendo no es que se estén privando de algo. Lo que está ocurriendo es que alguien está depositando en ti una fe absoluta. Una fe que dice: «De entre todos los seres humanos del mundo, yo creo que tú eres la persona con la que quiero recorrer la vida. Y creo que tú eres capaz de ser el hombre, la mujer, que yo he visto en ti.»

Eso no es una jaula. Eso es el mayor voto de confianza que un ser humano puede recibir.

\section{La fe que te hace mejor}

Piensa un momento en las personas que han creído en ti a lo largo de tu vida.

Quizá fue un profesor que vio algo en ti que tú no veías. Que te dijo, cuando estabas a punto de tirar la toalla: «Tú puedes con esto.» Y algo cambió. No porque de repente fueras más listo o más capaz. Sino porque alguien te miró con ojos de fe, y esa mirada te hizo crecer.

Quizá fue un entrenador que te puso en el equipo cuando nadie más apostaba por ti. O un padre, o una madre, que nunca dejó de confiar en ti, ni siquiera cuando tú mismo habías dejado de hacerlo. O un amigo que, en tu peor momento, te dijo: «Yo sé quién eres de verdad. Y esto no eres tú.»

Todos hemos tenido, si hemos tenido suerte, al menos a una persona así. Alguien que creyó en nosotros más de lo que nosotros creíamos en nosotros mismos. Y todos sabemos lo que eso produce: no quieres fallarle. No porque te castigue si lo haces. No por miedo. Sino porque su fe en ti te eleva. Porque cuando alguien te mira así, traicionar esa confianza sería traicionarte a ti mismo.

Ahora multiplica eso por la relación más íntima de tu vida.

Cuando tu esposa te dice «sí, quiero», lo que está diciendo es: «Creo en ti con todo lo que tengo. Creo en ti con mi cuerpo, con mi futuro, con mis hijos, con mis años. Te entrego la fe más grande que puedo entregar: la fe de que tú vas a estar a la altura.» Cuando tu esposo te dice «sí, quiero», está haciendo exactamente lo mismo.

Eso no es una restricción. Es una responsabilidad enorme, sí. Pero es la clase de responsabilidad que te hace mejor persona. Porque la fe que alguien deposita en ti tiene un poder que nada más tiene: te muestra una versión de ti mismo que quizá no conocías. Y te hace querer ser esa versión.

\section{«No puedo» y «no quiero fallarte»}

Hay una diferencia abismal entre dos frases que, desde fuera, pueden parecer iguales.

La primera: «No puedo estar con otra persona.» La segunda: «No quiero fallarte.»

La primera es prohibición. Viene de fuera. Me la imponen. La obedezco ---si la obedezco--- a regañadientes, porque es la norma, porque es lo que toca, porque me metí en esto y ahora tengo que aguantar. La primera genera resentimiento. Acumula frustración. Convierte la fidelidad en una carga que pesa un poco más cada año.

La segunda es respuesta. Viene de dentro. Nace de haber recibido algo tan grande que no quiero destrozarlo. La segunda no necesita vigilancia ni control. No necesita que nadie me vigile, porque no es miedo lo que me mueve: es amor. Es la consciencia de que alguien me ha confiado lo más valioso que tiene ---su fe en mí--- y yo no quiero ser indigno de eso.

¿Ves la diferencia? En un caso, la fidelidad es una ley que me limita. En el otro, es una respuesta libre a un amor que me eleva. En un caso, soy un preso. En el otro, soy alguien que ha sido creído, y que elige estar a la altura de esa fe.

La fidelidad cristiana no funciona como una prohibición. Funciona como una correspondencia.\footnote{Benedicto XVI lo expresó con claridad al hablar del amor: «Él nos ha amado primero y sigue amándonos primero; por eso, nosotros podemos corresponder con el amor.» Benedicto XVI, Carta encíclica \textit{Deus Caritas Est}, n. 17 (2005). El mismo principio se aplica a la fidelidad: es respuesta a una fe recibida, no cumplimiento de una norma impuesta.} No digo «no» a otros porque me lo prohíban. Digo «sí» a ti, cada día, porque tu fe en mí merece mi mejor versión.

\section{La confianza como suelo firme}

Si habéis leído el capítulo 2, recordaréis lo que vimos sobre el apego. Bowlby descubrió que un niño necesita una \textit{base segura}: alguien predecible, constante, que esté ahí cuando lo necesite. Desde esa seguridad, el niño puede explorar el mundo. Sin ella, se paraliza o se desespera.\footnote{Cf. capítulo 2 de este libro, especialmente la sección sobre Bowlby y la teoría del apego.}

La fidelidad es la versión adulta, consciente y elegida de esa misma dinámica.

Cuando sabes que tu pareja te es fiel ---no porque le obligan, sino porque cree en ti y en lo vuestro---, ocurre algo parecido a lo que le ocurre al niño con apego seguro: puedes soltar la guardia. Puedes mostrarte como eres, con tus miserias, con tus contradicciones, con esas partes de ti que no enseñas a nadie. Puedes ser vulnerable sin miedo, porque sabes que el otro no va a usar tu vulnerabilidad como excusa para marcharse.

La biología ya nos inclinaba hacia ahí. Las hormonas del vínculo ---la oxitocina, la vasopresina--- nos empujan a la permanencia, a la confianza, al cuidado.\footnote{Cf. capítulo 3 de este libro, sobre la oxitocina y la vasopresina como hormonas del apego de pareja.} Pero la biología es un empujón, no una decisión. La fidelidad toma ese empujón biológico y lo convierte en promesa libre. Lo eleva de instinto a elección. De química a don.

Y ahí está la conexión con lo que vimos en el capítulo anterior: la indisolubilidad crea el espacio seguro; la fidelidad lo habita. La promesa de que no te vas a ir crea el marco; la fe que depositas en el otro cada día lo llena de contenido. No son dos cosas distintas. Son dos caras de lo mismo: un amor que dice «estoy aquí, creo en ti, y no me voy a ningún sitio».

\section{Lo que la infidelidad destruye de verdad}

Quizá ahora se entienda mejor por qué la infidelidad duele tanto. No es solo que duela el engaño, aunque duele. No es solo que duelan los celos, aunque duelen. Lo que de verdad destroza la infidelidad es que destruye la fe.

Cuando alguien te es infiel, lo que se rompe no es un acuerdo. Lo que se rompe es la creencia de que eras suficiente. La confianza de que el otro te miraba y veía a alguien por quien merecía la pena apostar la vida entera. Lo que se rompe es esa voz que decía «creo en ti» y que de repente se queda muda.

Por eso la infidelidad se siente como una traición ---porque lo es---. Pero no es solo una traición al pacto. Es una traición a la persona. A la fe que esa persona había puesto en ti. A la versión de ti mismo que el otro te ayudaba a ser.

Y por eso ---no por moralismo, no por rigidez--- la Iglesia la toma tan en serio.\footnote{«La fidelidad expresa la constancia en mantener la palabra dada. Dios es fiel. El sacramento del Matrimonio hace entrar al hombre y a la mujer en la fidelidad de Cristo a su Iglesia.» \textit{Catecismo de la Iglesia Católica}, n. 2365.} No porque le importe controlar la vida sexual de nadie, sino porque sabe lo que está en juego: la fe de una persona en otra persona. Y esa fe es sagrada.

\section{Creer en ti cuando tú no crees en ti}

Hay algo más, y quizá sea lo más hermoso de todo.

La fidelidad no consiste solo en creer en el otro cuando es fácil hacerlo. Cuando las cosas van bien, cuando estáis enamorados, cuando todo fluye. Eso no requiere mucha fe: es como creer que el sol saldrá mañana.

La verdadera fe se prueba cuando es difícil creer. Cuando el otro ha fallado. Cuando ha mostrado su peor versión. Cuando te ha decepcionado o se ha decepcionado a sí mismo. Cuando piensa que no vale la pena, que no es capaz, que no da la talla.

En esos momentos, la fidelidad es decir: «Yo sigo creyendo en ti. Aunque tú ahora mismo no creas en ti, yo sí. Yo sé quién eres. He visto lo mejor de ti, y sé que eso es lo real, aunque ahora esté escondido.»

Eso es fidelidad. No la fidelidad fácil de los días buenos, sino la fidelidad tozuda de los días malos. La que sostiene al otro cuando no puede sostenerse solo. La que le recuerda quién es cuando se le ha olvidado.

¿No es eso, en el fondo, lo que Dios hace con nosotros? ¿No es eso lo que significa que Dios sea fiel? Que sigue creyendo en nosotros cuando nosotros hemos dejado de creer. Que nos mira con ojos de fe incluso cuando nosotros nos miramos con ojos de vergüenza. Que su fidelidad no depende de nuestro merecimiento, sino de su amor.\footnote{«Si somos infieles, Él permanece fiel, porque no puede negarse a sí mismo» (2~Tim 2,13). La fidelidad de Dios no es una respuesta a nuestra bondad, sino una expresión de su ser. Y el matrimonio cristiano está llamado a reflejar, imperfectamente pero realmente, esa misma fidelidad. Cf. Juan Pablo II, Exhortación apostólica \textit{Familiaris Consortio}, n. 20 (1981).}

El matrimonio nos invita a algo extraordinario: a ser para el otro ese espejo de Dios. A creer en el otro como Dios cree en nosotros. No porque el otro sea perfecto ---no lo es, y tú tampoco---, sino porque la fe no se basa en la perfección del otro. Se basa en el amor que eliges tener.

\section{La fidelidad que se elige cada mañana}

Una última cosa. La fidelidad no es algo que se decide una vez, el día de la boda, y que luego funciona sola, como un piloto automático. La fidelidad se elige. Y se elige cada día.

Cada mañana, al despertar junto a la misma persona, hay una elección silenciosa que casi nunca se verbaliza pero que lo sostiene todo: «Hoy, otra vez, creo en ti. Hoy, otra vez, apuesto por nosotros.» No es una elección heroica. No sale en las noticias. Nadie te aplaude por ella. Pero es la elección más importante que vas a tomar en todo el día.

Y aquí está la paradoja hermosa de la fidelidad: cuanto más la eliges, más fácil se vuelve. No porque te acostumbres ---la costumbre mata---. Sino porque cada día que eliges creer en el otro, cada día que respondes a su fe con tu fe, el vínculo se hace más fuerte. La confianza se hace más honda. Y lo que empezó como una decisión se convierte, poco a poco, en una segunda naturaleza. En una forma de estar en el mundo. En la persona que eres.\footnote{Santo Tomás de Aquino describe la fidelidad como un hábito de la voluntad: no basta un acto aislado, sino la disposición estable a mantener la fe prometida. La virtud, en la tradición tomista, no es una imposición sino una perfección de la libertad. Cf. \textit{Summa Theologiae}, II-II, q. 110, a. 3.}

El día de la boda, la fidelidad es una promesa. Diez años después, es una historia. Treinta años después, es una identidad. No es que seas fiel porque prometiste serlo. Es que eres fiel porque en eso te has convertido, día a día, elección a elección, fe a fe.

Y eso ---convertirte en alguien fiel--- es mucho más grande que simplemente «no estar con otros». Es haberte transformado en la clase de persona que cree en otra persona con todo lo que tiene. Y eso, si lo piensas bien, se parece bastante a lo que significa amar.

\section{Porque debo, porque quiero, porque te quiero}

Hay una progresión que lo resume todo. Tres razones para hacer las cosas, y cada una supone un salto respecto a la anterior.

La primera: \textit{porque debo}. Es la razón de la norma, de la ley, del deber externo. Hago las cosas porque me las mandan, porque es lo que toca, porque las consecuencias de no hacerlas serían peores. Es el nivel más bajo. Cumples, pero no amas. Es como pagar impuestos: necesario, pero no te hace mejor persona.

La segunda: \textit{porque quiero}. Aquí hay un salto enorme. Ya no es imposición: es elección. Es la libertad en acto. Decía san Josemaría Escrivá que ``porque quiero'' es la razón más sobrenatural que existe,\footnote{Cf. capítulo 6 de este libro.} y tenía razón: en esas dos palabras se manifiesta la libertad que Dios nos ha dado, la misma que nos hace imagen suya. ``Porque quiero'' transforma el deber en don.

Pero hay una tercera, y es la que cambia las reglas del juego: \textit{porque te quiero}. Ya no es solo libertad ---es libertad orientada. Dirigida a una persona concreta. Con nombre, con cara, con defectos que conoces de memoria y virtudes que solo tú has visto. ``Porque te quiero'' es el amor que deja de ser abstracto y se encarna en alguien.

Y ahí está el misterio: ``porque te quiero'' se parece mucho a cómo ama Dios. Porque Dios no dice ``amo al mundo'' en genérico, como quien firma un manifiesto. Dios ama a cada uno de nosotros en particular, por nuestro nombre, de un modo irrepetible e insustituible.\footnote{«No temas, que yo te he rescatado, te he llamado por tu nombre y eres mío» (Is 43,1). El amor de Dios no es un amor de masas: es un amor personal, dirigido, con nombre propio.} El cónyuge hace exactamente lo mismo. No ama ``el amor''. No ama una idea. Te ama a ti. Con todo lo que eso implica.

La fidelidad recorre ese camino. Empieza en el ``debo'' de quien no conoce otra razón. Madura en el ``quiero'' de quien descubre la libertad. Y culmina en el ``te quiero'' de quien ha aprendido que el amor más alto no es el que se da a todos por igual, sino el que se entrega entero a una persona ---esta persona, aquí, ahora, para siempre---.

\paracaminar

\begin{enumerate}
  \item ¿Alguna vez alguien creyó en vosotros cuando vosotros no creíais en vosotros mismos? ¿Qué efecto tuvo esa fe? ¿Cómo os hizo sentir? Compartidlo.

  \item ¿Hay diferencia, en vuestra experiencia, entre el «no puedo» de la prohibición y el «no quiero fallarte» de la correspondencia? ¿Desde cuál de los dos vivís ---o queréis vivir--- la fidelidad?

  \item La fidelidad se elige cada día. ¿Qué gestos concretos y cotidianos os ayudan a renovar esa elección? ¿Qué podríais hacer para que vuestra fe mutua sea más visible, más tangible?
\end{enumerate}

\vinetacapitulo{ch09}

\chapter{Sal de la tierra}

\epigraph{«Vosotros sois la sal de la tierra. [\ldots] Vosotros sois la luz del mundo.»}{Mateo 5,13-14}

Si hay algo que la gente cree saber sobre lo que la Iglesia piensa del matrimonio, es esto: que quiere que tengas muchos hijos. Cuantos más, mejor. Sin plantearse nada, sin pensárselo dos veces. Abrir la puerta a la vida ---dicen--- significa no cerrarla nunca, como si la fecundidad fuera una cuestión de aritmética y el valor de un matrimonio se midiera por el número de sillas que necesita alrededor de la mesa.

La caricatura es conocida. Y como toda caricatura, tiene un rasgo reconocible ---algo de verdad--- rodeado de exageración. Sí, la Iglesia valora la apertura a la vida. Sí, los hijos son un bien enorme. Pero reducir la fecundidad del matrimonio a la cantidad de hijos es como reducir la calidad de un árbol al número de frutos que da en una temporada, sin mirar si da sombra, si sus raíces sujetan la tierra, si sus ramas acogen nidos.

La fecundidad de un matrimonio es algo mucho más ancho, mucho más hondo y mucho más hermoso de lo que esa caricatura sugiere. Y para verlo, hay que cambiar la mirada.

\section{La familia que todos recuerdan}

Dejadme que os cuente algo que probablemente habéis vivido. Pensad un momento en las familias que habéis conocido a lo largo de vuestra vida. No las familias perfectas ---esas no existen fuera de los anuncios de Navidad---. Las familias \emph{reales}. Las que os han marcado.

Seguro que hay alguna. Una casa a la que siempre os apetecía ir. Donde os sentíais acogidos sin necesidad de avisar. Donde la puerta estaba abierta y la mesa se estiraba sin drama para un comensal más. Una casa donde se reían mucho, donde los padres discutían a veces pero se miraban con una complicidad que daba gusto ver. Una familia donde los amigos de los hijos acababan quedándose a dormir y, con el tiempo, siendo casi de la familia.

¿La tenéis? Bien. Ahora fijaos en algo curioso: esa familia que recordáis con cariño, esa familia que os hizo sentir que el mundo era un lugar más habitable, probablemente no hacía nada espectacular. No organizaba grandes eventos. No tenía un programa de acción social. Simplemente \emph{vivía} de una manera que generaba vida a su alrededor. Como la sal, que no se ve en el plato pero cambia todo el sabor.

Eso es fecundidad.

\section{¿Y si fuera al revés?}

Porque la percepción habitual ---fecundidad igual a hijos, punto--- no solo es reduccionista. Es que, a veces, el enfoque se invierte hasta decir lo contrario de lo que debería.

Pensemos en un matrimonio que tiene cinco hijos pero que vive encerrado en sí mismo. Una familia que no se abre, que no acoge, que no comparte. Una casa donde no entra nadie de fuera. Unos padres que no tienen tiempo para nadie más ---ni para ellos mismos---, no porque estén ocupados dándose, sino porque han convertido la familia en una fortaleza con el puente levadizo siempre alzado. ¿Es fecundo ese matrimonio? Tiene frutos, sí. Pero ¿da vida?

Ahora pensemos en otro matrimonio. Quizá tiene dos hijos, quizá tiene uno, quizá ---por razones que solo ellos y Dios conocen--- no ha podido tener ninguno. Pero es una casa donde la gente quiere estar. Una pareja que acompaña a otros matrimonios en crisis. Unos esposos que acogen al sobrino que lo está pasando mal, que visitan al vecino enfermo, que abren su salón para que un grupo de amigos pueda reunirse a rezar o simplemente a compartir la vida. Una familia que \emph{irradia} algo que los demás reconocen sin poder explicarlo del todo: alegría, sentido, calor, esperanza.

¿Cuál de los dos es más fecundo?

La pregunta no pretende enfrentar a las familias grandes con las pequeñas. No se trata de eso. Se trata de entender que la fecundidad no se mide en números, sino en vida generada. En cuánta vida brota alrededor de ese matrimonio. En cuántas personas encuentran en esa familia un refugio, un referente, una razón para creer que el amor de verdad es posible.\footnote{«La fecundidad del amor matrimonial no se reduce a la sola procreación de los hijos, sino que debe extenderse a su educación moral y a su formación espiritual. [\ldots] Los esposos, a los que Dios no ha concedido tener hijos, pueden, no obstante, tener una vida conyugal plena de sentido, humana y cristianamente.» \textit{Catecismo de la Iglesia Católica}, n. 1654.}

\section{El amor que desborda}

Hay una ley no escrita que se cumple siempre: el amor verdadero no cabe dentro.

Cuando dos personas se aman de verdad ---no con el amor de las películas, que mira al otro como si el resto del mundo no existiera, sino con el amor maduro, el que elige y permanece---, ese amor tiende a salir. A desbordarse. A buscar a quién más darse. No porque los esposos se aburran el uno del otro, sino porque el amor es así: cuanto más auténtico, más expansivo. Cuanto más hondo, más ancho.

Los hijos son la expresión más natural de ese desbordamiento. No nacen porque «hay que tenerlos», ni porque «la Iglesia manda». Nacen porque el amor de dos personas se hace tan real, tan encarnado, que se convierte en una tercera persona. El hijo no es un proyecto ni una obligación: es el amor que cobra rostro, el abrazo que se hace carne.\footnote{«El hijo no viene de fuera a añadirse al amor mutuo de los esposos; brota del corazón mismo de ese don recíproco, del que es fruto y cumplimiento.» Juan Pablo II, Exhortación apostólica \textit{Familiaris Consortio}, n. 14 (1981).} Hay algo sobrecogedor en eso, si se piensa despacio: que dos personas puedan amarse tanto que de su amor nazca alguien nuevo, alguien que nunca antes existió, alguien que vivirá para siempre. Pocas cosas en la experiencia humana se acercan más al acto creador de Dios.

Pero el desbordamiento no se agota en los hijos. El amor fecundo sigue buscando cauces, sigue inventando formas de dar vida.

\section{La fecundidad que no se cuenta}

La tradición cristiana ha hablado siempre de varias dimensiones de la fecundidad, aunque a veces las hayamos olvidado.\footnote{El Concilio Vaticano II amplía la comprensión de la fecundidad matrimonial más allá de la procreación: «El matrimonio y el amor conyugal están ordenados por su propia naturaleza a la procreación y educación de la prole. [\ldots] Sin embargo, el matrimonio no ha sido instituido solamente para la procreación.» \textit{Gaudium et Spes}, n. 48 y 50.}

Hay una \textbf{fecundidad educativa}: la que forma personas. No basta con traer hijos al mundo; hay que ayudarles a crecer, a pensar, a ser libres, a amar. Un padre que enseña a su hijo a pedir perdón está siendo fecundo. Una madre que le enseña a mirar al que sufre está generando vida. Y no solo para sus hijos: la educación que una familia transmite se extiende a quienes la rodean, porque los valores se contagian como el acento.

Hay una \textbf{fecundidad espiritual}. Es la del matrimonio que reza junto, que busca a Dios, que transmite la fe no como una lección teórica sino como algo vivo, experimentado, agradecido. Es también la del matrimonio que acompaña a otros en su camino de fe ---parejas que preparan a novios para el matrimonio, que animan grupos de oración, que simplemente testimonian con su vida que creer tiene sentido---.\footnote{El papa Francisco habla de la «fecundidad ampliada» de la familia: «Una familia que no acoge y no sale de sí misma deja de ser familia.» Cf. Francisco, Exhortación apostólica \textit{Amoris Laetitia}, n. 181 (2016).}

Hay una \textbf{fecundidad social}: la familia que se implica en su comunidad, que participa, que aporta. La que acoge al inmigrante, la que cuida al anciano, la que se compromete con lo que pasa en su barrio. La familia no es una isla. Es ---debería ser--- un foco de calor que derrite un poco el hielo de alrededor.

Y hay una fecundidad que quizá no tiene nombre técnico pero que es la más visible de todas: la \textbf{fecundidad del testimonio}. Es la que se produce cuando alguien mira a una pareja y piensa: «Yo quiero eso.» No porque esa pareja sea perfecta, sino porque es real. Porque se nota que se quieren, que se respetan, que luchan juntos, que se perdonan, que disfrutan estando juntos. Esa envidia sana que despiertan ciertos matrimonios es, quizá, la forma más poderosa de evangelización que existe. No hace falta dar discursos. No hace falta repartir folletos. Basta con \emph{vivir} de manera que la gente mire y vea que el amor para siempre no es una fantasía de ingenuos, sino algo posible, real, deseable.

\section{La casa de la puerta abierta}

Me gusta pensar en la familia fecunda como una casa con la puerta abierta. No abierta de par en par a cualquier hora, que eso no es acogida, es caos. Pero sí abierta de verdad: una casa donde el que llega sabe que es bienvenido.

Todos hemos conocido casas así. Casas donde el sofá siempre tiene sitio para uno más. Donde la conversación fluye y nadie mira el reloj. Donde los hijos traen amigos con naturalidad y los amigos traen a los amigos de los amigos. Casas donde se ha llorado y se ha reído, donde se ha discutido a gritos y se ha hecho las paces en voz baja, donde la vida ha entrado con todo su peso y toda su belleza.

Esas casas no nacen por casualidad. Nacen de matrimonios que han entendido ---a veces sin saberlo formular--- que el amor no es un tesoro para guardar en una caja fuerte, sino un fuego que solo se mantiene vivo si se comparte. Que cuanto más das, más tienes. Que la generosidad no empobrece: multiplica.

En el capítulo primero de este libro hablamos de algo que tal vez recordéis: la supervivencia de la especie humana depende de la cooperación.\footnote{Cf. capítulo 1 de este libro.} Nuestros bebés nacen tan inmaduros, tan vulnerables, que sin una red de cuidado no sobreviven. Y esa red ---madre, padre, abuelos, comunidad--- es la primera forma de familia. La biología nos dice que no estamos hechos para criar solos, para vivir solos, para amar solos.

Pues bien: la fecundidad del matrimonio es la contribución de esa familia a la red. Es el hilo que los esposos tejen en la trama grande de la comunidad humana. Cada familia fecunda ---fecunda de verdad, no solo en número--- fortalece el tejido entero. Cada matrimonio que irradia vida hace que el mundo sea un poco más habitable para todos.

\section{La generosidad que piensa}

Y aquí hay que hablar de algo que a veces se omite por miedo a decir lo incorrecto: la paternidad responsable.

Porque si la fecundidad no es solo cuestión de números, tampoco la generosidad es cuestión de imprudencia. La Iglesia no pide a los matrimonios que tengan «todos los hijos que vengan» sin más, como quien deja un grifo abierto y se desentiende. Lo que pide es algo más exigente y más libre a la vez: que sean generosos \emph{y} prudentes. Que acojan a los hijos como un don, sí, pero que lo hagan con responsabilidad, discerniendo juntos ---los dos, en diálogo, ante Dios--- cuántos hijos pueden amar bien.\footnote{«Por razones justificadas, los esposos pueden querer espaciar los nacimientos de sus hijos. [\ldots] La paternidad responsable comporta [\ldots] el dominio que la razón y la voluntad deben ejercer.» Pablo VI, Encíclica \textit{Humanae Vitae}, n. 10 y 16 (1968). Cf. también \textit{Gaudium et Spes}, n. 50: «Los esposos [\ldots] de común acuerdo y esfuerzo, se formarán un juicio recto, atendiendo tanto a su propio bien como al de los hijos, ya nacidos o todavía por venir.»}

No es lo mismo tener tres hijos con el corazón abierto, bien educados, bien queridos, que tener ocho por inercia, saturados de cansancio y sin espacio emocional para ninguno. Y no es lo mismo tener un hijo único al que se ama con generosidad y al que se enseña a compartir, que no tener ninguno por puro egoísmo, porque «los niños molestan» o porque «hay que vivir la vida».

La clave, como siempre en el matrimonio cristiano, no está en el número. Está en el corazón. ¿Está abierto o cerrado? ¿Da vida o la retiene? ¿Es generoso o calculador?

Un matrimonio que tiene dos hijos y vive abierto al mundo es más fecundo que uno que tiene diez y vive replegado sobre sí mismo. Un matrimonio sin hijos que adopta, que acoge, que derrama vida a su alrededor, puede ser extraordinariamente fecundo. La fecundidad no es un dato demográfico. Es una actitud del alma. Eso sí, que la fecundidad de la vida sea amplia no borra el significado propio de cada gesto de amor entre los esposos; son dos niveles distintos, y en el próximo capítulo volveremos sobre ello.

Y esto merece un momento de atención especial, porque a veces se pasa de largo demasiado rápido. Un matrimonio que no tiene hijos ---por las razones que sean, que pertenecen al misterio de cada vida--- no es un matrimonio con la fecundidad cerrada. Es, si cabe, un matrimonio con la fecundidad \emph{especialmente} abierta. Porque toda la capacidad de entrega, toda la energía de amor que esos esposos llevan dentro, no se queda sin cauce: busca otros caminos. Y los encuentra. La acogida del que sufre. El servicio al que necesita. La entrega generosa al amigo, al vecino, al desconocido. La irradiación callada de una vida compartida que, precisamente porque no se cierra sobre los hijos propios, se abre de par en par al mundo.

No tener hijos biológicos no clausura la fecundidad. Puede, paradójicamente, lanzarla de una manera que pocos imaginan. Hay matrimonios sin hijos que han sido padres y madres de medio barrio, que han sostenido comunidades enteras, que han dado más vida a su alrededor que muchas familias numerosas. Porque la fecundidad, como venimos diciendo, no se mide por la biología. Se mide por la vida que brota alrededor.

\section{Sal y luz}

Jesús usó dos imágenes muy concretas para hablar de lo que sus discípulos estaban llamados a ser: sal y luz.\footnote{Mt 5,13-14.}

La sal no se ve. Nadie dice, al probar un plato bien sazonado: «Qué buena está la sal.» Dice: «Qué bueno está el plato.» La sal desaparece para que el sabor aparezca. No se impone, no grita, no llama la atención. Simplemente \emph{está}, y su presencia cambia todo lo que toca.

La luz tampoco pide permiso. No necesita explicar lo que hace. Simplemente ilumina. Entra en la habitación oscura y la oscuridad retrocede, no porque la luz luche contra ella, sino porque la oscuridad no puede sostenerse ante la luz.

Un matrimonio fecundo es sal y luz. No necesita dar lecciones sobre el amor. No necesita demostrar nada. Solo necesita vivir lo que es: una alianza de amor que genera vida. Y cuando lo hace, la gente lo nota. Los amigos lo notan. Los vecinos lo notan. Los hijos ---propios y ajenos--- lo notan. Y algo se enciende en ellos: el deseo de vivir así. La intuición de que eso es posible. La esperanza de que el amor no es una mentira.

Esa es la fecundidad más profunda del matrimonio. No la que se ve en el certificado de familia numerosa. La que se ve en los ojos de quienes han sido tocados por esa familia. La que germina en silencio, como la semilla en la tierra, y da fruto treinta, sesenta, cien por uno.\footnote{Cf. Mt 13,8.}

\section{Lo que el mundo necesita ver}

Vivimos en un tiempo raro. Un tiempo en el que el amor se consume rápido y se descarta fácil. En el que las relaciones se miden por lo que me aportan y se abandonan cuando dejan de ser rentables. En el que la palabra «para siempre» se ha vuelto sospechosa, como una promesa que nadie se atreve a tomarse en serio.

En ese contexto, un matrimonio que dura, que crece, que se abre a los demás, que genera vida en todas sus formas, es un acto casi revolucionario. No porque sea perfecto ---ya lo hemos dicho: la perfección no existe en la vida real---. Sino porque es \emph{posible}. Porque demuestra, con hechos más que con palabras, que el amor fiel, indisoluble y fecundo no es una reliquia del pasado ni una imposición de la Iglesia: es lo que el corazón humano desea en lo más hondo.

Cada matrimonio que vive así ---con sus luchas, con sus caídas, con sus reconciliaciones--- está diciendo al mundo algo que el mundo necesita oír desesperadamente: que vale la pena apostar por el amor. Que vale la pena quedarse. Que vale la pena dar vida.

No hace falta ser héroes. No hace falta ser santos de estampa. Basta con ser sal. Basta con ser luz. Basta con vivir el amor con el corazón abierto y dejar que ese amor haga su trabajo.

Porque un matrimonio fecundo no es el que llena la casa de hijos. Es el que llena el mundo de esperanza.

\paracaminar

\begin{enumerate}
  \item Pensad en una familia que os haya marcado positivamente. ¿Qué tenía esa familia que la hacía especial? ¿Podéis identificar formas de fecundidad que iban más allá de los hijos?

  \item ¿Cómo imagináis vuestra familia dentro de diez años? No solo en número de hijos, sino en estilo de vida: ¿será vuestra casa un lugar abierto o cerrado? ¿Qué os gustaría que la gente sintiera al cruzar vuestra puerta?

  \item La imagen de la sal y la luz: ¿en qué ámbitos concretos ---amigos, barrio, parroquia, trabajo--- creéis que vuestro matrimonio puede «dar sabor» y «dar luz»?
\end{enumerate}

\vinetacapitulo{ch10}

\chapter{La sexualidad como don}

\epigraph{El cuerpo, y solo él, es capaz de hacer visible lo que es invisible: lo espiritual y lo divino.}{Juan Pablo II}

Hay pocos temas sobre los que se hable tanto y se diga tan poco.

La sexualidad humana está en todas partes ---en la publicidad, en las series, en las conversaciones de bar, en los titulares--- y, sin embargo, rara vez se habla de ella con la profundidad que merece. Nos movemos entre dos extremos que parecen opuestos pero que comparten una misma raíz: no entender qué es. De un lado, la represión: la sexualidad como algo sucio, peligroso, de lo que es mejor no hablar. Del otro, la banalización: la sexualidad como consumo, como entretenimiento, como un acto más entre otros que no compromete a nada ni a nadie.

Entre el tabú y el escaparate hay un abismo. Y en ese abismo se ha perdido algo esencial. Algo que nuestros cuerpos saben pero que la cultura ha olvidado. Algo que la fe cristiana custodia, no como prohibición, sino como buena noticia.

Vamos a intentar recuperarlo.

\section{Ni represión ni consumo}

Empecemos por desmontar los dos extremos, porque ambos hacen daño.

La represión tiene una historia larga. Durante siglos, una cierta lectura del cristianismo ---más influida por el neoplatonismo y el jansenismo que por el Evangelio--- presentó la sexualidad como una concesión a la debilidad de la carne, algo tolerado en el matrimonio solo porque servía para tener hijos, pero sospechoso en sí mismo. El placer era visto con desconfianza. El cuerpo, como cárcel del alma. El deseo, como enemigo.\footnote{Esta visión negativa del cuerpo y la sexualidad debe más al dualismo platónico y a corrientes rigoristas que al mensaje bíblico. El Cantar de los Cantares, incluido en el canon de las Escrituras, es un canto apasionado al amor erótico entre los esposos. La tradición de la Iglesia siempre ha afirmado la bondad de la creación ---incluido el cuerpo y la sexualidad---, aunque no siempre su predicación haya estado a la altura de esa afirmación.}

El resultado fue generaciones enteras para las que hablar de sexualidad equivalía a hablar de pecado. Y muchos que crecieron así ---quizá vuestros padres, quizá vosotros mismos--- arrastran heridas reales: vergüenza, culpa, incapacidad de vivir la intimidad con libertad y alegría.

En el otro extremo, la cultura contemporánea ha hecho algo distinto pero igualmente empobrecedor: ha vaciado la sexualidad de significado. Si antes el cuerpo era sospechoso, ahora es irrelevante. El sexo se presenta como una actividad recreativa, un intercambio de sensaciones entre cuerpos que se usan mutuamente, sin que el acto diga nada ni comprometa a nada. «Es solo sexo», se dice, como si esas tres palabras no fueran una contradicción en sí mismas.

Porque el sexo nunca es «solo» sexo. Cualquiera que haya vivido la experiencia lo sabe. La intimidad sexual nos deja vulnerables de una forma que ninguna otra actividad humana consigue. Nos expone. Nos desnuda ---y no hablo solo de la ropa---. Hay algo en el acto sexual que toca una fibra profunda, algo que no se puede reducir a biología ni a deporte.

Los dos extremos, curiosamente, comparten el mismo error: tratar la sexualidad como si fuera \emph{solo} un asunto del cuerpo. Para la represión, el cuerpo es el problema. Para la banalización, el cuerpo es el instrumento. Ninguno de los dos entiende que el cuerpo es el \emph{lenguaje}.

\section{Lo que el cuerpo dice}

En los capítulos anteriores hemos visto que el cuerpo humano no es un accidente biológico ni una máquina indiferente. El cuerpo habla. Tiene su propia gramática, su propia elocuencia. La oxitocina que se dispara con el contacto piel con piel, las hormonas que empujan al vínculo, el rostro que se sonroja, las manos que tiemblan cuando rozan por primera vez las de la persona amada: todo eso es lenguaje. El cuerpo dice cosas que la boca a veces no sabe decir.

Pues bien: si hay un momento en que el cuerpo habla con toda su elocuencia, ese momento es la unión sexual.

Juan Pablo II dedicó años a pensar sobre esto. Entre 1979 y 1984, en una serie de audiencias generales que pasaron relativamente desapercibidas, desarrolló lo que se conoce como la \emph{Teología del Cuerpo}: una reflexión sobre el significado del cuerpo humano y la sexualidad que muchos consideran una de las aportaciones más originales de la teología del siglo XX.\footnote{Juan Pablo II, \textit{Hombre y mujer los creó. Catequesis sobre el amor humano}, 129 audiencias generales entre septiembre de 1979 y noviembre de 1984. La obra completa ha sido editada en múltiples volúmenes. Para una introducción accesible en español, cf. Yves Semen, \textit{La sexualidad según Juan Pablo II}, Desclée de Brouwer, 2005.} No es un texto fácil ---su estilo es denso, filosófico, a ratos casi impenetrable---, pero la idea central es de una sencillez luminosa.

La idea es esta: el cuerpo humano tiene un «significado nupcial». Es decir, está hecho para la entrega. No solo puede entregarse: está \emph{diseñado} para ello. Toda la configuración del cuerpo humano ---la diferencia sexual, la complementariedad, el deseo que nos empuja hacia el otro--- apunta a un mismo destino: la donación de uno mismo a la persona amada.

Y esa donación alcanza su expresión más plena, más radical, más elocuente, en la unión sexual de los esposos.

¿Qué dice el cuerpo cuando dos personas se unen en el acto sexual? Dice algo enorme. Dice algo que, si nos paramos a escucharlo, resulta casi sobrecogedor.

El cuerpo dice: \emph{Te doy todo lo que soy.}

No una parte. No un rato. No con condiciones. \emph{Todo.} La desnudez física es signo de una desnudez más profunda: me muestro ante ti tal como soy, sin máscaras, sin defensas, sin reservas. Me entrego entero. Te acojo entera.

Por eso la sexualidad no puede ser trivial. Por eso «solo sexo» es una frase imposible. Porque el cuerpo no sabe mentir ---o, mejor dicho, cuando miente, algo se rompe por dentro---. Cuando dos cuerpos se unen diciendo «te doy todo» mientras la voluntad piensa «te doy un rato», hay una fractura. Una incoherencia entre lo que el cuerpo proclama y lo que la persona decide. Y esa fractura, aunque no siempre sea consciente, deja huella.

\section{Las tres palabras del cuerpo}

Aquí es donde todo lo que hemos visto en los capítulos anteriores se hace carne. Literalmente.

Hemos hablado de las tres propiedades del matrimonio cristiano: indisolubilidad, fidelidad, fecundidad.\footnote{Cf. capítulos 8, 9 y 10 de este libro.} Las hemos visto como regalos, como expresiones del amor total. Pues bien: en la unión sexual, esas tres propiedades dejan de ser conceptos teológicos y se convierten en lenguaje del cuerpo. Lo que los esposos se prometieron con palabras el día de la boda, sus cuerpos lo repiten cada vez que se unen.

\subsection*{El cuerpo dice: «Para siempre»}

Cuando dos esposos se unen, sus cuerpos dicen indisolubilidad. La entrega sexual es, por su propia naturaleza, una entrega sin fecha de caducidad. No es un préstamo. No es un alquiler. Es un don total, sin reservas, sin letra pequeña que diga «hasta que me canse» o «mientras me convenga».

«Serán una sola carne»,\footnote{Gn 2,24.} dice el Génesis. No «una sola carne de momento». No «una sola carne a tiempo parcial». Una sola carne, sin más. La unión de los cuerpos dice para siempre con una elocuencia que las palabras no alcanzan. Cada vez que los esposos se unen, renuevan con sus cuerpos la promesa que hicieron con sus labios: \emph{No me voy. Me quedo contigo. Pase lo que pase.}

\subsection*{El cuerpo dice: «Solo a ti»}

La entrega sexual dice también fidelidad. Y no la fidelidad entendida como restricción ---«no puedo estar con nadie más»--- sino como concentración de toda la capacidad de donación en una sola persona: «todo lo que soy, para ti; todo lo que tengo, para ti; toda mi capacidad de amar, volcada en ti».

Es exclusiva no por prohibición, sino por lógica interna. Si me doy \emph{entero} a ti, no me queda nada para dar a otro. No porque alguien me lo prohíba, sino porque así funciona la entrega total: es indivisible. No se puede dar todo a dos personas a la vez, igual que no se puede estar entero en dos sitios al mismo tiempo.

Los cuerpos que se unen dicen \textit{fides}, fe: «Creo en ti. Me fío de ti. Deposito en tus manos lo más íntimo que tengo ---mi cuerpo, mi vulnerabilidad, mi desnudez--- porque confío en que lo vas a custodiar.»\footnote{Cf. capítulo 9 de este libro, donde se desarrolla la fidelidad como fe depositada en el otro.}

\subsection*{El cuerpo dice: «Estamos abiertos a lo que venga»}

Y los cuerpos que se unen dicen fecundidad. Pero aquí conviene hacer una pausa, porque en el capítulo anterior hemos visto que la fecundidad del matrimonio es algo mucho más amplio que tener hijos. Un matrimonio fecundo es el que irradia vida a su alrededor, el que es sal de la tierra, el que genera esa envidia sana que hace que otros deseen amar así. Los hijos, cuando los hay, son parte de esa fecundidad ---una parte hermosa y central---, pero no la agotan.

Lo que la unión sexual expresa es una dimensión específica de esa fecundidad amplia: la apertura a la vida biológica. Los esposos no se unen «a pesar de» la posibilidad de engendrar, ni «para» engendrar como si fuera una producción industrial. Se unen como don total ---y un don total incluye la fertilidad, incluye la posibilidad de que ese amor sea tan real, tan concreto, que tome carne en una persona nueva---.

La Teología del Cuerpo de Juan Pablo II desarrolló esta idea con una profundidad que no tenía precedentes.\footnote{Cf. Juan Pablo II, \textit{Hombre y mujer los creó}, audiencias 117-133; Pablo VI, Encíclica \textit{Humanae Vitae} (1968), especialmente nn. 11-12, donde se desarrolla la inseparabilidad del significado unitivo y procreativo del acto conyugal. También Juan Pablo II, \textit{Familiaris Consortio} (1981), nn. 28-35.} Para él, si la entrega sexual es total, si los cuerpos dicen «te doy todo lo que soy», entonces cerrar deliberadamente la puerta a la procreación es introducir una reserva en el lenguaje del cuerpo. Es decir «te doy todo» mientras se retiene algo.

Conviene no confundir dos planos distintos. Como vimos en el capítulo anterior, la fecundidad global de un matrimonio es inmensamente amplia: abarca lo social, lo espiritual, el testimonio, la acogida. Un matrimonio sin hijos puede ser extraordinariamente fecundo. Pero aquí no hablamos de la fecundidad de la vida entera, sino del significado de un gesto concreto: el acto de unión conyugal. Y ese gesto tiene su propia gramática. No se trata de que cada encuentro deba producir un hijo ---eso nadie lo sostiene---, sino de que el lenguaje que hablan los cuerpos al unirse sea el de una entrega sin letra pequeña, donde la apertura a la vida es una dimensión de la donación total, no una obligación de resultado.

Es una enseñanza exigente, y conviene recibirla con honestidad: es la visión de san Juan Pablo II, coherente con toda su antropología del cuerpo. Pero probablemente la fecundidad a la que está llamado el matrimonio se pueda ---y se deba--- ver como un todo más amplio que la procreación. La procreación no queda excluida de esa visión más amplia: queda \emph{incluida} en ella, como una de sus expresiones más bellas, pero no la única. Un matrimonio sin hijos no es un matrimonio estéril si su amor genera vida a su alrededor.\footnote{Cf. capítulo 10 de este libro, donde se desarrolla la fecundidad como irradiación. También la exhortación del Papa Francisco: «El amor siempre da vida. [...] El amor conyugal no se agota dentro de la pareja [...] sino que los habilita para la máxima donación.» Francisco, \textit{Amoris Laetitia}, n. 165 (2016).} Y un matrimonio con hijos, si vive replegado sobre sí mismo, puede estar traicionando una fecundidad que debería desbordarlo.

La apertura a la vida ---en su sentido más pleno--- no es una carga impuesta desde fuera. Es la consecuencia natural de un amor que se da sin reservas. Es lo que pasa cuando el «te doy todo» se dice de verdad, con el cuerpo y con la vida entera. Y esa vida que nace del amor toma muchas formas: a veces tiene la cara de un hijo, a veces la de un hogar acogedor, a veces la de un servicio callado al que lo necesita.

\section{La castidad como integración}

Si hay una palabra que la cultura contemporánea ha malentendido más que ninguna otra, esa palabra es \emph{castidad}.

Para la mayoría de la gente, castidad significa abstinencia. No tener relaciones sexuales. Punto. Y como la abstinencia no suele ser un proyecto de vida demasiado atractivo, la castidad ha quedado relegada al rincón de las virtudes polvorientas, de las que nadie habla excepto para burlarse.

Pero castidad no es abstinencia. La castidad es algo mucho más amplio, más dinámico, más vivo. La castidad es \emph{integración}.\footnote{«La castidad significa la integración lograda de la sexualidad en la persona, y por ello en la unidad interior del hombre en su ser corporal y espiritual.» \textit{Catecismo de la Iglesia Católica}, n. 2337.}

¿Qué significa eso? Significa que la sexualidad no es una fuerza salvaje que hay que encerrar en una jaula para que no haga daño, ni una energía neutra que se puede gastar como uno quiera sin consecuencias. La sexualidad es una dimensión profunda de la persona ---tan profunda como la inteligencia o la libertad--- que necesita ser integrada en el conjunto de lo que somos. Ordenada. No reprimida: \emph{ordenada}. Puesta al servicio del amor, que es para lo que existe.

Una persona casta no es una persona que no desea. Es una persona que desea \emph{bien}. Que ha aprendido a mirar al otro como persona y no como objeto. Que ha integrado su deseo sexual en su capacidad de amar, de modo que la atracción no le ciega sino que le ilumina, no le esclaviza sino que le libera para amar mejor.

La castidad es, para los esposos, la capacidad de vivir su sexualidad como lenguaje de donación y no de apropiación. De unirse al otro para \emph{darse}, no para \emph{tomarse}. De buscar en la intimidad el bien del otro antes que la propia satisfacción ---lo cual, paradójicamente, es el camino más seguro hacia una intimidad más plena y más gozosa---.

No es fácil. Nadie dijo que lo fuera. Pero es profundamente liberador, porque transforma el deseo de fuerza centrípeta ---todo hacia mí--- en fuerza centrífuga ---todo hacia ti---. Y eso es exactamente lo que hace el amor cuando es amor de verdad.

\section{La escuela de todos los días}

Y aquí quiero hablar de algo que va más allá de la sexualidad, pero que nace de ella. Porque la entrega que los esposos aprenden en la intimidad no se queda entre las sábanas. Se extiende, como una marea, a toda la vida compartida.

El matrimonio transforma. No deja a nadie igual.

Pensadlo un momento: no sois la misma persona que erais cuando os casasteis ---o cuando decidisteis casaros---. Y no me refiero solo a las canas o a los kilos. Me refiero a algo más profundo. La convivencia diaria con otra persona, la entrega concreta y cotidiana, la fricción inevitable de dos vidas que se funden sin dejar de ser distintas\ldots{} todo eso os ha ido cambiando. Os ha ido puliendo. Os ha ido haciendo ---si os habéis dejado--- mejores personas de las que habríais sido solos.

Porque el matrimonio es una escuela. No una escuela con pupitres y exámenes, sino una escuela de \emph{virtudes}, de esas que se aprenden en la piel, no en los libros.

Se aprende \textbf{humildad}, porque tarde o temprano ---más temprano que tarde--- hay que reconocer que te has equivocado. Que has dicho algo injusto. Que has sido egoísta. Y pedir perdón. Y el perdón pedido a la persona que duerme a tu lado es más costoso y más verdadero que cualquier examen de conciencia abstracto.

Se aprende \textbf{paciencia}, porque el otro no es como tú querrías que fuera. Tiene sus manías, sus ritmos, sus silencios que no entiendes y sus palabras que sobran. Y aprender a convivir con todo eso sin pretender cambiar al otro a tu imagen y semejanza es un ejercicio de paciencia que ningún libro puede enseñar.

Se aprende \textbf{justicia}, porque hay que ceder. Hay que repartir. Hay que escuchar la perspectiva del otro y reconocer que quizá ---solo quizá--- tiene razón. La vida compartida exige un sentido de la justicia que no se aprende en la facultad de Derecho sino en la cocina, en la logística de quién lleva a los niños al colegio, en la negociación silenciosa de cada día.

Se aprende \textbf{generosidad}, porque amar es dar sin llevar cuentas. Dar tiempo, dar atención, dar presencia, dar renuncia. Dar sin esperar retorno, sin pasar factura, sin el cálculo mezquino de «yo hice esto, ahora te toca a ti».

Y se aprende \textbf{fortaleza}, porque hay momentos en que todo es difícil. En que el cansancio pesa más que el cariño. En que la rutina ha borrado el brillo. En que uno querría salir corriendo. Y decidir quedarse ---no por inercia, no por cobardía, sino por amor--- es un acto de fortaleza que deja huella en el alma.

No es que el matrimonio enseñe estas virtudes en abstracto, como quien lee un catálogo de buenas intenciones. Es que la convivencia diaria, el roce constante con otra persona, te las va tallando en la carne. Te pule, te lima, te moldea. Te hace mejor persona ---si te dejas---. Y ese «si te dejas» es importante, porque la escuela del matrimonio solo funciona cuando uno acepta ser alumno. Cuando uno reconoce que no ha llegado terminado, que le queda mucho por aprender, que el otro es maestro tanto como compañero.

El matrimonio no solo une a dos personas. Las transforma. Y la persona que sale de esa transformación no es la que entró, ni tampoco la que habría sido si hubiera caminado sola.

\section{El sacramento hecho carne}

Volvamos a la sexualidad para cerrar el círculo.

Decíamos al principio que el cuerpo habla, y que en la unión sexual habla con su máxima elocuencia. Ahora podemos ver el cuadro completo: la sexualidad de los esposos es el lugar donde el sacramento del matrimonio se hace carne de la forma más literal posible.

Cuando los esposos se casaron, dijeron tres «sí» ante el altar: sí a la indisolubilidad, sí a la fidelidad, sí a la fecundidad.\footnote{El rito del matrimonio incluye tres preguntas a los contrayentes antes del consentimiento: si vienen libremente, si se comprometen a amarse y respetarse durante toda la vida, y si están dispuestos a recibir con amor los hijos que Dios les dé. Cf. \textit{Ritual del Matrimonio}, Praenotanda, nn. 2-3.} Esos tres «sí» no fueron solo palabras. Fueron el compromiso de una vida entera.

Y cada vez que se unen, sus cuerpos repiten esos tres «sí»: me doy a ti para siempre, me doy a ti en exclusiva, me doy a ti abierto a la vida que pueda venir. La unión sexual de los esposos no es un apéndice del sacramento, un añadido más o menos relevante. Es el sacramento \emph{en acto}. Es la alianza renovada en la carne. Es el «sí» del altar traducido al lenguaje del cuerpo.

Por eso la sexualidad conyugal es sagrada. No sagrada en el sentido de intocable o de solemne. Sagrada en el sentido más profundo: es un lugar de encuentro con el misterio. Un lugar donde dos personas, en su vulnerabilidad y su ternura, hacen visible algo del amor de Dios. Ese Dios que se entrega sin reservas, que es fiel hasta el final, y cuyo amor es tan fecundo que crea vida.

No hacen falta velas ni música de fondo para que esto sea verdad. Es verdad en la intimidad de un martes por la noche, con el cansancio de la jornada encima y los niños por fin dormidos. Es verdad en la ternura torpe de las primeras veces y en la complicidad callada de los muchos años. Es verdad siempre que dos esposos se unen con la intención de darse, no de tomarse; de amar, no de usar.

Y esa verdad ---que el cuerpo habla, que la sexualidad es don, que la intimidad conyugal dice cosas que las palabras no alcanzan--- es quizá la mejor noticia que la fe cristiana tiene que ofrecer en un mundo que ha olvidado escuchar lo que el cuerpo dice.

\paracaminar

\begin{enumerate}
  \item ¿Os reconocéis en alguno de los dos extremos ---represión o banalización--- a la hora de entender la sexualidad? ¿Cómo ha influido vuestra educación, vuestra familia o vuestra cultura en la forma en que vivís la intimidad?

  \item Pensad en las virtudes que el matrimonio os ha ido enseñando ---o que intuís que os enseñará---. ¿Cuál os cuesta más? ¿En cuál habéis crecido más?

  \item La idea de que el cuerpo «habla» en la unión sexual, y que dice indisolubilidad, fidelidad y apertura a la vida: ¿cómo la recibís? ¿Cambia en algo vuestra forma de entender la intimidad conyugal?
\end{enumerate}

\vinetacapitulo{ch11}

\chapter{La familia, comunidad de amor}

\epigraph{«Donde hay dos o tres reunidos en mi nombre, allí estoy yo en medio de ellos.»}{Mateo 18,20}

Este libro empezó con un bebé.

Un bebé que nacía con la cabeza sin cerrar, incapaz de sostenerla por sí mismo, incapaz de sobrevivir una sola noche sin que alguien le abrazara, le alimentara, le mantuviera con vida. Dijimos entonces que esa fragilidad no era un fallo de diseño, sino el diseño mismo. Que nacemos inacabados porque estamos hechos para ser completados por otros. Que el amor no es un adorno de la vida humana, sino su condición de posibilidad.

Ahora, al final del camino, quiero que volvamos a ese bebé. Pero esta vez vamos a ampliar el encuadre. Vamos a ver la foto completa.

Porque ese bebé no estaba solo con su madre. Nunca lo estuvo.

\section{La tribu que sostiene al niño}

Hay un proverbio africano que se ha hecho famoso ---tanto que ya nadie sabe a qué pueblo pertenece exactamente---: «Hace falta una aldea entera para criar a un niño.» Como muchos proverbios, dice en una línea lo que la ciencia ha necesitado décadas para demostrar.

En el capítulo 1 hablamos del dilema obstétrico: la pelvis se estrecha, el cerebro crece, el bebé nace inmaduro. La madre no puede sola. Necesita ayuda. Y dijimos que el padre que se queda, el que cuida, el que baja su testosterona y sube su oxitocina, fue una pieza clave en la supervivencia de nuestra especie.

Pero la historia no termina ahí. Hay otra pieza del rompecabezas que entonces solo mencionamos de pasada y que ahora merece toda nuestra atención.

Se llama la \textbf{hipótesis de la abuela}.\footnote{La hipótesis fue propuesta formalmente por la antropóloga Kristen Hawkes, de la Universidad de Utah, a partir de sus estudios con los Hadza de Tanzania en los años noventa. Cf. Hawkes et al., ``Grandmothering, menopause, and the evolution of human life histories'', \textit{PNAS}, 1998.}

El enigma es el siguiente: la menopausia humana es una rareza biológica. En casi todas las especies, las hembras son fértiles hasta el final de su vida. ¿Por qué las mujeres humanas dejan de ser fértiles décadas antes de morir? Desde un punto de vista puramente evolutivo, ¿no sería más «eficiente» seguir reproduciéndose?

La respuesta que propone esta hipótesis es fascinante: las abuelas que dejaban de tener hijos propios y se dedicaban a ayudar a criar a los hijos de sus hijas contribuían más a la supervivencia de sus genes que si hubieran seguido teniendo crías. Una abuela que recolecta alimentos, que vigila a los niños, que cuida a la madre agotada, que enseña a la nieta a identificar raíces comestibles, que consuela al nieto asustado, multiplica las posibilidades de que esos nietos sobrevivan y se reproduzcan.

La evolución, en otras palabras, no favoreció solo a la madre que cuidaba. Ni solo a la pareja que permanecía unida. Favoreció a la \emph{familia extensa}. A la red de cuidado intergeneracional. A la tribu.

\section{Las abuelas de Dios}

Hay algo que me conmueve profundamente de esta hipótesis, y es que la ciencia, sin pretenderlo, está describiendo algo que la fe ya conocía.

Porque lo que la hipótesis de la abuela dice, traducido al lenguaje de la fe, es que el ser humano está diseñado para que el amor desborde. No basta el amor de la madre. No basta el amor de la pareja. Hace falta un amor que se extienda más allá, que cruce generaciones, que cree redes donde cada persona cuida de otra y es cuidada a su vez.

La abuela que deja de engendrar para dedicarse a cuidar está haciendo algo que se parece mucho a lo que la tradición cristiana llama \emph{fecundidad espiritual}: una vida que da fruto no a través de la reproducción, sino a través de la entrega. Una vida que se multiplica cuidando.

Pensad en vuestras propias abuelas. O en esas figuras que, sin ser abuelas de sangre, ejercieron ese papel: la vecina que os recogía del colegio, la tía que os enseñó a cocinar, el abuelo que os contaba historias que todavía recordáis. Cada una de esas personas fue, a su manera, parte de la aldea que os crió. Cada una de ellas hizo posible que llegarais hasta aquí.

La familia no es la pareja encerrada en su piso. La familia es una red. Una red que se teje con hilos de sangre, pero también de cercanía, de costumbre, de cariño. Una red que, cuando funciona, sostiene a los débiles, acoge a los nuevos, acompaña a los que se van. Y cuando no funciona ---porque a veces no funciona---, su ausencia se nota como se nota el frío cuando falta el techo.

\section{La Iglesia doméstica}

El Concilio Vaticano II, en uno de esos hallazgos de lenguaje que parecen sencillos pero encierran siglos de reflexión, llamó a la familia «Iglesia doméstica».\footnote{Concilio Vaticano II, Constitución dogmática \textit{Lumen Gentium}, n. 11 (1964): «En esta especie de Iglesia doméstica, los padres deben ser para sus hijos los primeros predicadores de la fe, mediante la palabra y el ejemplo.» Juan Pablo II retomó y desarrolló ampliamente esta idea en la Exhortación apostólica \textit{Familiaris Consortio} (1981), especialmente en los nn. 49-62.}

\textit{Iglesia doméstica.} La expresión es más fuerte de lo que parece. No dice que la familia sea «como» una iglesia, en sentido metafórico. Dice que la familia \emph{es} Iglesia. La versión más pequeña, más cotidiana, más real de la Iglesia. El lugar donde la fe no se explica, sino que se vive. Donde no se predica con palabras, sino con gestos.

Porque, pensémoslo bien: ¿dónde aprende un niño lo que es el perdón? No en una clase de catequesis. Lo aprende cuando ve a su padre pedirle perdón a su madre después de una discusión. ¿Dónde descubre lo que es la generosidad? Cuando ve a sus padres acoger en casa a alguien que lo necesita. ¿Dónde intuye, por primera vez, que existe algo más grande que él, algo que sostiene la vida y le da sentido? Quizá en esa oración sencilla antes de cenar que nadie le obliga a decir, pero que todo el mundo dice.

La familia es el primer lugar donde se experimenta ---o no se experimenta--- que el amor es real. Y esa experiencia marca para siempre. No porque determine fatalmente el futuro, sino porque establece el suelo sobre el que todo lo demás se construye. Un niño que ha visto amor sabe que el amor existe, aunque luego le cueste encontrarlo. Un niño que no lo ha visto tendrá que descubrirlo por sí mismo, que es más difícil ---pero no imposible, porque la gracia llega donde la naturaleza no alcanza---.

\section{Los suegros y otras pruebas de fe}

Y ahora viene el momento de hablar de algo que todo el que se ha casado conoce, que todos los manuales de pastoral mencionan con delicadeza y que la vida real sirve sin anestesia: las familias de origen.

Porque cuando dos personas se casan, no se casan solas. Se casan con sus historias, sus costumbres, sus formas de celebrar la Navidad y de doblar las toallas. Se casan con sus padres, sus hermanos, sus tradiciones familiares. Y todo eso tiene que encontrar un encaje que, a veces, rechina.

El Génesis lo dice con una claridad que impresiona: «Dejará el hombre a su padre y a su madre y se unirá a su mujer.»\footnote{Gn 2,24.} Dejar. No abandonar, no renegar, no olvidar. Dejar. Soltar el lugar de hijo para ocupar el lugar de esposo. Soltar el lugar de hija para ocupar el lugar de esposa. Crear un espacio nuevo, una casa nueva ---aunque sea el mismo piso de alquiler---, un «nosotros» que es distinto del «nosotros» de la familia en la que creciste.

Y eso, que suena muy bonito en un libro, en la práctica significa aprender a navegar entre dos familias que os quieren, que quieren lo mejor para vosotros, que a veces no se ponen de acuerdo sobre qué es lo mejor para vosotros, y que ---seamos honestos--- a veces meten la nariz donde no las llaman.

No voy a dar recetas, porque cada familia es un mundo y cada suegra es un universo.\footnote{Y cada suegro. Los suegros también existen, aunque la cultura popular haya concentrado todo su humor en las suegras. Quede aquí constancia de que el desafío es ecuménico.} Pero sí voy a decir dos cosas que he visto confirmadas una y otra vez.

La primera: \textbf{honrar a vuestros padres no significa obedecerlos en todo}. Sois adultos. Habéis fundado una familia nueva. Las decisiones sobre vuestra casa, vuestra economía, vuestra forma de educar a vuestros hijos, son \emph{vuestras}. Pedir consejo es sabio. Dejarse gobernar es un error. Y cuando haya que poner un límite, es mejor que lo ponga el hijo o la hija de sangre, no el cónyuge. Si tu madre se mete demasiado, eres tú quien tiene que hablar con ella. No tu marido. Eso no es cobardía: es cuidar a las dos personas que quieres.

La segunda: \textbf{la familia política es un regalo que no habéis elegido, pero que podéis aprender a querer}. No os casasteis con ellos, pero forman parte de la vida de la persona con la que sí os casasteis. Y esa persona es, en parte, lo que es gracias a ellos. Así que miradlos con agradecimiento, aunque a veces os cueste. Y tened paciencia. Mucha paciencia. La misma que esperáis que ellos tengan con vosotros.

\section{La puerta abierta}

Si hay algo que define a una familia cristiana, no es que rece mucho ni que vaya a misa todos los domingos ---aunque ambas cosas estén muy bien---. Lo que define a una familia cristiana es que su puerta está abierta.

Abierta al hijo que llega, esperado o por sorpresa. Abierta al amigo que necesita un plato de comida y una conversación. Abierta al vecino que está solo. Abierta al sobrino que necesita un lugar donde quedarse mientras encuentra su camino. Abierta al mundo.

Hablamos en el capítulo 10 de la fecundidad como irradiación: un matrimonio fecundo es el que genera vida a su alrededor.\footnote{Cf. capítulo 10 de este libro.} Pues la familia es el lugar donde esa fecundidad se hace concreta. Una familia abierta es una familia que no vive para sí misma. Que no se agota en sus propias necesidades. Que tiene la generosidad de mirar más allá de sus paredes y preguntarse: ¿a quién podemos acoger? ¿A quién podemos servir?

No hace falta ser héroes. A veces basta con invitar a cenar a esa pareja que acaba de mudarse al barrio y no conoce a nadie. Con acompañar al anciano del tercero al médico. Con preparar la habitación de los hijos ---que ya se fueron--- para quien la necesite. La hospitalidad no es un gran gesto: es una disposición del corazón que se traduce en cosas muy pequeñas.

San Juan Pablo II escribió que la familia está llamada a ser «la primera y vital célula de la sociedad».\footnote{Juan Pablo II, Exhortación apostólica \textit{Familiaris Consortio}, n. 42 (1981), citando el Decreto \textit{Apostolicam Actuositatem} del Concilio Vaticano II, n. 11.} No la única célula, pero sí la primera. El lugar donde se aprende a convivir, a compartir, a resolver conflictos sin destruirse. Si la familia funciona, la sociedad tiene una oportunidad. Si la familia se cierra sobre sí misma, se asfixia. Como la semilla que no rompe la tierra: tiene toda la vida dentro, pero si no se abre, no da fruto.

\section{El perdón de cada día}

Hay una frase que, si la familia fuera una empresa, sería su lema corporativo. No es «te quiero», aunque esa es importante. No es «gracias», aunque esa es imprescindible. La frase más necesaria en una familia, la que más se usa en las casas donde la gente es feliz de verdad, es:

\textbf{«Perdona.»}

No hablo del gran perdón, el perdón dramático, el perdón de las películas donde alguien ha hecho algo terrible y el otro, tras una noche de lluvia y violines, decide perdonar. Ese perdón existe, y a veces hace falta. Pero no es el que sostiene una familia.

El que sostiene una familia es el perdón pequeño. El de cada día. El «perdona, he sido brusco». El «perdona, no te estaba escuchando». El «perdona, me he pasado con lo que he dicho». El «perdona, tenías razón y yo no quería verlo». Ese perdón que no sale en las películas porque no es espectacular, pero que es el aceite que mantiene la maquinaria en marcha.

Vivir juntos es rozarse. Es inevitable. Compartir espacio, tiempo, dinero, decisiones, cansancio, hijos, preocupaciones ---todo eso genera fricción---. Y la fricción, si no se lubrica con perdón, acaba desgastando. No de golpe. Poco a poco. Como el agua que erosiona la piedra: no por su fuerza, sino por su constancia.

El perdón en la familia no es un acto heroico. Es una necesidad higiénica. Tan básica como ducharse. Y, como la ducha, hay que practicarlo todos los días, no una vez al año.\footnote{San Pablo lo dijo con menos humor pero más profundidad: «Que no se ponga el sol sobre vuestra ira» (Ef 4,26). En otras palabras: no dejéis que un enfado pernocte en vuestra casa. Resolvedlo antes de dormir. No siempre es posible, pero como ideal es imbatible.}

Lo contrario del perdón no es el rencor espectacular. Es la cuenta que se va acumulando en silencio. Es el «no digo nada, pero me acuerdo». Es la lista mental de agravios que un día, sin que nadie lo espere, estalla. Las familias no se rompen por una gran traición. Se rompen por mil pequeñas heridas que nadie curó.

Perdonar no es sentir que no ha pasado nada. Perdonar es decidir que lo que ha pasado no va a definir la relación. Es elegir al otro por encima del agravio. Y eso, en una familia, hay que hacerlo todos los días. A veces varias veces al día.

Es agotador. Es necesario. Es, probablemente, lo más parecido a amar como Dios ama que podemos hacer los seres humanos corrientes.

\section{La familia imperfecta}

Y aquí quiero decir algo que ojalá alguien me hubiera dicho a mí antes de casarme: \textbf{vuestra familia no va a ser perfecta. Y eso no es un fracaso. Es lo normal.}

Habrá días en los que todo fluya: la casa estará recogida, los niños serán adorables, la cena saldrá rica y os miraréis a los ojos y pensaréis: «Qué suerte tengo.» Pero habrá muchos más días en los que la casa será un caos, los niños se pelearán, la cena se quemará y os miraréis a los ojos y pensaréis: «¿Cómo hemos llegado hasta aquí?»

Ambos días son reales. Ambos son vuestra familia.

La tentación más peligrosa no es la infidelidad ni el desamor. La tentación más peligrosa es la comparación. Es mirar a esa otra familia ---la del amigo, la del vecino, la que sale en Instagram--- y pensar que ellos sí lo tienen todo resuelto. Que su casa siempre está limpia, sus hijos siempre obedecen, su matrimonio siempre brilla. Es mentira, por supuesto. Nadie publica las fotos de la discusión de las once de la noche ni del niño llorando en el supermercado. Pero la mentira es eficaz, y siembra la duda: «Quizá lo estamos haciendo mal.»

No. No lo estáis haciendo mal. Lo estáis haciendo como se hace: con lo que tenéis, con lo que sabéis, con lo que podéis. Con más buena voluntad que aciertos, y eso es suficiente. Porque la familia no es un escaparate. Es un taller. Un lugar donde se trabaja, donde se falla, donde se vuelve a intentar. Donde la materia prima es el amor imperfecto de personas imperfectas que, sin embargo, han decidido estar juntas.

El papa Francisco, con esa franqueza suya que a veces incomoda y siempre ilumina, lo dijo así: «No existe la familia perfecta. No hay maridos perfectos, ni esposas perfectas. No nos hablemos de suegras perfectas. [\ldots] Solo existen familias en camino.»\footnote{Papa Francisco, audiencia general del 16 de septiembre de 2015.} En camino. No en destino. La perfección no es el punto de partida ni el requisito. Es el horizonte hacia el que se camina, sabiendo que nunca se llega del todo, y que eso está bien.

\section{Diseñados para amar}

Permitidme que, al cerrar este capítulo ---y con él, este libro---, vuelva una vez más al principio.

Empezamos hablando de un bebé que nacía indefenso, con la cabeza sin soldar y los ojos que apenas veían. Dijimos que esa fragilidad era el origen de todo: de la necesidad de cuidar, de vincularse, de amar. La biología nos mostraba un cuerpo diseñado para la relación: hormonas que empujan al vínculo, cerebros que se transforman con la paternidad y la maternidad, una química que favorece la permanencia.

Después la fe nos dio la clave de lectura: no estamos hechos para amar por casualidad ni por presión evolutiva. Estamos hechos para amar porque quien nos hizo es Amor. Somos imagen de un Dios que es, en sí mismo, comunión, entrega, don. Y amando ---entregándonos, saliéndonos de nosotros mismos, haciéndonos don para el otro--- nos parecemos a Él.

Vimos que el matrimonio no es un contrato con fecha de caducidad, sino una alianza sin letra pequeña. Que la indisolubilidad no es una cadena, sino la promesa más grande que una persona puede recibir. Que la fidelidad es fe, la fe absoluta que el otro deposita en ti. Que la fecundidad desborda los hijos y se derrama en todo lo que toca una familia que vive abierta. Que la sexualidad es un lenguaje que dice con el cuerpo lo que la alianza dice con palabras: «Todo lo mío es tuyo, sin reservas, para siempre.»

Y ahora, al final, vemos el cuadro completo: aquel bebé indefenso crece. Es sostenido por una madre, por un padre, por unos abuelos, por una comunidad que le cuida, le enseña, le quiere. Poco a poco aprende a hablar, a caminar, a pensar. Aprende ---si tiene suerte--- lo que es ser querido. Y un día, ya adulto, mira a otra persona a los ojos y siente algo que reconoce sin que nadie se lo explique: el deseo de entregarse. La vocación al amor. Y dice: «Sí, quiero.»

Y empieza de nuevo. Un hogar. Una alianza. Una vida compartida. Y quizá, un día, otro bebé indefenso que necesitará una cabeza que le sostengan, unos brazos que le abracen, una aldea entera que le críe. El ciclo continúa. La misma historia, en otra clave, con otros rostros, con los mismos miedos y la misma esperanza.

\medskip

Dijimos al hablar del matrimonio que la Cruz no es principalmente sufrimiento: es entrega.\footnote{Cf. capítulo 7 de este libro.} Que la persona a la que amas es tu cruz ---la razón por la que entregas tu vida--- y que Dios nos da una cruz a la que podemos abrazar cada noche. Pues bien: la familia es la extensión de ese abrazo. Es la Cruz que se hace mesa, que se hace casa, que se hace sobremesa un domingo, que se hace discusión un lunes y reconciliación un martes. Es la entrega hecha cotidianidad.

No será perfecta. Será real. Habrá días luminosos y noches oscuras. Habrá risas y portazos, cenas memorables y leche derramada, oraciones a medias y silencios que pesan. Habrá momentos en los que pensaréis que no podéis más, y momentos en los que sabréis, con una certeza que no necesita argumentos, que esto ---exactamente esto, con todo su desorden--- es lo más importante que habéis hecho en la vida.

Porque estáis diseñados para esto. No para la perfección, sino para el amor. No para acertar siempre, sino para volver a intentarlo cada mañana. No para ser una familia de postal, sino para ser una familia de verdad: torpe, ruidosa, desordenada, y llena ---llena hasta los bordes--- de algo que se parece mucho a Dios.

\medskip

Y quiero dejaros con una última imagen. Pensad en el camino de una vida entera. Desde aquel bebé con la cabeza sin cerrar hasta el último día. Es un camino largo, lleno de etapas, de giros, de subidas y bajadas. Pues bien: vuestro cónyuge es la persona que camina con vosotros ese camino. No un tramo. No una etapa. El camino entero. Hasta la misma puerta del cielo.

Porque si lo que dijimos en el capítulo 7 es verdad ---que tu esposo, tu esposa, es tu cruz, y que la Cruz es el camino hacia el Padre---, entonces tu cónyuge es, literalmente, quien te lleva al cielo. No en sentido figurado. En el sentido más real y más hondo que existe. Cada día que os amáis, cada perdón que os dais, cada sacrificio silencioso, cada noche en la que elegís al otro aunque cueste, es un paso más en ese camino de santidad que recorréis juntos. Tu cónyuge te acompaña, te sostiene, te empuja, te corrige, te abraza. Y al final del camino ---cuando llegue el momento de cruzar esa puerta definitiva---, será quien te tome de la mano y te la ponga en la mano de Dios para empezar la vida nueva.

Del bebé indefenso del primer capítulo al esposo que llega a las puertas del cielo de la mano de quien amó: esa es la historia completa. Esa es la historia para la que estáis diseñados.

Que así sea.

\paracaminar

\begin{enumerate}
  \item Pensad juntos en vuestra «aldea»: ¿quiénes son las personas ---de vuestra familia o de fuera de ella--- que os sostienen, que os cuidan, que hacen posible vuestra vida en común? ¿Les habéis dado las gracias alguna vez?

  \item El perdón diario: ¿cómo gestionáis los roces y los pequeños conflictos en vuestra vida juntos? ¿Hay alguna cuenta pendiente, alguna herida pequeña que no se ha curado, que podríais atender hoy? Antes de cerrar este libro, miraos a los ojos y decid en voz alta algo que necesitéis perdonar o algo por lo que necesitéis pedir perdón.

  \item Este libro se titula \textit{Diseñados para amar}. Ahora que lo habéis recorrido juntos, escribid cada uno en un papel ---sin enseñárselo al otro--- qué ha cambiado en vuestra forma de ver el matrimonio, la familia o el amor. Después, intercambiad los papeles y leed en silencio lo que ha escrito el otro. No hace falta comentar nada. A veces basta con saber.
\end{enumerate}

\vinetacapitulo{ch12}


% ===========================
% BIBLIOGRAFÍA
% ===========================
\chapter*{Para saber más}
\addcontentsline{toc}{chapter}{Para saber más}
\markboth{Para saber más}{Para saber más}

Lo que sigue no es una bibliografía académica, sino una selección de lecturas para quien quiera seguir tirando del hilo. He incluido solo los textos que me parecen más accesibles o más importantes, con una breve nota sobre cada uno. Las referencias completas de los estudios científicos citados en las notas a pie de página se encuentran en las propias notas.

\section*{Sobre la biología del vínculo}

\begin{description}[style=nextline, leftmargin=0pt, labelwidth=0pt]

\item[Helen Fisher, \textit{Why We Love} (2004)]
La mejor introducción a la neurociencia del enamoramiento. Fisher explica con claridad cómo la dopamina, la noradrenalina y la serotonina crean ese estado de «locura transitoria» que llamamos enamorarnos. Muy ameno y riguroso a la vez.

\item[Helen Fisher, \textit{Anatomy of Love} (1992, ed. rev. 2016)]
Más ambicioso que el anterior: recorre la historia evolutiva del emparejamiento humano, desde la sabana africana hasta las apps de citas. Imprescindible para entender por qué somos como somos en pareja.

\item[John Bowlby, \textit{Attachment and Loss} (1969--1980)]
La obra fundacional de la teoría del apego. Son tres volúmenes densos, pero el primero (\textit{Attachment}) es el más accesible y el que más directamente conecta con lo que contamos en los capítulos 1 y 2. Para quien prefiera algo más breve, el artículo de Hazan y Shaver (1987), que aplica la teoría del apego al amor de pareja, es un buen punto de partida.

\item[Richard Wrangham, \textit{Catching Fire} (2009)]
Un libro fascinante sobre cómo el dominio del fuego y la cocción de alimentos cambiaron el curso de la evolución humana. Lo citamos brevemente en el capítulo 1, pero merece una lectura completa: cambia la forma de ver nuestra propia especie.

\item[Ashley Montagu, \textit{Touching} (1971)]
Un clásico sobre el sentido del tacto y su papel en el desarrollo humano. Sorprende descubrir cuánto depende nuestra salud física y emocional del contacto piel con piel.

\end{description}

\section*{Sobre el amor y la pareja}

\begin{description}[style=nextline, leftmargin=0pt, labelwidth=0pt]

\item[Erich Fromm, \textit{El arte de amar} (1956)]
Un pequeño gran libro. Fromm argumenta que el amor no es un sentimiento que te sucede, sino un arte que se aprende y se practica. Su distinción entre amor inmaduro («te amo porque te necesito») y amor maduro («te necesito porque te amo») sigue siendo una de las mejores frases escritas sobre el tema.

\item[C.\,S. Lewis, \textit{Los cuatro amores} (1960)]
Lewis distingue entre afecto, amistad, eros y caridad con una lucidez y un humor que solo él tenía. Es breve, profundo y muy disfrutable. Ideal para leer en pareja.

\item[Deborah Tannen, \textit{You Just Don't Understand} (1990)]
Un estudio sobre cómo hombres y mujeres usan el lenguaje de formas distintas, y cómo eso genera malentendidos crónicos en las parejas. Práctico y revelador. Existe traducción al español.

\end{description}

\section*{Sobre la visión cristiana del matrimonio}

\begin{description}[style=nextline, leftmargin=0pt, labelwidth=0pt]

\item[Juan Pablo II, \textit{Hombre y mujer los creó} (1979--1984)]
Las catequesis de los miércoles sobre la «Teología del Cuerpo». Es la obra más ambiciosa jamás escrita sobre el significado del cuerpo, la sexualidad y el matrimonio desde una perspectiva cristiana. No es lectura fácil ---el estilo de Wojtyła es filosófico y denso---, pero las ideas que contiene son revolucionarias. Para una introducción más accesible, véase el libro de Yves Semen que se cita más abajo.

\item[Yves Semen, \textit{La sexualidad según Juan Pablo II} (2005)]
La mejor puerta de entrada a la Teología del Cuerpo. Semen traduce el pensamiento de Wojtyła a un lenguaje asequible sin traicionar su profundidad. Muy recomendable como primer contacto.

\item[Benedicto XVI, \textit{Deus Caritas Est} (2005)]
La primera encíclica de Benedicto XVI es una meditación luminosa sobre el amor: desde el eros hasta el ágape, desde la experiencia humana hasta la caridad de Dios. Breve, bella y sorprendentemente personal.

\item[Francisco, \textit{Amoris Laetitia} (2016)]
La exhortación del Papa Francisco sobre la familia. El capítulo 4, que es un comentario al himno de la caridad de san Pablo (1~Cor 13), es una de las páginas más hermosas escritas sobre el amor conyugal en un documento eclesial. Muy pastoral, muy realista, muy humana.

\item[Juan Pablo II, \textit{Familiaris Consortio} (1981)]
El documento de referencia sobre la familia en la doctrina católica contemporánea. Más sistemático que \textit{Amoris Laetitia}, pero complementario. Especialmente valioso para entender la indisolubilidad, la fidelidad y la apertura a la vida.

\item[\textit{Gaudium et Spes}, nn. 47--52 (1965)]
Los números del Concilio Vaticano II dedicados al matrimonio y la familia siguen siendo un texto de referencia. Breve, denso y profundamente renovador para su época.

\end{description}

\section*{Para pensar más}

\begin{description}[style=nextline, leftmargin=0pt, labelwidth=0pt]

\item[Emmanuel Levinas, \textit{Totalidad e infinito} (1961)]
Para lectores filosóficamente audaces. Levinas desarrolla la idea de que el encuentro con el «rostro del otro» es la experiencia ética fundamental. Su concepto de alteridad ---el otro como irreductiblemente otro--- ilumina de forma profunda lo que decimos en el capítulo 5 sobre la diferencia como condición del encuentro.

\item[Gabriel Marcel, \textit{Être et avoir} (1935)]
Marcel distingue entre «tener» y «ser», entre poseer y entregarse. Su filosofía del compromiso y la fidelidad creadora conecta directamente con lo que exploramos en los capítulos 8 y 9. Breve y luminoso.

\end{description}


\end{document}
